% =============================================================
%  MAIN — Common Lisp: Guida Pratica
%  Compilare con:  pdflatex main  (due volte per i riferimenti)
% =============================================================
\documentclass[a4paper,11pt,twoside,openright]{book}

% =============================================================
%  PREAMBLE — Common Lisp: Guida Pratica
%  Compilare con: pdflatex main  (x2)
% =============================================================

% --- Encoding e lingua ---
\usepackage[T1]{fontenc}
\usepackage[utf8]{inputenc}
\usepackage[italian]{babel}

% --- Font ---
\usepackage{lmodern}
\usepackage{microtype}

% --- Layout pagina ---
\usepackage[
  a4paper,
  top=2.5cm, bottom=2.8cm,
  inner=2.8cm, outer=2.2cm,
  marginparwidth=1.5cm
]{geometry}

% --- Colori ---
\usepackage[dvipsnames,svgnames,table]{xcolor}

\definecolor{clBg}{RGB}{248,248,248}       % sfondo listing
\definecolor{clKeyword}{RGB}{0,0,180}      % keyword Lisp
\definecolor{clComment}{RGB}{100,100,100}  % commenti
\definecolor{clString}{RGB}{0,120,0}       % stringhe

% Colori difficoltà
\definecolor{colOne}{RGB}{130,130,130}     % ★      grigio
\definecolor{colTwo}{RGB}{34,110,34}       % ★★     verde
\definecolor{colThree}{RGB}{30,80,160}     % ★★★    blu
\definecolor{colFour}{RGB}{200,100,0}      % ★★★★   arancio
\definecolor{colFive}{RGB}{180,30,30}      % ★★★★★  rosso

% Colori box
\definecolor{boxEssFrame}{RGB}{30,80,160}
\definecolor{boxEssBack}{RGB}{235,241,252}
\definecolor{boxColFrame}{RGB}{110,110,110}
\definecolor{boxColBack}{RGB}{245,245,245}
\definecolor{boxHintFrame}{RGB}{180,130,0}
\definecolor{boxHintBack}{RGB}{255,252,220}
\definecolor{boxSolFrame}{RGB}{34,110,34}
\definecolor{boxSolBack}{RGB}{236,248,236}
\definecolor{boxRagFrame}{RGB}{140,50,140}
\definecolor{boxRagBack}{RGB}{248,236,248}

% --- Matematica ---
\usepackage{amsmath}
\usepackage{amssymb}

% --- Listings (codice Lisp) ---
\usepackage{listings}

\lstdefinelanguage{CommonLisp}{%
  sensitive=false,
  morecomment=[l]{;},
  morestring=[b]",
  morekeywords=[1]{%
    defun, defmacro, defparameter, defvar, defconstant,
    defstruct, defclass, defmethod, defgeneric,
    let, let*, flet, labels, setf, setq,
    if, cond, when, unless, case, ecase,
    and, or, not,
    progn, prog1, prog2,
    loop, dotimes, dolist, do, do*,
    lambda, function,
    funcall, apply,
    mapcar, maplist, mapcan, mapc,
    remove-if, remove-if-not, find-if, find-if-not,
    every, some, notany, notevery,
    reduce, count-if, count,
    sort, stable-sort,
    return, return-from,
    values, multiple-value-bind, multiple-value-list,
    destructuring-bind,
    handler-case, handler-bind, error, warn, signal,
    with-open-file, open, close,
    format, print, princ, prin1, terpri, fresh-line,
    read, read-line, read-char,
    car, cdr, cons, list, append, reverse, length,
    first, second, third, fourth, fifth,
    rest, last, nth, nthcdr, butlast,
    null, atom, listp, consp, numberp, stringp,
    symbolp, functionp, integerp, floatp, characterp,
    eq, eql, equal, equalp,
    quote, backquote,
    make-hash-table, gethash, remhash, maphash, hash-table-count,
    make-array, aref, array-dimensions,
    make-instance, slot-value, initialize-instance,
    in-package, defpackage, use-package, export, import,
    compile, load, require, provide,
    t, nil%
  },
  morekeywords=[2]{%
    +, -, *, /, mod, rem, expt, sqrt, abs,
    floor, ceiling, truncate, round,
    min, max, 1+, 1-,
    =, /=, <, >, <=, >=,
    char, char-code, code-char, char-upcase, char-downcase,
    string, string-upcase, string-downcase, string-length,
    coerce, type-of, typep, check-type%
  }
}

\lstdefinestyle{lispstyle}{%
  language=CommonLisp,
  basicstyle=\ttfamily\small,
  keywordstyle=[1]{\color{clKeyword}\bfseries},
  keywordstyle=[2]{\color{clKeyword}},
  commentstyle=\itshape\color{clComment},
  stringstyle=\color{clString},
  backgroundcolor=\color{clBg},
  frame=single,
  framesep=4pt,
  framerule=0.4pt,
  rulecolor=\color{gray!50},
  xleftmargin=8pt,
  xrightmargin=4pt,
  breaklines=true,
  breakatwhitespace=true,
  showstringspaces=false,
  tabsize=2,
  numbers=none,
  upquote=true,
  extendedchars=true,
  literate=
    {→}{{$\rightarrow$}}1
    {★}{{\textcolor{colOne}{$\bigstar$}}}1
    {à}{{\`a}}1  {è}{{\`e}}1  {é}{{\'{e}}}1  {ì}{{\`{\i}}}1
    {ò}{{\`o}}1  {ù}{{\`u}}1
    {À}{{\`A}}1  {È}{{\`E}}1  {É}{{\'{E}}}1
    {Ì}{{\`I}}1  {Ò}{{\`O}}1  {Ù}{{\`U}}1
}

\lstset{style=lispstyle}

% Shorthand per listing inline
\newcommand{\cl}[1]{\lstinline[style=lispstyle]{#1}}

% --- tcolorbox ---
\usepackage[most]{tcolorbox}
\tcbuselibrary{listings, breakable, skins}

% Stelle di difficoltà
\newcommand{\unaStella}{\textcolor{colOne}{$\bigstar$}}
\newcommand{\dueStelle}{\textcolor{colTwo}{$\bigstar\bigstar$}}
\newcommand{\treStelle}{\textcolor{colThree}{$\bigstar\bigstar\bigstar$}}
\newcommand{\quattroStelle}{\textcolor{colFour}{$\bigstar\bigstar\bigstar\bigstar$}}
\newcommand{\cinqueStelle}{\textcolor{colFive}{$\bigstar\bigstar\bigstar\bigstar\bigstar$}}

% Mappa numero di stelle → colore frame
\newcommand{\starcolor}[1]{%
  \ifcase#1
    \color{colOne}%   0 (non usato)
  \or\color{colOne}%  1
  \or\color{colTwo}%  2
  \or\color{colThree}% 3
  \or\color{colFour}% 4
  \else\color{colFive}% 5
  \fi
}

% ---- Box ESERCIZIO ESSENZIALE ◆ ----
% Uso: \begin{essbox}{n}{titolo}{stelle} ... \end{essbox}
\newenvironment{essbox}[3]{%
  \medskip
  \begin{tcolorbox}[
    enhanced,
    breakable,
    colframe=boxEssFrame,
    colback=boxEssBack,
    fonttitle=\bfseries\small,
    title={{\small\textcolor{boxEssFrame}{$\blacklozenge$}}\enspace Esercizio~#1\enspace---\enspace#2\enspace
      \hfill\ifcase#3\or\unaStella\or\dueStelle\or\treStelle\or\quattroStelle\or\cinqueStelle\fi},
    arc=3pt,
    boxrule=1.2pt,
    left=6pt, right=6pt, top=4pt, bottom=4pt
  ]
}{%
  \end{tcolorbox}
  \medskip
}

% ---- Box ESERCIZIO COLLATERALE ○ ----
% Uso: \begin{colbox}{n}{titolo}{stelle} ... \end{colbox}
\newenvironment{colbox}[3]{%
  \medskip
  \begin{tcolorbox}[
    enhanced,
    breakable,
    colframe=boxColFrame,
    colback=boxColBack,
    fonttitle=\bfseries\small,
    title={{\small\textcolor{boxColFrame}{$\circ$}}\enspace Esercizio~#1\enspace---\enspace#2\enspace
      \hfill\ifcase#3\or\unaStella\or\dueStelle\or\treStelle\or\quattroStelle\or\cinqueStelle\fi},
    arc=3pt,
    boxrule=0.8pt,
    borderline={0.8pt}{0pt}{boxColFrame, dashed},
    left=6pt, right=6pt, top=4pt, bottom=4pt
  ]
}{%
  \end{tcolorbox}
  \medskip
}

% ---- Box SUGGERIMENTO ----
\newenvironment{hintbox}[1]{%
  \medskip
  \begin{tcolorbox}[
    enhanced,
    breakable,
    colframe=boxHintFrame,
    colback=boxHintBack,
    fonttitle=\bfseries\small,
    title={Suggerimento --- Esercizio~#1},
    arc=3pt,
    boxrule=0.8pt,
    left=6pt, right=6pt, top=4pt, bottom=4pt
  ]
  \small
}{%
  \end{tcolorbox}
  \medskip
}

% ---- Box SOLUZIONE ----
\newenvironment{solbox}[1]{%
  \medskip
  \begin{tcolorbox}[
    enhanced,
    breakable,
    colframe=boxSolFrame,
    colback=boxSolBack,
    fonttitle=\bfseries\small,
    title={Soluzione --- Esercizio~#1},
    arc=3pt,
    boxrule=0.8pt,
    left=6pt, right=6pt, top=4pt, bottom=4pt
  ]
  \small
}{%
  \end{tcolorbox}
  \medskip
}

% ---- Box RAGIONAMENTO (per ★★★★/★★★★★) ----
\newenvironment{ragionamentoBox}[1]{%
  \medskip
  \begin{tcolorbox}[
    enhanced,
    breakable,
    colframe=boxRagFrame,
    colback=boxRagBack,
    fonttitle=\bfseries\small,
    title={Ragionamento --- Esercizio~#1},
    arc=3pt,
    boxrule=0.8pt,
    left=6pt, right=6pt, top=4pt, bottom=4pt
  ]
  \small
}{%
  \end{tcolorbox}
  \medskip
}

% Alias brevi per compatibilita' (comandi -> ambienti)
% Non usare questi per listing, usare direttamente gli ambienti
\newcommand{\essenzialeBox}[4]{\begin{essbox}{#1}{#2}{#3}#4\end{essbox}}
\newcommand{\collateraleBox}[4]{\begin{colbox}{#1}{#2}{#3}#4\end{colbox}}
\newcommand{\hintBox}[2]{\begin{hintbox}{#1}#2\end{hintbox}}
\newcommand{\soluzioneBox}[2]{\begin{solbox}{#1}#2\end{solbox}}

% --- Intestazioni capitolo ---
\usepackage{titlesec}
\titleformat{\chapter}[display]
  {\normalfont\Large\bfseries\color{boxEssFrame}}
  {\chaptertitlename\ \thechapter}{10pt}{\Huge}
\titlespacing*{\chapter}{0pt}{20pt}{30pt}

% --- Header/Footer ---
\usepackage{fancyhdr}
\pagestyle{fancy}
\fancyhf{}
\fancyhead[LE,RO]{\small\thepage}
\fancyhead[LO]{\small\nouppercase{\rightmark}}
\fancyhead[RE]{\small\nouppercase{\leftmark}}
\renewcommand{\headrulewidth}{0.4pt}
\fancypagestyle{plain}{%
  \fancyhf{}
  \fancyfoot[C]{\small\thepage}
  \renewcommand{\headrulewidth}{0pt}
}

% --- Hyperref (deve essere quasi l'ultimo) ---
\usepackage[
  colorlinks=true,
  linkcolor=boxEssFrame,
  urlcolor=colTwo,
  citecolor=colThree,
  bookmarksnumbered=true,
  bookmarksopen=true,
  pdftitle={Common Lisp --- Guida Pratica},
  pdfauthor={},
  pdfsubject={Common Lisp},
  pdfkeywords={Common Lisp, programmazione funzionale, esercizi}
]{hyperref}

% --- Multicol per cheat sheet ---
\usepackage{multicol}
\usepackage{booktabs}
\usepackage{array}
\usepackage{enumitem}

% Impaginazione: evita vedove e orfane
\widowpenalty=10000
\clubpenalty=10000

% Etichette esercizi, hint, soluzioni
% Uso: \label{ex:42}  \label{hint:42}  \label{sol:42}

% Comandi utili
\newcommand{\kw}[1]{\texttt{\bfseries #1}}  % keyword inline
\newcommand{\fn}[1]{\texttt{#1}}            % funzione inline
\newcommand{\result}{{\color{clComment}; $\rightarrow$\ }}

% Nomi simboli
\newcommand{\essential}{{\color{boxEssFrame}$\blacklozenge$}}
\newcommand{\optional}{{\color{boxColFrame}$\circ$}}


\begin{document}

% ---- Copertina -----------------------------------------------
\begin{titlepage}
  \centering
  \vspace*{3cm}
  {\Huge\bfseries\color{boxEssFrame} Common Lisp}\\[0.6cm]
  {\LARGE\itshape Guida Pratica}\\[0.4cm]
  {\large con 148 esercizi graduali}\\[3cm]
  \rule{0.4\textwidth}{1pt}\\[1cm]
  {\large\itshape Per studenti con basi di programmazione imperativa}\\[0.5cm]
  {\normalsize Ambiente consigliato: Allegro CL Free Express Edition}\\[2cm]
  \rule{0.4\textwidth}{0.4pt}\\[0.8cm]
  {\large\the\year}
  \vfill
\end{titlepage}

% ---- Pagina crediti ------------------------------------------
\thispagestyle{empty}
\vspace*{4cm}
\noindent\textbf{Note sull'ambiente}\\[0.3em]
Gli esempi e le soluzioni di questo manuale sono conformi allo
standard \textbf{ANSI Common Lisp} e sono stati sviluppati e
testati con \textbf{Allegro CL Free Express Edition}.
Qualsiasi implementazione conforme allo standard (SBCL, CCL,
CLISP, ECL, ecc.) è egualmente compatibile.

\vspace{1.5em}
\noindent\textbf{Uso del manuale}\\[0.3em]
Questo materiale è pensato per l'uso didattico in classe e per
lo studio individuale.

\vspace{1.5em}
\noindent\textit{Compilato con \LaTeX.}

\clearpage

% ---- Come leggere questo manuale ----------------------------
\chapter*{Come leggere questo manuale}
\addcontentsline{toc}{chapter}{Come leggere questo manuale}

\noindent
Il manuale è diviso in quattro sezioni:

\begin{enumerate}[leftmargin=2em]
  \item \textbf{Cheat sheet} — una o due pagine da tenere
        aperte durante gli esercizi, con i costrutti
        più usati e link alle sezioni della reference.
  \item \textbf{Quick reference} — descrizione concisa di
        tutti i costrutti del linguaggio, con esempi.
  \item \textbf{Esercizi} — 148 esercizi in 10 capitoli,
        dalla sintassi base ai progetti complessi.
  \item \textbf{Suggerimenti e soluzioni} — hint per gli
        esercizi più difficili e soluzioni commentate.
\end{enumerate}

\bigskip
\noindent\textbf{Legenda simboli}

\medskip
\begin{tabular}{@{}lp{10cm}@{}}
  \essential & \textbf{Esercizio essenziale} --- introduce un costrutto
               o un pattern nuovo. Da svolgere sempre.\\[6pt]
  \optional  & \textbf{Esercizio collaterale} --- consolida o approfondisce
               quanto già visto. Consigliato per chi vuole
               padronanza completa del capitolo.\\[6pt]
  \unaStella & Facile\\[2pt]
  \dueStelle & Medio\\[2pt]
  \treStelle & Difficile\\[2pt]
  \quattroStelle & Avanzato\\[2pt]
  \cinqueStelle  & Progetto\\
\end{tabular}

\bigskip
\noindent\textbf{Percorso minimo (53 esercizi)}\\
Svolgi solo gli esercizi \essential: sono sufficienti a
coprire tutti i concetti fondamentali del linguaggio.

\noindent\textbf{Percorso completo (148 esercizi)}\\
Svolgi anche i collaterali \optional\ per consolidare
ogni argomento e acquisire una padronanza più solida.

\bigskip
\noindent\textbf{Suggerimenti}\\
Gli esercizi di difficoltà \treStelle\ o superiore hanno
un suggerimento nella sezione apposita. Cercalo prima di
guardare la soluzione.

\bigskip
\noindent\textbf{Soluzioni}\\
Le soluzioni sono nella sezione finale. Per gli esercizi
\unaStella/\dueStelle\ è riportato solo il codice.
Per \treStelle\ ci sono commenti inline.
Per \quattroStelle/\cinqueStelle\ c'è anche una sezione di
\textit{ragionamento} che spiega la strategia e la
complessità algoritmica.

\clearpage

% ---- Indice --------------------------------------------------
\tableofcontents
\clearpage

% ---- Sezioni principali --------------------------------------
% =============================================================
%  00_cheatsheet.tex — Common Lisp Cheat Sheet
%  1-2 pagine dense, due colonne, link interni alla reference
% =============================================================
\chapter*{Cheat Sheet}
\addcontentsline{toc}{chapter}{Cheat Sheet}
\markboth{Cheat Sheet}{Cheat Sheet}

{\small
\begin{multicols}{2}

% ---- Aritmetica --------------------------------------------
\noindent\textbf{\hyperref[sec:numeri]{Aritmetica}}\\
\begin{tabular}{@{}ll@{}}
\fn{(+ a b \ldots)}    & somma \\
\fn{(- a b)}           & differenza \\
\fn{(* a b \ldots)}    & prodotto \\
\fn{(/ a b)}           & quoziente (esatto) \\
\fn{(mod a b)}         & resto (segno di \texttt{b}) \\
\fn{(rem a b)}         & resto (segno di \texttt{a}) \\
\fn{(expt b e)}        & $b^e$ \\
\fn{(sqrt x)}          & radice quadrata \\
\fn{(abs x)}           & valore assoluto \\
\fn{(1+ x)} / \fn{(1- x)} & inc/dec \\
\fn{(floor x)}         & parte intera $\lfloor x \rfloor$ \\
\fn{(ceiling x)}       & $\lceil x \rceil$ \\
\fn{(round x)}         & arrotonda \\
\fn{(max a b)} / \fn{(min a b)} & massimo/minimo \\
\end{tabular}

\medskip
% ---- Confronti --------------------------------------------
\noindent\textbf{\hyperref[sec:numeri]{Confronti numerici}}\\
\begin{tabular}{@{}ll@{}}
\fn{(= a b)}  & uguale \\
\fn{(/= a b)} & diverso \\
\fn{(< a b)}  \fn{(> a b)} \fn{(<= a b)} \fn{(>= a b)} & ordine \\
\end{tabular}

\medskip
% ---- Booleani ---------------------------------------------
\noindent\textbf{\hyperref[sec:booleani]{Booleani}}\\
\begin{tabular}{@{}ll@{}}
\fn{T} / \fn{NIL}    & vero / falso (e lista vuota) \\
\fn{(and a b \ldots)} & e logico (short-circuit) \\
\fn{(or a b \ldots)}  & o logico (short-circuit) \\
\fn{(not x)}          & negazione \\
\end{tabular}

\medskip
% ---- Variabili e binding ----------------------------------
\noindent\textbf{\hyperref[sec:variabili]{Definizioni e binding}}\\
\begin{tabular}{@{}ll@{}}
\fn{(defun nome (p\ldots) corpo)} & funzione \\
\fn{(defparameter *x* val)}       & variabile globale \\
\fn{(defvar *x* val)}             & globale (no-overwrite) \\
\fn{(defconstant +k+ val)}        & costante \\
\fn{(let ((x v)\ldots) corpo)}    & binding locale \\
\fn{(let* ((x v)\ldots) corpo)}   & binding sequenziale \\
\fn{(setf x val)}                 & assegnamento \\
\end{tabular}

\medskip
% ---- Condizioni -------------------------------------------
\noindent\textbf{\hyperref[sec:condizioni]{Condizioni}}\\
\begin{tabular}{@{}ll@{}}
\fn{(if test then else)}   & ramo singolo \\
\fn{(cond (t1 e1)\ldots)}  & rami multipli \\
\fn{(when test corpo\ldots)}   & solo se vero \\
\fn{(unless test corpo\ldots)} & solo se falso \\
\fn{(case val (k e)\ldots)}    & switch su valore \\
\end{tabular}

\columnbreak

% ---- Liste ------------------------------------------------
\noindent\textbf{\hyperref[sec:liste]{Liste — costruzione e accesso}}\\
\begin{tabular}{@{}ll@{}}
\fn{(cons x lst)}        & prependi \\
\fn{(list a b \ldots)}   & crea lista \\
\fn{(append l1 l2)}      & concatena \\
\fn{(reverse lst)}       & inverti \\
\fn{(car lst)} / \fn{(first lst)}  & primo \\
\fn{(cdr lst)} / \fn{(rest lst)}   & resto \\
\fn{(cadr lst)} / \fn{(second lst)}& secondo \\
\fn{(nth n lst)}         & n-esimo (0-based) \\
\fn{(last lst)}          & ultima cons-cell \\
\fn{(length lst)}        & lunghezza \\
\fn{(null lst)}          & lista vuota? \\
\fn{(member x lst)}      & appartiene? \\
\end{tabular}

\medskip
% ---- Higher-order -----------------------------------------
\noindent\textbf{\hyperref[sec:higherorder]{Higher-order}}\\
\begin{tabular}{@{}ll@{}}
\fn{(mapcar \#'f lst)}       & trasforma \\
\fn{(mapcan \#'f lst)}       & trasforma+appiattisce \\
\fn{(remove-if \#'p lst)}    & rimuovi se vero \\
\fn{(remove-if-not \#'p lst)}& tieni se vero \\
\fn{(find-if \#'p lst)}      & primo match \\
\fn{(every \#'p lst)}        & tutti veri? \\
\fn{(some \#'p lst)}         & almeno uno vero? \\
\fn{(reduce \#'f lst)}       & fold sinistro \\
\fn{(sort lst \#'<)}         & ordina (distruttivo) \\
\fn{(count-if \#'p lst)}     & conta match \\
\fn{(funcall f a b)}         & chiama f \\
\fn{(apply f args)}          & chiama f con lista args \\
\fn{\#'nome}                 & oggetto funzione \\
\fn{(lambda (x) corpo)}      & funzione anonima \\
\end{tabular}

\medskip
% ---- Iterazione -------------------------------------------
\noindent\textbf{\hyperref[sec:iterazione]{Iterazione}}\\
\begin{tabular}{@{}ll@{}}
\fn{(dotimes (i n) corpo)}         & 0..n-1 \\
\fn{(dolist (x lst) corpo)}        & su lista \\
\fn{(loop for x in lst do \ldots)} & loop su lista \\
\fn{(loop for i from a to b do \ldots)} & loop numerico \\
\fn{(loop \ldots collect e)}       & raccogli risultati \\
\fn{(loop \ldots sum e)}           & somma \\
\end{tabular}

\medskip
% ---- Strutture dati specializzate -------------------------
\noindent\textbf{\hyperref[sec:hashtable]{Hash table}}\\
\begin{tabular}{@{}ll@{}}
\fn{(make-hash-table)}              & crea \\
\fn{(gethash k ht)}                 & leggi \\
\fn{(setf (gethash k ht) v)}        & scrivi \\
\fn{(remhash k ht)}                 & elimina \\
\fn{(maphash \#'f ht)}              & itera \\
\end{tabular}

\medskip
% ---- Array ------------------------------------------------
\noindent\textbf{\hyperref[sec:array]{Array e vettori}}\\
\begin{tabular}{@{}ll@{}}
\fn{\#(1 2 3)}                    & vettore letterale \\
\fn{(make-array n)}               & crea vettore \\
\fn{(make-array '(r c))}          & crea matrice \\
\fn{(aref a i)} / \fn{(aref m r c)} & accesso \\
\fn{(setf (aref a i) v)}          & assegna \\
\fn{(length a)}                   & dimensione \\
\end{tabular}

\medskip
% ---- Defstruct --------------------------------------------
\noindent\textbf{\hyperref[sec:strutture]{Strutture (defstruct)}}\\
\begin{tabular}{@{}ll@{}}
\fn{(defstruct nome slot\ldots)}       & definisci \\
\fn{(make-nome :slot val)}             & crea istanza \\
\fn{(nome-slot ist)}                   & accedi \\
\fn{(setf (nome-slot ist) v)}          & modifica \\
\end{tabular}

\medskip
% ---- I/O --------------------------------------------------
\noindent\textbf{\hyperref[sec:io]{I/O}}\\
\begin{tabular}{@{}ll@{}}
\fn{(read)}                    & leggi da stdin \\
\fn{(read-line)}               & leggi riga \\
\fn{(print x)}                 & stampa + newline \\
\fn{(princ x)}                 & stampa senza quote \\
\fn{(format t "\textasciitilde A\textasciitilde \%"\ x)} & formattato \\
\fn{\textasciitilde A} \fn{\textasciitilde D} \fn{\textasciitilde F} \fn{\textasciitilde \%} & gen/int/float/newline \\
\end{tabular}

\medskip
% ---- Macro ------------------------------------------------
\noindent\textbf{\hyperref[sec:macro]{Macro}}\\
\begin{tabular}{@{}ll@{}}
\fn{(defmacro nome (p\ldots) corpo)} & definisci \\
\fn{`(lista ,expr ,@lista)}          & backquote/unquote \\
\fn{(gensym)}                        & simbolo unico \\
\end{tabular}

\medskip
% ---- Predicati --------------------------------------------
\noindent\textbf{\hyperref[sec:predicati]{Predicati comuni}}\\
\begin{tabular}{@{}ll@{}}
\fn{(null x)}      & NIL? \\
\fn{(atom x)}      & atomo (non cons)? \\
\fn{(listp x)}     & lista? \\
\fn{(numberp x)}   & numero? \\
\fn{(symbolp x)}   & simbolo? \\
\fn{(stringp x)}   & stringa? \\
\fn{(eq a b)}      & identità oggetto \\
\fn{(eql a b)}     & identità + numeri/char \\
\fn{(equal a b)}   & uguaglianza strutturale \\
\end{tabular}

\end{multicols}

\vfill
\noindent\rule{\textwidth}{0.4pt}\\[2pt]
{\footnotesize
\textbf{Legenda:}
\essential\ = esercizio essenziale\quad
\optional\ = esercizio collaterale\quad
\unaStella\ facile\quad
\dueStelle\ medio\quad
\treStelle\ difficile\quad
\quattroStelle\ avanzato\quad
\cinqueStelle\ progetto
}
} % fine \small

% =============================================================
%  01_reference.tex — Quick Reference dettagliata
% =============================================================
\chapter*{Quick Reference}
\addcontentsline{toc}{chapter}{Quick Reference}
\markboth{Quick Reference}{Quick Reference}

% --- Box per le macro-aree della reference ---
% Usa gli stessi colori semantici del cheatsheet
% Testo interno: grigio e più piccolo per non competere col codice
% UNBREAKABLE per evitare rotture interne
\newtcolorbox{refbox}[2][]{%
  enhanced,
  colback=#2!4,
  colframe=#2!70!black,
  fonttitle=\bfseries\sffamily\normalsize,
  coltitle=white,
  title={#1},
  boxrule=0.8pt,
  arc=4pt,
  left=12pt, right=12pt, top=10pt, bottom=10pt,
  toptitle=4pt, bottomtitle=4pt,
  before skip=14pt plus 4pt,
  after skip=14pt plus 4pt,
  fontupper=\small\color{black!75},  % testo grigio, più piccolo
  parskip=6pt                         % spazio tra paragrafi interni
}

% Box per note/avvertimenti interni alle sezioni
\newtcolorbox{refnote}[1][]{%
  enhanced,
  colback=black!3,
  colframe=black!30,
  fonttitle=\bfseries\small,
  coltitle=black!70,
  title={#1},
  boxrule=0.4pt,
  arc=3pt,
  left=8pt, right=8pt, top=6pt, bottom=6pt,
  before skip=10pt,
  after skip=10pt,
  fontupper=\small\color{black!65}   % testo ancora più leggero
}

% Comando per sottotitoli interni ai box (evita break dopo)
\newcommand{\reftitle}[1]{\textbf{#1}\nopagebreak\par\nopagebreak}

\noindent
Questa sezione descrive i costrutti fondamentali di Common Lisp
in forma concisa. Per ogni argomento: sintassi, spiegazione e
almeno un esempio.

\vspace{16pt}

\begin{tcolorbox}[
  enhanced,
  colback=black!2,
  colframe=black!25,
  boxrule=0.4pt,
  arc=3pt,
  left=12pt, right=12pt, top=10pt, bottom=10pt
]
\centering
\textbf{Legenda colori}\\[8pt]
\textcolor{csBlue}{\rule{1.2em}{1.2em}}\enspace Basi e stile\qquad
\textcolor{csOrange}{\rule{1.2em}{1.2em}}\enspace Controllo del flusso\qquad
\textcolor{csGreen}{\rule{1.2em}{1.2em}}\enspace Liste e sequenze\\[6pt]
\textcolor{csTeal}{\rule{1.2em}{1.2em}}\enspace Strutture dati e I/O\qquad
\textcolor{csPurple}{\rule{1.2em}{1.2em}}\enspace Astrazione e OOP
\end{tcolorbox}

\vspace{16pt}

% =============================================================
% MACRO-AREA 1: BASI DEL LINGUAGGIO
% =============================================================
\addcontentsline{toc}{section}{Basi del linguaggio}

% --- REPL e S-expression ---
\begin{refbox}[REPL e S-expression]{csBlue}
\label{sec:repl}

Common Lisp usa la notazione \textbf{S-expression} (prefix):
il primo elemento di ogni lista è l'operatore, i successivi
sono gli argomenti.

\begin{lstlisting}
(+ 1 2 3)        ; -> 6
(* 2 (+ 3 4))    ; -> 14
\end{lstlisting}

I commenti si scrivono con \fn{;} (riga singola) o
\fn{\#|~...|~\#} (multiriga):

\begin{lstlisting}
; commento su una riga
#| commento
   su piu' righe |#
\end{lstlisting}

Il REPL legge, valuta e stampa ogni espressione.
Le variabili \fn{*}, \fn{**}, \fn{***} contengono gli ultimi
tre risultati restituiti.

\end{refbox}

% --- Numeri ---
\begin{refbox}[Numeri]{csBlue}
\label{sec:numeri}

Common Lisp distingue interi, razionali, float e complessi:

\begin{lstlisting}
42      -7       ; interi
3.14    1.5e10   ; float
1/3     22/7     ; razionali esatti
#C(1 2)          ; complesso 1+2i
\end{lstlisting}

\textbf{Operatori aritmetici:}\nopagebreak

\begin{lstlisting}
(+ a b)   (- a b)   (* a b)   (/ a b)
(mod 10 3)     ; -> 1  (segno del divisore)
(rem 10 3)     ; -> 1  (segno del dividendo)
(expt 2 10)    ; -> 1024
(sqrt 9)       ; -> 3.0
(abs -5)       ; -> 5
(1+ x) (1- x)  ; incrementa / decrementa
(incf x)       ; setf di (1+ x)
(decf x)       ; setf di (1- x)
(floor 7/2)    ; -> 3  (verso -inf)
(ceiling 7/2)  ; -> 4  (verso +inf)
(truncate 7/2) ; -> 3  (verso zero)
(round 2.5)    ; -> 2  (round-to-even)
(max 3 1 4)    ; -> 4
(min 3 1 4)    ; -> 1
\end{lstlisting}

\textbf{Confronti numerici} (usare \fn{=}, non \fn{eq}/\fn{equal}):\nopagebreak

\begin{lstlisting}
(= a b)  (/= a b)  (< a b)  (> a b)  (<= a b)  (>= a b)
\end{lstlisting}

\end{refbox}

% --- Booleani ---
\begin{refbox}[Booleani]{csBlue}
\label{sec:booleani}

\fn{NIL} è l'unico valore falso; qualsiasi altra cosa è vera.
\fn{NIL} è anche la lista vuota.

\begin{lstlisting}
T    ; vero
NIL  ; falso e lista vuota
\end{lstlisting}

\begin{lstlisting}
(and T NIL)    ; -> NIL
(or  NIL 42)   ; -> 42  (primo valore non-NIL)
(not NIL)      ; -> T
\end{lstlisting}

\fn{and} e \fn{or} usano \textit{short-circuit evaluation} e
restituiscono il valore che ha determinato il risultato,
non necessariamente \fn{T} o \fn{NIL}.

\end{refbox}

% --- Simboli, Stringhe, Caratteri ---
\begin{refbox}[Simboli, stringhe e caratteri]{csBlue}
\label{sec:simboli}

\textbf{Simboli:} identificatori unici — lo stesso nome riferisce sempre lo stesso oggetto.
La virgoletta singola \fn{'} previene la valutazione:\nopagebreak

\begin{lstlisting}
'foo          ; simbolo FOO
(eq 'a 'a)   ; -> T  (stessa identita')
\end{lstlisting}

\textbf{Stringhe:}\nopagebreak

\begin{lstlisting}
"hello"
(length "hello")           ; -> 5
(char "hello" 0)           ; -> #\h
(subseq "hello" 1 3)       ; -> "el"
(string-upcase "ciao")     ; -> "CIAO"
(string-downcase "CIAO")   ; -> "ciao"
(concatenate 'string "foo" "bar") ; -> "foobar"
(string= "abc" "abc")      ; -> T  (case-sensitive)
(string-equal "ABC" "abc") ; -> T  (case-insensitive)
(parse-integer "42")       ; -> 42
\end{lstlisting}

\textbf{Caratteri:}\nopagebreak

\begin{lstlisting}
#\A  #\space  #\newline  #\tab
(char-code #\A)          ; -> 65
(code-char 65)           ; -> #\A
(char-upcase #\a)        ; -> #\A
(alpha-char-p #\a)       ; -> T
(digit-char-p #\5)       ; -> 5  (o NIL)
\end{lstlisting}

\end{refbox}

% --- Variabili ---
\begin{refbox}[Variabili e binding]{csBlue}
\label{sec:variabili}

\textbf{Globali:}\nopagebreak
\begin{lstlisting}
(defparameter *contatore* 0)  ; globale, sovrascritta al reload
(defvar *debug* nil)          ; globale, non sovrascritta se gia' legata
(defconstant +pi-greco+ 3.14159265) ; costante
\end{lstlisting}

Convenzione: variabili globali dinamiche tra \fn{*...*} (earmuffs);
costanti tra \fn{+...+}.

\medskip
\textbf{Binding locale — \fn{let}:} tutte le variabili sono
inizializzate in parallelo, non possono riferirsi l'una all'altra:\nopagebreak
\begin{lstlisting}
(let ((x 10)
      (y 20))
  (+ x y))   ; -> 30
\end{lstlisting}

\textbf{Binding sequenziale — \fn{let*}:} ogni variabile può
usare quelle definite prima:\nopagebreak
\begin{lstlisting}
(let* ((r 5)
       (area (* pi r r)))
  area)       ; -> 78.539...
\end{lstlisting}

\textbf{Assegnamento:}\nopagebreak
\begin{lstlisting}
(setf *contatore* (1+ *contatore*))
(incf *contatore*)      ; equivalente piu' conciso
(decf *contatore* 3)    ; decrementa di 3
\end{lstlisting}

\end{refbox}

% --- Funzioni ---
\begin{refbox}[Definizione di funzioni]{csBlue}
\label{sec:funzioni}

\begin{lstlisting}
(defun quadrato (x)
  "Restituisce il quadrato di X."  ; docstring (opzionale)
  (* x x))

(quadrato 5)   ; -> 25
\end{lstlisting}

Il valore di ritorno è quello dell'ultima espressione.
\fn{return-from} permette un'uscita anticipata:

\begin{lstlisting}
(defun cerca (item lst)
  (dolist (x lst nil)          ; nil = valore se lista esaurita
    (when (equal x item)
      (return-from cerca t))))
\end{lstlisting}

\textbf{Argomenti opzionali, rest, keyword:}\nopagebreak
\begin{lstlisting}
; &optional con valore di default
(defun saluta (&optional (nome "mondo"))
  (format t "Ciao, ~A!~%" nome))

; &rest raccoglie argomenti variabili in una lista
(defun somma (&rest numeri)
  (reduce #'+ numeri :initial-value 0))

; &key usa argomenti nominati
(defun crea-rettangolo (&key (w 1) (h 1))
  (* w h))
(crea-rettangolo :w 3 :h 4)   ; -> 12
\end{lstlisting}

\textbf{Funzioni locali} con \fn{flet} e \fn{labels}:\nopagebreak
\begin{lstlisting}
; flet: funzioni locali non ricorsive
(flet ((doppio (x) (* 2 x)))
  (doppio 5))    ; -> 10

; labels: funzioni locali, anche ricorsive e mutuamente ricorsive
(labels ((pari-p (n) (if (= n 0) t (dispari-p (1- n))))
         (dispari-p (n) (if (= n 0) nil (pari-p (1- n)))))
  (pari-p 4))    ; -> T
\end{lstlisting}

\textbf{\fn{progn}, \fn{prog1}, \fn{prog2}:} valutano sequenze di espressioni:\nopagebreak
\begin{lstlisting}
(progn e1 e2 e3)   ; ritorna e3
(prog1 e1 e2 e3)   ; ritorna e1
(prog2 e1 e2 e3)   ; ritorna e2
\end{lstlisting}

\end{refbox}

% =============================================================
% MACRO-AREA 2: CONTROLLO DEL FLUSSO
% =============================================================
\addcontentsline{toc}{section}{Controllo del flusso}

% --- Condizioni ---
\begin{refbox}[Condizioni]{csOrange}
\label{sec:condizioni}

\begin{lstlisting}
; if: ramo singolo (else opzionale)
(if (> x 0) 'positivo 'non-positivo)

; cond: rami multipli, t e' il default
(cond
  ((< x 0) 'negativo)
  ((= x 0) 'zero)
  (t        'positivo))

; when/unless: no ramo else, corpo multiplo
(when (> x 0)
  (print "positivo")
  (+ x 1))

(unless (null lst)
  (car lst))

; case: switch su valore (eql), otherwise = default
(case giorno
  ((1) "lunedi")
  ((2) "martedi")
  (otherwise "altro"))

; and/or come condizioni: ritornano il valore determinante
(and (> x 0) x)      ; -> x se positivo, NIL altrimenti
(or  val-a val-b)    ; -> val-a se non-NIL, altrimenti val-b
\end{lstlisting}

\end{refbox}

% --- Iterazione ---
\begin{refbox}[Iterazione]{csOrange}
\label{sec:iterazione}

\textbf{\fn{dotimes} e \fn{dolist}:}\nopagebreak
\begin{lstlisting}
(dotimes (i 5)          ; i da 0 a 4
  (print i))

(dolist (x '(a b c) 'fatto)  ; terzo arg = valore di ritorno
  (print x))
\end{lstlisting}

\textbf{\fn{loop} — forme comuni:}\nopagebreak
\begin{lstlisting}
; iterazione su lista
(loop for x in '(1 2 3)
      collect (* x x))          ; -> (1 4 9)

; iterazione numerica
(loop for i from 1 to 5
      sum i)                     ; -> 15

; passo personalizzato
(loop for i from 0 below 10 by 2
      collect i)                 ; -> (0 2 4 6 8)

; filtro
(loop for x in '(1 2 3 4 5)
      when (evenp x)
      collect x)                 ; -> (2 4)

; con accumulatore esplicito
(loop with acc = 0
      for x in '(1 2 3 4)
      do (incf acc x)
      finally (return acc))      ; -> 10

; loop infinito con uscita esplicita
(loop
  (when (done-p stato) (return risultato))
  (setf stato (aggiorna stato)))
\end{lstlisting}

\end{refbox}

% --- Ricorsione ---
\begin{refbox}[Ricorsione]{csOrange}
\label{sec:ricorsione}

\textbf{Schema base su lista:}\nopagebreak
\begin{lstlisting}
(defun lunghezza (lst)
  (if (null lst)
      0                                ; caso base
      (1+ (lunghezza (cdr lst)))))     ; passo ricorsivo
\end{lstlisting}

\textbf{Tail recursion con accumulatore} (evita stack overflow):\nopagebreak
\begin{lstlisting}
(defun lunghezza (lst)
  (labels ((iter (lst acc)
             (if (null lst)
                 acc
                 (iter (cdr lst) (1+ acc)))))
    (iter lst 0)))
\end{lstlisting}

Una chiamata è \textit{tail-recursive} se è l'ultima operazione
prima del ritorno. Allegro CL può ottimizzare la tail recursion (TCO),
eliminando i frame dallo stack — ma questo non è garantito
dallo standard ANSI.

\begin{refnote}[Schema ricorsione su lista]
\begin{tabular}{@{}ll@{}}
\textbf{Caso base:} & \fn{(null lst)} \\
\textbf{Caso ricorsivo:} & opera su \fn{(car lst)}, continua su \fn{(cdr lst)} \\
\textbf{Con accumulatore:} & aggiunge parametro \fn{acc}, inizializzato nel wrapper \\
\end{tabular}
\end{refnote}

\end{refbox}

% =============================================================
% MACRO-AREA 3: LISTE E SEQUENZE
% =============================================================
\addcontentsline{toc}{section}{Liste e sequenze}

% --- Liste ---
\begin{refbox}[Liste]{csGreen}
\label{sec:liste}

Una lista è una catena di \textit{cons-cell}. Ogni cella contiene
un valore (\fn{car}) e il puntatore al resto (\fn{cdr}).
\fn{NIL} termina la lista.

\begin{lstlisting}
'(a b c)           ; lista letterale (= (cons 'a (cons 'b (cons 'c nil))))
(cons 'a '(b c))   ; -> (A B C)
(car '(a b c))     ; -> A
(cdr '(a b c))     ; -> (B C)
\end{lstlisting}

\textbf{Coppia puntata} (cons non terminante in NIL):\nopagebreak
\begin{lstlisting}
(cons 'a 'b)       ; -> (A . B)  lista impropria
\end{lstlisting}

\textbf{Accesso:}\nopagebreak
\begin{lstlisting}
(first lst)  (rest lst)             ; = car, cdr
(second lst) (third lst)            ; cadr, caddr
(nth 2 lst)                         ; terzo elemento (0-based)
(nthcdr 2 lst)                      ; cdr applicato 2 volte
(last lst)                          ; ultima cons-cell (non elemento!)
(butlast lst)                       ; tutti tranne l'ultimo
\end{lstlisting}

\textbf{Costruzione e modifica:}\nopagebreak
\begin{lstlisting}
(list 'a 'b 'c)             ; -> (A B C)
(append '(a b) '(c d))      ; -> (A B C D)  (non distruttivo)
(reverse '(a b c))          ; -> (C B A)    (non distruttivo)
(nreverse (list 'a 'b 'c))  ; -> (C B A)    (distruttivo!)
(copy-list lst)             ; copia superficiale
\end{lstlisting}

\textbf{Push e pop:}\nopagebreak
\begin{lstlisting}
(push 'x lst)    ; aggiunge in testa (modifica lst)
(pop lst)        ; rimuove e restituisce la testa
\end{lstlisting}

\textbf{Ricerca e test:}\nopagebreak
\begin{lstlisting}
(length '(a b c))           ; -> 3
(member 'b '(a b c))        ; -> (B C)  (non solo T!)
(member 'b '(a b c) :test #'equal)  ; con comparatore custom
(find 'b '(a b c))          ; -> B
(position 'b '(a b c))      ; -> 1
\end{lstlisting}

\textbf{Sostituzione:}\nopagebreak
\begin{lstlisting}
(subst 'x 'b '(a b (b c)))  ; -> (A X (X C))  ricorsivo
(remove 'b '(a b b c))      ; -> (A C)
(remove-duplicates '(a b a c b)) ; -> (A C B)
\end{lstlisting}

\end{refbox}

% --- Predicati ---
\begin{refbox}[Predicati]{csGreen}
\label{sec:predicati}

\textbf{Predicati di tipo:}\nopagebreak
\begin{lstlisting}
(null '())       ; -> T   (lista vuota o NIL)
(atom 'a)        ; -> T   (atomo = non-cons)
(consp '(a))     ; -> T   (cons non vuota)
(listp '(a))     ; -> T   (lista, anche NIL)
(numberp 3)      ; -> T
(integerp 3)     ; -> T
(floatp 3.0)     ; -> T
(symbolp 'foo)   ; -> T
(stringp "x")    ; -> T
(functionp #'+)  ; -> T
(characterp #\a) ; -> T
\end{lstlisting}

\textbf{Uguaglianza:}\nopagebreak
\begin{lstlisting}
(eq 'a 'a)              ; T -- identita' oggetto (simboli, T, NIL)
(eql 3 3)               ; T -- eq + numeri e caratteri
(equal '(a b) '(a b))   ; T -- strutturale (ricorsiva)
(equalp "Ciao" "ciao")  ; T -- strutturale + case-insensitive
\end{lstlisting}

\begin{refnote}[Quando usare quale uguaglianza]
\begin{tabular}{@{}ll@{}}
Numeri:         & \fn{=} \\
Simboli:        & \fn{eq} \\
Caratteri:      & \fn{eql} \\
Liste/strutture: & \fn{equal} \\
Stringhe (no case): & \fn{equalp} o \fn{string-equal} \\
\end{tabular}
\end{refnote}

\end{refbox}

% --- Sequenze ---
\begin{refbox}[Sequenze (liste, vettori, stringhe)]{csGreen}
\label{sec:sequenze}

Molte funzioni funzionano su qualsiasi sequenza:

\begin{lstlisting}
(elt '(a b c) 1)             ; -> B  (lista)
(elt "hello" 1)              ; -> #\e (stringa)
(elt #(10 20 30) 1)          ; -> 20 (vettore)

(length "hello")             ; -> 5
(subseq "hello" 1 3)         ; -> "el"
(concatenate 'list '(a) '(b)) ; -> (A B)
(concatenate 'string "fo" "o") ; -> "foo"
(position #\l "hello")       ; -> 2
(count #\l "hello")          ; -> 2
(remove #\l "hello")         ; -> "heo"

; coerce converte tra tipi di sequenza
(coerce "abc" 'list)         ; -> (#\a #\b #\c)
(coerce '(#\a #\b) 'string)  ; -> "ab"
(coerce #(1 2 3) 'list)      ; -> (1 2 3)
\end{lstlisting}

\end{refbox}

% --- Destructuring ---
\begin{refbox}[Destructuring]{csGreen}
\label{sec:destructuring}

\fn{destructuring-bind} lega variabili alla struttura di una lista:

\begin{lstlisting}
(destructuring-bind (a b c) '(1 2 3)
  (+ a b c))             ; -> 6

(destructuring-bind (head . tail) '(1 2 3)
  tail)                  ; -> (2 3)

(destructuring-bind (a (b c) d) '(1 (2 3) 4)
  (list a b c d))        ; -> (1 2 3 4)
\end{lstlisting}

\end{refbox}

% =============================================================
% MACRO-AREA 4: STRUTTURE DATI
% =============================================================
\addcontentsline{toc}{section}{Strutture dati}

% --- Alist e Plist ---
\begin{refbox}[Association list e Property list]{csTeal}

\textbf{Association list (alist):} lista di coppie \fn{(chiave~.~valore)}:
\label{sec:alist}
\begin{lstlisting}
(defparameter *rubrica*
  '((mario . "333-1111")
    (luigi . "333-2222")))

(assoc 'mario *rubrica*)          ; -> (MARIO . "333-1111")
(cdr (assoc 'mario *rubrica*))    ; -> "333-1111"
(rassoc "333-1111" *rubrica* :test #'equal) ; cerca per valore
(cons '(anna . "333-3333") *rubrica*)  ; aggiunge coppia in testa
(acons 'anna "333-3333" *rubrica*)    ; equivalente, più leggibile
\end{lstlisting}

\textbf{Property list (plist):} chiavi e valori alternati nella stessa lista:\nopagebreak
\begin{lstlisting}
(defparameter *p* '(:nome "Mario" :eta 23))
(getf *p* :nome)           ; -> "Mario"
(getf *p* :via "n/d")     ; -> "n/d"  (default se assente)
(setf (getf *p* :eta) 24)  ; modifica
\end{lstlisting}

\begin{refnote}[Alist vs Plist]
\begin{tabular}{@{}lll@{}}
\textbf{} & \textbf{Alist} & \textbf{Plist} \\
Struttura & \fn{((k . v)...)} & \fn{(k v k v...)} \\
Ricerca   & \fn{assoc}  & \fn{getf} \\
Uso tipico & dati associativi & proprietà oggetti \\
\end{tabular}
\end{refnote}

\end{refbox}

% --- Hash table ---
\begin{refbox}[Hash table]{csTeal}
\label{sec:hashtable}

\begin{lstlisting}
; creazione (:test specifica il comparatore di chiavi)
(make-hash-table)                    ; default: eql
(make-hash-table :test #'equal)      ; chiavi stringa/lista

; lettura (due valori: valore e flag "trovato?")
(gethash 'foo ht)                    ; -> NIL, NIL  (assente)
(gethash 'foo ht "default")          ; -> "default", NIL

; scrittura e rimozione
(setf (gethash 'foo ht) 42)
(remhash 'foo ht)

; utilità
(hash-table-count ht)                ; numero di chiavi presenti
(clrhash ht)                         ; svuota la tabella

; iterazione
(maphash (lambda (k v)
           (format t "~A -> ~A~%" k v))
         ht)
\end{lstlisting}

\begin{refnote}[Ricevere entrambi i valori di gethash]
\begin{lstlisting}
(multiple-value-bind (valore trovato)
    (gethash 'foo ht)
  (if trovato
      (format t "Trovato: ~A~%" valore)
      (format t "Assente~%")))
\end{lstlisting}
\end{refnote}

\end{refbox}

% --- Array ---
\begin{refbox}[Array e vettori]{csTeal}
\label{sec:array}

\begin{lstlisting}
#(10 20 30)          ; vettore letterale (immutabile)

; creazione (sempre con :initial-element per evitare letture indefinite)
(make-array 5 :initial-element 0)           ; vettore di 5 zeri
(make-array 5 :initial-element nil)         ; vettore di 5 NIL
(make-array '(3 3) :initial-element 0)      ; matrice 3x3

; accesso e modifica
(aref v 2)                                  ; terzo elemento
(aref m 1 2)                                ; riga 1, colonna 2
(setf (aref v 0) 99)

; dimensioni
(length v)                                  ; lunghezza (vettori 1D)
(array-dimensions m)                        ; -> (3 3)
(array-dimension m 0)                       ; -> 3 (numero righe)

; vettore dinamico (adjustable)
(defparameter *v*
  (make-array 0 :fill-pointer 0 :adjustable t))
(vector-push-extend 42 *v*)                 ; aggiunge in coda
(vector-pop *v*)                            ; rimuove dall'ultima posizione
\end{lstlisting}

\end{refbox}

% --- Defstruct ---
\begin{refbox}[Strutture con defstruct]{csTeal}
\label{sec:strutture}

\fn{defstruct} genera automaticamente costruttore, accessori,
predicato e funzione di copia:

\begin{lstlisting}
(defstruct persona
  nome
  (eta 18)        ; slot con valore di default
  citta)

; costruzione
(defparameter *p*
  (make-persona :nome "Mario" :eta 23 :citta "Roma"))

; accesso agli slot
(persona-nome *p*)           ; -> "Mario"
(persona-eta *p*)            ; -> 23

; modifica
(setf (persona-eta *p*) 24)

; predicato e copia
(persona-p *p*)              ; -> T
(copy-persona *p*)           ; copia superficiale
\end{lstlisting}

\begin{refnote}[Cosa genera defstruct automaticamente]
\begin{tabular}{@{}ll@{}}
\fn{make-nome}    & costruttore \\
\fn{nome-slot}    & lettore per ogni slot \\
\fn{(setf (nome-slot ...) ...)} & modificatore \\
\fn{nome-p}       & predicato di tipo \\
\fn{copy-nome}    & copia superficiale \\
\end{tabular}
\end{refnote}

\end{refbox}

% --- Multiple values ---
\begin{refbox}[Valori multipli]{csTeal}
\label{sec:multiplevalues}

Una funzione può restituire più di un valore con \fn{values}:

\begin{lstlisting}
(values 3 4)                ; restituisce due valori

; cattura con multiple-value-bind
(multiple-value-bind (q r)
    (floor 17 5)
  (format t "~A resto ~A~%" q r))  ; "3 resto 2"

; raccoglie i valori in una lista
(multiple-value-list (floor 17 5)) ; -> (3 2)

; passa i valori multipli come argomenti
(multiple-value-call #'+ (values 1 2) (values 3 4)) ; -> 10
\end{lstlisting}

\end{refbox}

% =============================================================
% MACRO-AREA 5: ASTRAZIONE
% =============================================================
\addcontentsline{toc}{section}{Funzioni, closures e macro}

% --- Higher-order ---
\begin{refbox}[Funzioni higher-order]{csPurple}
\label{sec:higherorder}

\fn{\#'} ottiene l'oggetto funzione; \fn{lambda} crea funzioni anonime:

\begin{lstlisting}
#'quadrato           ; oggetto funzione (= (function quadrato))
(lambda (x) (* x x)) ; funzione anonima

(funcall #'quadrato 5)         ; -> 25
(funcall (lambda (x) (* x x)) 5) ; -> 25
(apply #'+ '(1 2 3))           ; -> 6
(apply #'max 3 '(1 5 2))       ; -> 5  (argomenti misti)
\end{lstlisting}

\textbf{Funzioni sulle liste:}\nopagebreak
\begin{lstlisting}
(mapcar #'quadrato '(1 2 3 4))          ; -> (1 4 9 16)
(mapcar #'+ '(1 2 3) '(10 20 30))       ; -> (11 22 33)
(mapcan #'list '(1 2 3))                ; -> (1 2 3)
(remove-if #'evenp '(1 2 3 4 5))        ; -> (1 3 5)
(remove-if-not #'oddp '(1 2 3 4 5))     ; -> (1 3 5)
(find-if #'evenp '(1 3 4 5))            ; -> 4
(find-if-not #'evenp '(1 3 4 5))        ; -> 1
(every #'numberp '(1 2 3))              ; -> T
(some  #'null '(1 nil 3))               ; -> T
(count-if #'evenp '(1 2 3 4))           ; -> 2
(reduce #'+ '(1 2 3 4))                 ; -> 10
(reduce #'+ '() :initial-value 0)       ; -> 0  (lista vuota)
(sort (copy-list '(3 1 2)) #'<)         ; -> (1 2 3)  (distruttivo!)
(stable-sort lst #'<)                   ; preserva ordine degli uguali
\end{lstlisting}

\begin{refnote}[sort è distruttivo]
\fn{sort} può modificare la lista originale. Usare sempre
\fn{(sort (copy-list lst) \#'<)} se si vuole preservare \fn{lst}.
\end{refnote}

\end{refbox}

% --- Closures ---
\begin{refbox}[Closures]{csPurple}
\label{sec:closures}

Una \textit{closure} è una funzione che cattura le variabili del
proprio ambiente lessicale. È il meccanismo principale per
incapsulare stato privato:

\begin{lstlisting}
(defun crea-contatore ()
  (let ((n 0))
    (lambda ()
      (incf n)     ; n e' privata, non accessibile dall'esterno
      n)))

(defparameter *c1* (crea-contatore))
(defparameter *c2* (crea-contatore))  ; stato indipendente da *c1*

(funcall *c1*)   ; -> 1
(funcall *c1*)   ; -> 2
(funcall *c2*)   ; -> 1  (contatore separato)
\end{lstlisting}

\textbf{Closure con più operazioni:}\nopagebreak
\begin{lstlisting}
(defun crea-generatore (start step)
  (let ((curr start))
    (lambda ()
      (prog1 curr        ; ritorna curr, poi lo incrementa
        (incf curr step)))))

(defparameter *gen* (crea-generatore 0 2))
(funcall *gen*)  ; -> 0
(funcall *gen*)  ; -> 2
(funcall *gen*)  ; -> 4
\end{lstlisting}

\end{refbox}

% --- Macro ---
\begin{refbox}[Macro]{csPurple}
\label{sec:macro}

Le macro trasformano codice \textit{prima} della valutazione
(espansione a compile-time). Ricevono le forme non valutate.

\begin{lstlisting}
(defmacro my-when (test &body body)
  `(if ,test
       (progn ,@body)
       nil))

; uso
(my-when (> x 0)
  (print "positivo")
  x)

; verifica dell'espansione
(macroexpand-1 '(my-when (> x 0) (print x)))
; -> (IF (> X 0) (PROGN (PRINT X)) NIL)
\end{lstlisting}

\begin{itemize}[leftmargin=1.5em, itemsep=2pt]
  \item \fn{`} (backquote) --- crea una lista quasi-letterale
  \item \fn{,} (unquote) --- valuta l'espressione nel contesto del chiamante
  \item \fn{,@} (splicing) --- inserisce gli elementi di una lista
\end{itemize}

\textbf{Igiene con \fn{gensym}:} evita la cattura accidentale di variabili:\nopagebreak
\begin{lstlisting}
(defmacro my-swap (a b)
  (let ((tmp (gensym "TMP")))  ; genera simbolo unico
    `(let ((,tmp ,a))
       (setf ,a ,b)
       (setf ,b ,tmp))))
\end{lstlisting}

\end{refbox}

% =============================================================
% MACRO-AREA 6: I/O, ERRORI, CLOS, PATTERN, CONVENZIONI
% =============================================================
\addcontentsline{toc}{section}{I/O e gestione errori}

% --- I/O ---
\begin{refbox}[Input/Output]{csTeal}
\label{sec:io}

\textbf{Lettura:}\nopagebreak
\begin{lstlisting}
(read)             ; legge una S-expression da stdin
(read-line)        ; legge una riga come stringa
(read-char)        ; legge un singolo carattere
\end{lstlisting}

\textbf{Stampa:}\nopagebreak
\begin{lstlisting}
(print x)          ; stampa con quote + newline prima
(princ x)          ; stampa senza quote (leggibile)
(prin1 x)          ; stampa con quote, senza newline prima
(terpri)           ; solo newline
(fresh-line)       ; newline solo se non gia' a inizio riga
\end{lstlisting}

\textbf{format:} funzione di output formattato:\nopagebreak
\begin{lstlisting}
(format t "~A~%"     x)    ; generico + newline
(format t "~D~%"     n)    ; intero decimale
(format t "~,2F~%"   f)    ; float con 2 decimali
(format t "~S~%"     x)    ; con quote (come prin1)
(format t "~{~A ~}~%" lst) ; itera su lista
(format t "~A e ~A~%" a b) ; argomenti multipli

; format su stringa (nil = non stampare, ritorna stringa)
(format nil "Risultato: ~A" val)  ; -> "Risultato: 42"
\end{lstlisting}

\textbf{Direttive format principali:}

\smallskip
\begin{tabular}{@{}l@{\hspace{8pt}}l@{\hspace{16pt}}l@{\hspace{8pt}}l@{}}
\fn{\textasciitilde A} & generico (aesthetic) &
\fn{\textasciitilde S} & con escape \\
\fn{\textasciitilde D} & intero dec &
\fn{\textasciitilde F} & float \\
\fn{\textasciitilde\%} & newline &
\fn{\textasciitilde\&} & fresh-line \\
\fn{\textasciitilde\textasciitilde} & tilde letterale &
\fn{\textasciitilde nT} & tab a colonna n \\
\end{tabular}

\medskip
\textbf{File:}\nopagebreak
\begin{lstlisting}
; lettura riga per riga
(with-open-file (s "dati.txt" :direction :input)
  (loop for riga = (read-line s nil nil)
        while riga
        collect riga))

; scrittura
(with-open-file (s "out.txt"
                   :direction :output
                   :if-exists :supersede)
  (format s "Ciao~%"))
\end{lstlisting}

\end{refbox}

% --- Gestione errori ---
\begin{refbox}[Gestione errori]{csTeal}
\label{sec:errori}

\begin{lstlisting}
; segnalare un errore
(error "Valore non valido: ~A" x)
(warn  "Attenzione: ~A" msg)      ; segnalazione non fatale

; gestire errori con handler-case
(handler-case
    (/ 10 0)
  (division-by-zero (e)
    (format t "Divisione per zero~%"))
  (error (e)
    (format t "Errore: ~A~%" e)))

; ignorare errori (ritorna NIL + condizione)
(ignore-errors
  (/ 10 0))    ; -> NIL, #<DIVISION-BY-ZERO ...>

; handler-bind (non unwind lo stack)
(handler-bind
    ((error (lambda (e)
              (format t "Intercettato: ~A~%" e))))
  (/ 10 0))
\end{lstlisting}

\end{refbox}

% =============================================================
\addcontentsline{toc}{section}{CLOS (programmazione a oggetti)}
% =============================================================

\begin{refbox}[CLOS --- Object-Oriented \normalfont\small(argomento avanzato)]{csPurple}
\label{sec:clos}

Il \textbf{Common Lisp Object System} (CLOS) è il sistema di
classi integrato in ANSI Common Lisp. Supporta ereditarietà
multipla e metodi polimorfici.

\begin{lstlisting}
; definizione di classe
(defclass animale ()                       ; () = nessuna superclasse
  ((nome  :initarg :nome  :accessor animale-nome)
   (verso :initarg :verso :reader   animale-verso)))

; sottoclasse
(defclass cane (animale)
  ())

; generic function e metodo
(defgeneric descrivi (a))

(defmethod descrivi ((a animale))
  (format t "~A dice ~A~%"
          (animale-nome a)
          (animale-verso a)))

; specializzazione nella sottoclasse
(defmethod descrivi ((c cane))
  (format t "Il cane ~A abbaia!~%" (animale-nome c)))

; creazione istanza
(defparameter *fido*
  (make-instance 'cane :nome "Fido" :verso "Bau"))

(descrivi *fido*)       ; chiama il metodo specializzato per CANE

; predicato di tipo
(typep *fido* 'cane)    ; -> T
(typep *fido* 'animale) ; -> T  (ereditarieta')
\end{lstlisting}

\fn{defmethod} può anche definire metodi \fn{:before},
\fn{:after}, \fn{:around} per decorare il comportamento
senza sovrascrivere il metodo principale.

\end{refbox}

% =============================================================
\addcontentsline{toc}{section}{Pattern e convenzioni}
% =============================================================

\begin{refbox}[Pattern essenziali]{csOrange}
\label{sec:pattern}

\textbf{Elaborazione ricorsiva di lista:}\nopagebreak
\begin{lstlisting}
(defun trasforma-lista (lst)
  (if (null lst)
      '()
      (cons (trasforma (car lst))
            (trasforma-lista (cdr lst)))))
\end{lstlisting}

\textbf{Accumulatore tail-recursive:}\nopagebreak
\begin{lstlisting}
(defun somma (lst)
  (labels ((iter (lst acc)
             (if (null lst)
                 acc
                 (iter (cdr lst) (+ acc (car lst))))))
    (iter lst 0)))
\end{lstlisting}

\textbf{Memoization:}\nopagebreak
\begin{lstlisting}
(defun memoize (f)
  (let ((cache (make-hash-table :test #'equal)))
    (lambda (&rest args)
      (multiple-value-bind (val trovato)
          (gethash args cache)
        (if trovato val
            (setf (gethash args cache)
                  (apply f args)))))))
\end{lstlisting}

\textbf{Backtracking:}\nopagebreak
\begin{lstlisting}
(defun risolvi (stato)
  (cond
    ((soluzione-p stato) (list stato))
    ((vicolo-cieco-p stato) nil)
    (t (loop for mossa in (mosse-valide stato)
             append (risolvi (applica mossa stato))))))
\end{lstlisting}

\end{refbox}

\begin{refbox}[Convenzioni di stile]{csBlue}
\label{sec:stile}

\begin{itemize}[leftmargin=1.5em, itemsep=3pt]
  \item Nomi con \textbf{kebab-case}: \fn{mia-funzione}, \fn{conta-elementi}
  \item Variabili globali dinamiche con \textbf{earmuffs}: \fn{*contatore*}
  \item Costanti con \textbf{più segni}: \fn{+max-dimensione+}
  \item Predicati terminano in \textbf{-p}: \fn{null}, \fn{evenp}, \fn{palindromo-p}
  \item Reimplementazioni di built-in: suffisso \fn{*}: \fn{length*}, \fn{reverse*}
  \item Ogni funzione pubblica dovrebbe avere una \textit{docstring}
  \item Funzioni piccole e composabili; evitare stato globale
  \item Gestire esplicitamente i casi limite: lista vuota, \fn{nil}, zero
\end{itemize}

\end{refbox}

\begin{refbox}[Package]{csBlue}
\label{sec:package}

\begin{lstlisting}
(defpackage :mio-progetto
  (:use :cl)
  (:export #:funzione-pubblica
           #:altra-funzione))

(in-package :mio-progetto)
\end{lstlisting}

Durante la fase di apprendimento è sufficiente lavorare nel
package \fn{cl-user} (il default al REPL).
Per esercizi multi-file, definire un package per progetto
ed esportare solo le funzioni pubbliche.

\end{refbox}

% =============================================================
%  02_esercizi.tex — Tutti gli esercizi (148 totali, 10 capitoli)
%  \part{Esercizi} e la legenda sono dichiarati in main.tex
% =============================================================

% =============================================================
\setcounter{section}{0}
\section{Fondamenti}
% =============================================================

\noindent
In questo capitolo si prende confidenza con la sintassi di
Common Lisp: S-expression, \fn{defun}, aritmetica, predicati,
condizionali (\fn{if}, \fn{cond}), binding locale (\fn{let},
\fn{let*}) e iterazione base (\fn{dotimes}, \fn{loop}).

\bigskip

\label{ex:1}
\begin{essbox}{1}{Somma}{1}
\exlabel{Implementare:}
\nopagebreak
\begin{lstlisting}
(add a b)
\end{lstlisting}
\begin{lstlisting}
(add 3 5)      ; -> 8
(add -2 7)     ; -> 5
(add 1.5 2.5)  ; -> 4.0
\end{lstlisting}


\small\color{black!55}\hyperref[sol:1]{$\rightarrow$ Soluzione}\normalsize\color{black}
\end{essbox}

\label{ex:2}
\begin{stdbox}{2}{Geometria del cerchio}{1}
\exlabel{Implementare:}
\nopagebreak
\begin{lstlisting}
(circle-area radius)
(circle-circumference radius)
\end{lstlisting}
Usare la costante \fn{pi} per $\pi$.
\begin{lstlisting}
(circle-area 2)           ; -> 12.566...
(circle-circumference 2)  ; -> 12.566...
\end{lstlisting}


\small\color{black!55}\hyperref[sol:2]{$\rightarrow$ Soluzione}\normalsize\color{black}
\end{stdbox}

\label{ex:3}
\begin{stdbox}{3}{Conversione Celsius/Fahrenheit}{1}
\exlabel{Implementare:}
\nopagebreak
\begin{lstlisting}
(celsius->fahrenheit c)
(fahrenheit->celsius f)
\end{lstlisting}
Formule: $F = C \times 9/5 + 32$;\quad $C = (F - 32) \times 5/9$.


\small\color{black!55}\hyperref[sol:3]{$\rightarrow$ Soluzione}\normalsize\color{black}
\end{stdbox}

\label{ex:4}
\begin{essbox}{4}{Ipotenusa}{1}
\exlabel{Implementare:}
\nopagebreak
\begin{lstlisting}
(ipotenusa a b)
\end{lstlisting}
\begin{lstlisting}
(ipotenusa 3 4)   ; -> 5.0
\end{lstlisting}
\vincolo{usare \fn{let} per i quadrati intermedi.}


\small\color{black!55}\hyperref[sol:4]{$\rightarrow$ Soluzione}\normalsize\color{black}
\end{essbox}

\label{ex:5}
\begin{colbox}{5}{Distanza euclidea}{1}
\exlabel{Implementare:}
\nopagebreak
\begin{lstlisting}
(distanza-euclidea x1 y1 x2 y2)
\end{lstlisting}
\begin{lstlisting}
(distanza-euclidea 0 0 3 4)   ; -> 5.0
\end{lstlisting}
\vincolo{usare \fn{let} per i valori intermedi.}


\small\color{black!55}\hyperref[sol:5]{$\rightarrow$ Soluzione}\normalsize\color{black}
\end{colbox}

\label{ex:6}
\begin{essbox}{6}{Pari o dispari}{1}
\exlabel{Implementare:}
\nopagebreak
\begin{lstlisting}
(pari-p n)
\end{lstlisting}
senza usare \fn{evenp}.
\begin{lstlisting}
(pari-p 8)    ; -> T
(pari-p 7)    ; -> NIL
\end{lstlisting}


\small\color{black!55}\hyperref[sol:6]{$\rightarrow$ Soluzione}\normalsize\color{black}
\end{essbox}

\label{ex:7}
\begin{colbox}{7}{Divisibilita'}{1}
\exlabel{Implementare:}
\nopagebreak
\begin{lstlisting}
(divisibile-p a b)
\end{lstlisting}
\begin{lstlisting}
(divisibile-p 12 4)   ; -> T
(divisibile-p 7  3)   ; -> NIL
\end{lstlisting}


\small\color{black!55}\hyperref[sol:7]{$\rightarrow$ Soluzione}\normalsize\color{black}
\end{colbox}

\label{ex:8}
\begin{essbox}{8}{Segno}{1}
\exlabel{Implementare:}
\nopagebreak
\begin{lstlisting}
(segno n)
\end{lstlisting}
che restituisce \fn{'positivo}, \fn{'negativo} o \fn{'zero}.
\begin{lstlisting}
(segno -7)   ; -> NEGATIVO
(segno 0)    ; -> ZERO
(segno 12)   ; -> POSITIVO
\end{lstlisting}
\vincolo{usare \fn{cond}.}


\small\color{black!55}\hyperref[sol:8]{$\rightarrow$ Soluzione}\normalsize\color{black}
\end{essbox}

\label{ex:9}
\begin{stdbox}{9}{Massimo di tre}{1}
\exlabel{Implementare} \fn{(max-tre a b c)} senza usare \fn{max}.
\nopagebreak
\begin{lstlisting}
(max-tre 3 7 2)   ; -> 7
\end{lstlisting}
\vincolo{usare \fn{cond}.}


\small\color{black!55}\hyperref[sol:9]{$\rightarrow$ Soluzione}\normalsize\color{black}
\end{stdbox}

\label{ex:10}
\begin{colbox}{10}{Triangolo valido}{1}
\exlabel{Implementare} \fn{(triangolo-valido-p a b c)}.
Ogni lato deve essere minore della somma degli altri due.
\begin{lstlisting}
(triangolo-valido-p 3 4 5)   ; -> T
(triangolo-valido-p 1 2 10)  ; -> NIL
\end{lstlisting}


\small\color{black!55}\hyperref[sol:10]{$\rightarrow$ Soluzione}\normalsize\color{black}
\end{colbox}

\label{ex:11}
\begin{essbox}{11}{Anno bisestile}{2}
\exlabel{Implementare:}
\nopagebreak
\begin{lstlisting}
(bisestile-p anno)
\end{lstlisting}
Divisibile per 4, eccetto multipli di 100 che non siano anche di 400.
\begin{lstlisting}
(bisestile-p 2000)   ; -> T
(bisestile-p 1900)   ; -> NIL
(bisestile-p 2024)   ; -> T
\end{lstlisting}
\vincolo{usare \fn{and}/\fn{or} senza \fn{cond}.}


\small\color{black!55}\hyperref[sol:11]{$\rightarrow$ Soluzione}\normalsize\color{black}
\end{essbox}

\label{ex:12}
\begin{stdbox}{12}{Giorni del mese}{2}
\exlabel{Implementare} \fn{(giorni-mese mese anno)} con \fn{mese} da 1 a 12.
\nopagebreak
\begin{lstlisting}
(giorni-mese 2 2024)   ; -> 29
(giorni-mese 2 2023)   ; -> 28
(giorni-mese 7 2024)   ; -> 31
\end{lstlisting}
Usare \fn{bisestile-p} per febbraio.


\small\color{black!55}\hyperref[sol:12]{$\rightarrow$ Soluzione}\normalsize\color{black}
\end{stdbox}

\label{ex:13}
\begin{colbox}{13}{Voto in lettere}{1}
\exlabel{Implementare} \fn{(voto->giudizio punteggio)}, punteggio 0--100.

90--100 $\to$ ``ottimo''; 75--89 $\to$ ``buono'';
60--74 $\to$ ``sufficiente''; $<60$ $\to$ ``insufficiente''.

\vincolo{usare \fn{cond}.}


\small\color{black!55}\hyperref[sol:13]{$\rightarrow$ Soluzione}\normalsize\color{black}
\end{colbox}

\label{ex:14}
\begin{colbox}{14}{Terna pitagorica}{1}
\exlabel{Implementare} \fn{(terna-pitagorica-p a b c)}.
\nopagebreak
\begin{lstlisting}
(terna-pitagorica-p 3 4 5)    ; -> T
(terna-pitagorica-p 1 2 3)    ; -> NIL
\end{lstlisting}


\small\color{black!55}\hyperref[sol:14]{$\rightarrow$ Soluzione}\normalsize\color{black}
\end{colbox}

\label{ex:15}
\begin{essbox}{15}{Radici di equazione di secondo grado}{2}
\exlabel{Implementare:}
\nopagebreak
\begin{lstlisting}
(radici-quadratiche a b c)
\end{lstlisting}
Restituire lista con le due radici reali, oppure \fn{NIL} se $\Delta < 0$.
\begin{lstlisting}
(radici-quadratiche 1 -5 6)   ; -> (3.0 2.0)
(radici-quadratiche 1  0 1)   ; -> NIL
\end{lstlisting}
\vincolo{usare \fn{let*} per discriminante e radici.}


\small\color{black!55}\hyperref[sol:15]{$\rightarrow$ Soluzione}\normalsize\color{black}
\end{essbox}

\label{ex:16}
\begin{stdbox}{16}{Calcolatrice simbolica}{2}
\exlabel{Implementare} \fn{(calcola op a b)} dove \fn{op} e' un simbolo
\fn{'+}, \fn{'-}, \fn{'*}, \fn{'/}.
\begin{lstlisting}
(calcola '+ 3 4)   ; -> 7
(calcola '* 5 6)   ; -> 30
\end{lstlisting}
Gestire la divisione per zero restituendo \fn{NIL}.


\small\color{black!55}\hyperref[sol:16]{$\rightarrow$ Soluzione}\normalsize\color{black}
\end{stdbox}

\label{ex:17}
\begin{essbox}{17}{Somma delle cifre}{2}
\exlabel{Implementare:}
\nopagebreak
\begin{lstlisting}
(somma-cifre n)
\end{lstlisting}
\begin{lstlisting}
(somma-cifre 1234)   ; -> 10
(somma-cifre 99)     ; -> 18
\end{lstlisting}
\textbf{Suggerimento:} estrarre le cifre con \fn{mod} e \fn{floor}.


\small\color{black!55}\hyperref[sol:17]{$\rightarrow$ Soluzione}\enspace\textcolor{black!30}{|}\enspace\hyperref[hint:17]{$\rightarrow$ Suggerimento}\normalsize\color{black}
\end{essbox}

\label{ex:18}
\begin{essbox}{18}{Conta fino a N}{1}
\exlabel{Implementare} \fn{(conta-fino-a n)} che stampa i numeri da 1 a \fn{n}.
\nopagebreak
\begin{lstlisting}
(conta-fino-a 4)
; 1
; 2
; 3
; 4
\end{lstlisting}
\vincolo{usare \fn{dotimes} o \fn{loop}.}


\small\color{black!55}\hyperref[sol:18]{$\rightarrow$ Soluzione}\normalsize\color{black}
\end{essbox}

\label{ex:19}
\begin{colbox}{19}{Tabellina}{1}
\exlabel{Implementare} \fn{(tabellina n)} che stampa la tabellina di \fn{n}
da 1 a 10.
\begin{lstlisting}
(tabellina 3)
; 3 x 1 = 3
; 3 x 2 = 6
; ...
\end{lstlisting}


\small\color{black!55}\hyperref[sol:19]{$\rightarrow$ Soluzione}\normalsize\color{black}
\end{colbox}

\label{ex:20}
\begin{colbox}{20}{Sconto}{1}
\exlabel{Implementare} \fn{(sconto prezzo)}:
$\geq 200$ $\to$ 20\%; $\geq 100$ $\to$ 10\%; altrimenti 0\%.
\begin{lstlisting}
(sconto 250)   ; -> 200.0
(sconto 120)   ; -> 108.0
(sconto 50)    ; -> 50
\end{lstlisting}


\small\color{black!55}\hyperref[sol:20]{$\rightarrow$ Soluzione}\normalsize\color{black}
\end{colbox}

\label{ex:21}
\begin{colbox}{21}{Paga settimanale}{1}
\exlabel{Implementare} \fn{(paga-settimanale ore tariffa)}.
Fino a 40 ore: tariffa normale.
Ore eccedenti: tariffa $\times$ 1.5.
\begin{lstlisting}
(paga-settimanale 35 10)   ; -> 350
(paga-settimanale 45 10)   ; -> 475
\end{lstlisting}


\small\color{black!55}\hyperref[sol:21]{$\rightarrow$ Soluzione}\normalsize\color{black}
\end{colbox}

% =============================================================
\setcounter{section}{1}
\section{Liste}
% =============================================================

\noindent
Le liste sono la struttura dati fondamentale di Common Lisp.
In questo capitolo si implementano da zero le operazioni
principali usando \fn{car}, \fn{cdr}, \fn{cons} e la ricorsione.
\textbf{Non usare} le funzioni built-in che risolvono direttamente
l'esercizio, salvo diversa indicazione.

\bigskip

\label{ex:22}
\begin{essbox}{22}{Costruire liste con cons e list}{1}
\textbf{Parte A.} Costruire \fn{'(a (b c) d)} usando solo \fn{cons}.

\textbf{Parte B.} Costruire la stessa lista con \fn{list}.

Verificare che le due espressioni siano \fn{equal}.


\small\color{black!55}\hyperref[sol:22]{$\rightarrow$ Soluzione}\normalsize\color{black}
\end{essbox}

\label{ex:23}
\begin{stdbox}{23}{Accesso con car e cdr}{1}
Data \fn{'(a b c d e)}, estrarre il terzo elemento
usando solo \fn{car} e \fn{cdr} (non \fn{nth}, \fn{third}, ecc.).
\begin{lstlisting}
; risultato atteso: C
\end{lstlisting}


\small\color{black!55}\hyperref[sol:23]{$\rightarrow$ Soluzione}\normalsize\color{black}
\end{stdbox}

\label{ex:24}
\begin{essbox}{24}{Primo e resto}{1}
\exlabel{Implementare:}
\nopagebreak
\begin{lstlisting}
(first-element lst)
(rest-elements lst)
\end{lstlisting}
Gestire la lista vuota restituendo \fn{NIL}.
\begin{lstlisting}
(first-element '(a b c))    ; -> A
(rest-elements '(a b c))    ; -> (B C)
(first-element '())         ; -> NIL
\end{lstlisting}


\small\color{black!55}\hyperref[sol:24]{$\rightarrow$ Soluzione}\normalsize\color{black}
\end{essbox}

\label{ex:25}
\begin{colbox}{25}{Primo e ultimo}{1}
\exlabel{Implementare} \fn{(primo-e-ultimo lst)} che restituisce una lista
con il primo e l'ultimo elemento.
\begin{lstlisting}
(primo-e-ultimo '(a b c d))   ; -> (A D)
(primo-e-ultimo '(a))         ; -> (A A)
\end{lstlisting}
Gestire la lista vuota.


\small\color{black!55}\hyperref[sol:25]{$\rightarrow$ Soluzione}\normalsize\color{black}
\end{colbox}

\label{ex:26}
\begin{colbox}{26}{Scambia i primi due}{1}
\exlabel{Implementare} \fn{(scambia-primi-due lst)}.
\nopagebreak
\begin{lstlisting}
(scambia-primi-due '(a b c d))   ; -> (B A C D)
\end{lstlisting}
Gestire liste con meno di due elementi.


\small\color{black!55}\hyperref[sol:26]{$\rightarrow$ Soluzione}\normalsize\color{black}
\end{colbox}

\label{ex:27}
\begin{essbox}{27}{Lunghezza}{2}
\exlabel{Implementare} \fn{(length* lst)} senza usare \fn{length}.
\nopagebreak
\begin{lstlisting}
(length* '())        ; -> 0
(length* '(a b c))   ; -> 3
\end{lstlisting}
\vincolo{ricorsione.}


\small\color{black!55}\hyperref[sol:27]{$\rightarrow$ Soluzione}\normalsize\color{black}
\end{essbox}

\label{ex:28}
\begin{stdbox}{28}{Somma degli elementi}{2}
\exlabel{Implementare} \fn{(sum-list lst)}.
\nopagebreak
\begin{lstlisting}
(sum-list '(1 2 3 4))   ; -> 10
(sum-list '())          ; -> 0
\end{lstlisting}
\vincolo{ricorsione.}


\small\color{black!55}\hyperref[sol:28]{$\rightarrow$ Soluzione}\normalsize\color{black}
\end{stdbox}

\label{ex:29}
\begin{colbox}{29}{Prodotto degli elementi}{2}
\exlabel{Implementare} \fn{(product-list lst)}.
\nopagebreak
\begin{lstlisting}
(product-list '(2 3 4))   ; -> 24
(product-list '())        ; -> 1
\end{lstlisting}


\small\color{black!55}\hyperref[sol:29]{$\rightarrow$ Soluzione}\normalsize\color{black}
\end{colbox}

\label{ex:30}
\begin{colbox}{30}{Massimo di una lista}{2}
\exlabel{Implementare} \fn{(max-list lst)}.
\nopagebreak
\begin{lstlisting}
(max-list '(4 8 2 9 1))   ; -> 9
\end{lstlisting}
Definire il comportamento per la lista vuota.


\small\color{black!55}\hyperref[sol:30]{$\rightarrow$ Soluzione}\normalsize\color{black}
\end{colbox}

\label{ex:31}
\begin{essbox}{31}{Appartenenza}{2}
\exlabel{Implementare} \fn{(member* item lst)} senza usare \fn{member}.
\nopagebreak
\begin{lstlisting}
(member* 'b '(a b c))   ; -> T
(member* 'x '(a b c))   ; -> NIL
\end{lstlisting}


\small\color{black!55}\hyperref[sol:31]{$\rightarrow$ Soluzione}\normalsize\color{black}
\end{essbox}

\label{ex:32}
\begin{stdbox}{32}{Conta occorrenze}{2}
\exlabel{Implementare} \fn{(count-element item lst)} senza usare \fn{count}.
\nopagebreak
\begin{lstlisting}
(count-element 'a '(a b a c a))   ; -> 3
\end{lstlisting}


\small\color{black!55}\hyperref[sol:32]{$\rightarrow$ Soluzione}\normalsize\color{black}
\end{stdbox}

\label{ex:33}
\begin{colbox}{33}{Posizione}{2}
\exlabel{Implementare} \fn{(position* item lst)}: indice 0-based
della prima occorrenza, oppure \fn{NIL}.
\begin{lstlisting}
(position* 'c '(a b c d))   ; -> 2
(position* 'x '(a b c d))   ; -> NIL
\end{lstlisting}


\small\color{black!55}\hyperref[sol:33]{$\rightarrow$ Soluzione}\normalsize\color{black}
\end{colbox}

\label{ex:34}
\begin{colbox}{34}{Ultimo elemento}{2}
\exlabel{Implementare} \fn{(last-element lst)} senza usare \fn{last}.
\nopagebreak
\begin{lstlisting}
(last-element '(a b c))   ; -> C
\end{lstlisting}


\small\color{black!55}\hyperref[sol:34]{$\rightarrow$ Soluzione}\normalsize\color{black}
\end{colbox}

\label{ex:35}
\begin{essbox}{35}{Incrementa tutti}{2}
\exlabel{Implementare} \fn{(incrementa-tutti lst)} che aggiunge 1
a ogni elemento.
\begin{lstlisting}
(incrementa-tutti '(1 2 3))   ; -> (2 3 4)
\end{lstlisting}
\vincolo{ricorsione; non usare \fn{mapcar}.}

\exnota{Obiettivo: capire il pattern di trasformazione, precursore di \fn{mapcar}.}


\small\color{black!55}\hyperref[sol:35]{$\rightarrow$ Soluzione}\normalsize\color{black}
\end{essbox}

\label{ex:36}
\begin{stdbox}{36}{Filtra i pari}{2}
\exlabel{Implementare} \fn{(solo-pari lst)}.
\nopagebreak
\begin{lstlisting}
(solo-pari '(1 2 3 4 5 6))   ; -> (2 4 6)
\end{lstlisting}
\vincolo{ricorsione; non usare \fn{remove-if}.}

\exnota{Obiettivo: capire il pattern di filtraggio, precursore di \fn{remove-if-not}.}


\small\color{black!55}\hyperref[sol:36]{$\rightarrow$ Soluzione}\normalsize\color{black}
\end{stdbox}

\label{ex:37}
\begin{stdbox}{37}{Elimina tutte le occorrenze}{2}
\exlabel{Implementare} \fn{(elimina-tutte item lst)} senza usare \fn{remove}.
\nopagebreak
\begin{lstlisting}
(elimina-tutte 'a '(a b a c a))   ; -> (B C)
\end{lstlisting}


\small\color{black!55}\hyperref[sol:37]{$\rightarrow$ Soluzione}\normalsize\color{black}
\end{stdbox}

\label{ex:38}
\begin{colbox}{38}{Seleziona posizioni dispari}{2}
\exlabel{Implementare} \fn{(seleziona-dispari lst)}: restituisce gli elementi
in posizione 1\textsuperscript{a}, 3\textsuperscript{a}, 5\textsuperscript{a},
\ldots\ nella numerazione naturale (indici 0, 2, 4, \ldots\ in notazione 0-based).
\begin{lstlisting}
(seleziona-dispari '(a b c d e))   ; -> (A C E)
\end{lstlisting}


\small\color{black!55}\hyperref[sol:38]{$\rightarrow$ Soluzione}\normalsize\color{black}
\end{colbox}

\label{ex:39}
\begin{essbox}{39}{Reverse con accumulatore}{2}
\exlabel{Implementare} \fn{(reverse* lst)} senza usare \fn{reverse}.

Prima scrivere la versione ricorsiva semplice, poi quella
con accumulatore (tail-recursive).
\begin{lstlisting}
(reverse* '(a b c d))   ; -> (D C B A)
\end{lstlisting}


\small\color{black!55}\hyperref[sol:39]{$\rightarrow$ Soluzione}\normalsize\color{black}
\end{essbox}

\label{ex:40}
\begin{essbox}{40}{Append}{3}
\exlabel{Implementare} \fn{(append* lst1 lst2)} senza usare \fn{append}.
\nopagebreak
\begin{lstlisting}
(append* '(a b) '(c d))   ; -> (A B C D)
(append* '() '(a b))      ; -> (A B)
\end{lstlisting}
Non modificare le liste originali; ricorsione su \fn{lst1}.



\small\color{black!55}\hyperref[sol:40]{$\rightarrow$ Soluzione}\enspace\textcolor{black!30}{|}\enspace\hyperref[hint:40]{$\rightarrow$ Suggerimento}\normalsize\color{black}
\end{essbox}

\label{ex:41}
\begin{stdbox}{41}{Rimozione prima occorrenza}{3}
\exlabel{Implementare} \fn{(remove-first item lst)}: rimuove solo la
prima occorrenza.
\begin{lstlisting}
(remove-first 'b '(a b c b d))   ; -> (A C B D)
\end{lstlisting}
Non usare \fn{remove}; non modificare la lista originale.



\small\color{black!55}\hyperref[sol:41]{$\rightarrow$ Soluzione}\enspace\textcolor{black!30}{|}\enspace\hyperref[hint:41]{$\rightarrow$ Suggerimento}\normalsize\color{black}
\end{stdbox}

\label{ex:42}
\begin{colbox}{42}{Sostituzione prima occorrenza}{2}
\exlabel{Implementare} \fn{(replace-first old new lst)}.
\nopagebreak
\begin{lstlisting}
(replace-first 'b 'x '(a b c b))   ; -> (A X C B)
\end{lstlisting}


\small\color{black!55}\hyperref[sol:42]{$\rightarrow$ Soluzione}\normalsize\color{black}
\end{colbox}

\label{ex:43}
\begin{colbox}{43}{Lista palindroma}{2}
\exlabel{Implementare} \fn{(palindroma-p lst)}.
\nopagebreak
\begin{lstlisting}
(palindroma-p '(a b c b a))   ; -> T
(palindroma-p '(a b c))       ; -> NIL
\end{lstlisting}
Usare la funzione \fn{reverse*} dell'es.~39 (o \fn{reverse}) e \fn{equal}.


\small\color{black!55}\hyperref[sol:43]{$\rightarrow$ Soluzione}\normalsize\color{black}
\end{colbox}

\label{ex:44}
\begin{stdbox}{44}{Range}{2}
\exlabel{Implementare} \fn{(range a b)}: interi da \fn{a} a \fn{b} inclusi.
\nopagebreak
\begin{lstlisting}
(range 1 5)    ; -> (1 2 3 4 5)
(range 3 3)    ; -> (3)
(range 5 3)    ; -> NIL
\end{lstlisting}


\small\color{black!55}\hyperref[sol:44]{$\rightarrow$ Soluzione}\normalsize\color{black}
\end{stdbox}

\label{ex:45}
\begin{colbox}{45}{Duplica elementi}{2}
\exlabel{Implementare} \fn{(duplica-elementi lst)}.
\nopagebreak
\begin{lstlisting}
(duplica-elementi '(a b c))   ; -> (A A B B C C)
\end{lstlisting}


\small\color{black!55}\hyperref[sol:45]{$\rightarrow$ Soluzione}\normalsize\color{black}
\end{colbox}

\label{ex:46}
\begin{colbox}{46}{Estrai chiavi da alist}{2}
\exlabel{Implementare} \fn{(estrai-chiavi alist)}: restituisce la
lista delle chiavi (\fn{car} di ogni coppia puntata).
\begin{lstlisting}
(estrai-chiavi '((a . 1) (b . 2) (c . 3)))
; -> (A B C)
\end{lstlisting}


\small\color{black!55}\hyperref[sol:46]{$\rightarrow$ Soluzione}\normalsize\color{black}
\end{colbox}

\label{ex:47}
\begin{colbox}{47}{Prefisso}{2}
\exlabel{Implementare} \fn{(prefisso-p sub lst)}.
\nopagebreak
\begin{lstlisting}
(prefisso-p '(a b) '(a b c d))   ; -> T
(prefisso-p '(a c) '(a b c d))   ; -> NIL
(prefisso-p '() '(a b))          ; -> T
\end{lstlisting}


\small\color{black!55}\hyperref[sol:47]{$\rightarrow$ Soluzione}\normalsize\color{black}
\end{colbox}

\label{ex:48}
\begin{colbox}{48}{Intercala}{2}
\exlabel{Implementare} \fn{(intercala lst1 lst2)} che alterna gli elementi.
\nopagebreak
\begin{lstlisting}
(intercala '(1 2 3) '(a b c))   ; -> (1 A 2 B 3 C)
\end{lstlisting}
Assumere lunghezze uguali.


\small\color{black!55}\hyperref[sol:48]{$\rightarrow$ Soluzione}\normalsize\color{black}
\end{colbox}

\label{ex:49}
\begin{stdbox}{49}{Run-length encoding}{3}
\textbf{Parte A.} Implementare \fn{(rle-encode lst)}.
\begin{lstlisting}
(rle-encode '(a a a b b c))   ; -> ((3 A) (2 B) (1 C))
\end{lstlisting}
\textbf{Parte B.} Implementare \fn{(rle-decode lst)}: l'inverso.
\begin{lstlisting}
(rle-decode '((3 A) (2 B) (1 C)))   ; -> (A A A B B C)
\end{lstlisting}


\small\color{black!55}\hyperref[sol:49]{$\rightarrow$ Soluzione}\enspace\textcolor{black!30}{|}\enspace\hyperref[hint:49]{$\rightarrow$ Suggerimento}\normalsize\color{black}
\end{stdbox}

\label{ex:50}
\begin{colbox}{50}{Rimozione duplicati}{3}
\exlabel{Implementare} \fn{(remove-duplicates* lst)}: mantiene la prima
occorrenza, senza usare \fn{remove-duplicates}.
\begin{lstlisting}
(remove-duplicates* '(a b a c b a))   ; -> (A B C)
\end{lstlisting}


\small\color{black!55}\hyperref[sol:50]{$\rightarrow$ Soluzione}\enspace\textcolor{black!30}{|}\enspace\hyperref[hint:50]{$\rightarrow$ Suggerimento}\normalsize\color{black}
\end{colbox}

% =============================================================
\setcounter{section}{2}
\section{Ricorsione avanzata}
% =============================================================

\noindent
Tail recursion con accumulatori espliciti, ricorsione su strutture
annidate (alberi impliciti), esponenziazione per quadratura
e primo uso degli array con il crivello di Eratostene.

\bigskip

\label{ex:51}
\begin{essbox}{51}{Fattoriale tail-recursive}{2}
\exlabel{Implementare} \fn{(fattoriale n)} in forma tail-recursive
con accumulatore, usando \fn{labels}.
\begin{lstlisting}
(fattoriale 0)   ; -> 1
(fattoriale 5)   ; -> 120
\end{lstlisting}
\exnota{Obiettivo: capire il pattern accumulatore.}


\small\color{black!55}\hyperref[sol:51]{$\rightarrow$ Soluzione}\normalsize\color{black}
\end{essbox}

\label{ex:52}
\begin{stdbox}{52}{Reverse tail-recursive}{2}
Riscrivere \fn{reverse*} (es.~39) in forma tail-recursive.


\small\color{black!55}\hyperref[sol:52]{$\rightarrow$ Soluzione}\normalsize\color{black}
\end{stdbox}

\label{ex:53}
\begin{colbox}{53}{Somma di un intervallo}{2}
\exlabel{Implementare} \fn{(somma-range a b)} in forma tail-recursive.
\nopagebreak
\begin{lstlisting}
(somma-range 1 100)   ; -> 5050
\end{lstlisting}


\small\color{black!55}\hyperref[sol:53]{$\rightarrow$ Soluzione}\normalsize\color{black}
\end{colbox}

\label{ex:54}
\begin{stdbox}{54}{Fibonacci O(n)}{3}
\exlabel{Implementare} \fn{(fibonacci-fast n)} con due accumulatori,
complessita' $O(n)$.
\begin{lstlisting}
(fibonacci-fast 10)   ; -> 55
(fibonacci-fast 50)   ; -> 12586269025
\end{lstlisting}


\small\color{black!55}\hyperref[sol:54]{$\rightarrow$ Soluzione}\enspace\textcolor{black!30}{|}\enspace\hyperref[hint:54]{$\rightarrow$ Suggerimento}\normalsize\color{black}
\end{stdbox}
\label{ex:55}
\begin{essbox}{55}{Flatten}{3}
\exlabel{Implementare} \fn{(flatten* tree)} che appiattisce strutture annidate.
\nopagebreak
\begin{lstlisting}
(flatten* '(a (b (c d) e) f))   ; -> (A B C D E F)
(flatten* '((1 2) (3 (4 5))))   ; -> (1 2 3 4 5)
\end{lstlisting}


\small\color{black!55}\hyperref[sol:55]{$\rightarrow$ Soluzione}\enspace\textcolor{black!30}{|}\enspace\hyperref[hint:55]{$\rightarrow$ Suggerimento}\normalsize\color{black}
\end{essbox}

\label{ex:56}
\begin{stdbox}{56}{Profondita' di una struttura}{3}
\exlabel{Implementare} \fn{(tree-depth tree)}: profondita' massima di nesting.
\nopagebreak
\begin{lstlisting}
(tree-depth '())             ; -> 0
(tree-depth '(a b c))        ; -> 1
(tree-depth '(a (b c)))      ; -> 2
(tree-depth '(a (b (c d))))  ; -> 3
\end{lstlisting}


\small\color{black!55}\hyperref[sol:56]{$\rightarrow$ Soluzione}\enspace\textcolor{black!30}{|}\enspace\hyperref[hint:56]{$\rightarrow$ Suggerimento}\normalsize\color{black}
\end{stdbox}
\label{ex:57}
\begin{colbox}{57}{Conta atomi}{2}
\exlabel{Implementare} \fn{(count-atoms tree)}.
\nopagebreak
\begin{lstlisting}
(count-atoms '(a (b c) ((d))))   ; -> 4
\end{lstlisting}


\small\color{black!55}\hyperref[sol:57]{$\rightarrow$ Soluzione}\normalsize\color{black}
\end{colbox}

\label{ex:58}
\begin{colbox}{58}{Somma profonda}{2}
\exlabel{Implementare} \fn{(sum-tree tree)}.
\nopagebreak
\begin{lstlisting}
(sum-tree '(1 (2 3) ((4 5))))   ; -> 15
\end{lstlisting}


\small\color{black!55}\hyperref[sol:58]{$\rightarrow$ Soluzione}\normalsize\color{black}
\end{colbox}

\label{ex:59}
\begin{colbox}{59}{Ricerca in struttura annidata}{3}
\exlabel{Implementare} \fn{(tree-member-p item tree)}.
\nopagebreak
\begin{lstlisting}
(tree-member-p 'c '(a (b c) d))   ; -> T
(tree-member-p 'x '(a (b c) d))   ; -> NIL
\end{lstlisting}


\small\color{black!55}\hyperref[sol:59]{$\rightarrow$ Soluzione}\normalsize\color{black}
\end{colbox}

\label{ex:60}
\begin{stdbox}{60}{Sostituzione profonda}{3}
\exlabel{Implementare} \fn{(sostituisci-profondo old new tree)}.
\nopagebreak
\begin{lstlisting}
(sostituisci-profondo 'b 'x '(a (b (b c)) b))
; -> (A (X (X C)) X)
\end{lstlisting}


\small\color{black!55}\hyperref[sol:60]{$\rightarrow$ Soluzione}\enspace\textcolor{black!30}{|}\enspace\hyperref[hint:60]{$\rightarrow$ Suggerimento}\normalsize\color{black}
\end{stdbox}
\label{ex:61}
\begin{stdbox}{61}{Potenza efficiente}{3}
\exlabel{Implementare} \fn{(power-fast base exp)} con complessita' $O(\log n)$:
\[
x^n = \begin{cases}
  (x^{n/2})^2 & \text{se } n \text{ pari} \\
  x \cdot x^{n-1} & \text{se } n \text{ dispari}
\end{cases}
\]
\begin{lstlisting}
(power-fast 2 10)   ; -> 1024
(power-fast 3  0)   ; -> 1
\end{lstlisting}


\small\color{black!55}\hyperref[sol:61]{$\rightarrow$ Soluzione}\enspace\textcolor{black!30}{|}\enspace\hyperref[hint:61]{$\rightarrow$ Suggerimento}\normalsize\color{black}
\end{stdbox}
\label{ex:62}
\begin{stdbox}{62}{MCD e MCM}{2}
Implementare \fn{(mcd* a b)} e \fn{(mcm* a b)} senza usare
\fn{gcd}/\fn{lcm}.
\begin{lstlisting}
(mcd* 48 18)   ; -> 6
(mcm* 4  6)    ; -> 12
\end{lstlisting}


\small\color{black!55}\hyperref[sol:62]{$\rightarrow$ Soluzione}\normalsize\color{black}
\end{stdbox}
\label{ex:63}
\begin{colbox}{63}{Numero primo}{2}
\exlabel{Implementare} \fn{(primo-p n)}.
\nopagebreak
\begin{lstlisting}
(primo-p 17)   ; -> T
(primo-p 18)   ; -> NIL
\end{lstlisting}
Gestire correttamente 0 e 1.


\small\color{black!55}\hyperref[sol:63]{$\rightarrow$ Soluzione}\normalsize\color{black}
\end{colbox}

\label{ex:64}
\begin{essbox}{64}{Crivello di Eratostene}{3}
\exlabel{Implementare} \fn{(crivello n)}: lista di tutti i primi $\leq n$.
\nopagebreak
\begin{lstlisting}
(crivello 20)
; -> (2 3 5 7 11 13 17 19)
\end{lstlisting}
Vincoli: usare array booleano di lunghezza $n+1$; marcare i
multipli con \fn{setf}/\fn{aref}.



\small\color{black!55}\hyperref[sol:64]{$\rightarrow$ Soluzione}\enspace\textcolor{black!30}{|}\enspace\hyperref[hint:64]{$\rightarrow$ Suggerimento}\normalsize\color{black}
\end{essbox}

\label{ex:65}
\begin{stdbox}{65}{Fattori primi}{3}
\exlabel{Implementare} \fn{(prime-factors n)}.
\nopagebreak
\begin{lstlisting}
(prime-factors 60)   ; -> (2 2 3 5)
\end{lstlisting}
Risultato ordinato; $n > 1$.



\small\color{black!55}\hyperref[sol:65]{$\rightarrow$ Soluzione}\enspace\textcolor{black!30}{|}\enspace\hyperref[hint:65]{$\rightarrow$ Suggerimento}\normalsize\color{black}
\end{stdbox}

\label{ex:66}
\begin{colbox}{66}{Sottolista per posizione}{2}
\exlabel{Implementare} \fn{(sublist lst inizio fine)} (indici 0-based,
\fn{fine} escluso).
\begin{lstlisting}
(sublist '(a b c d e) 1 4)   ; -> (B C D)
\end{lstlisting}


\small\color{black!55}\hyperref[sol:66]{$\rightarrow$ Soluzione}\normalsize\color{black}
\end{colbox}

\label{ex:67}
\begin{colbox}{67}{Raggruppa in blocchi}{3}
\exlabel{Implementare} \fn{(raggruppa lst n)}: sottoliste di lunghezza \fn{n}.
\nopagebreak
\begin{lstlisting}
(raggruppa '(a b c d e f g) 3)
; -> ((A B C) (D E F) (G))
\end{lstlisting}


\small\color{black!55}\hyperref[sol:67]{$\rightarrow$ Soluzione}\enspace\textcolor{black!30}{|}\enspace\hyperref[hint:67]{$\rightarrow$ Suggerimento}\normalsize\color{black}
\end{colbox}

\label{ex:68}
\begin{colbox}{68}{Intersezione di insiemi}{2}
\exlabel{Implementare} \fn{(interseca lst1 lst2)} senza usare \fn{intersection}.
\nopagebreak
\begin{lstlisting}
(interseca '(a b c d) '(b d e))   ; -> (B D)
\end{lstlisting}


\small\color{black!55}\hyperref[sol:68]{$\rightarrow$ Soluzione}\normalsize\color{black}
\end{colbox}

% =============================================================
\setcounter{section}{3}
\section{Higher-order e closure}
% =============================================================

\noindent
Funzioni di ordine superiore (\fn{mapcar}, \fn{remove-if},
\fn{reduce}, \fn{apply}, \fn{mapcan}) e closure: funzioni che
catturano variabili dall'ambiente lessicale.

\bigskip

\label{ex:69}
\begin{essbox}{69}{Quadrato degli elementi con mapcar}{1}
\exlabel{Implementare} \fn{(square-list lst)} usando \fn{mapcar} e \fn{lambda}.
\nopagebreak
\begin{lstlisting}
(square-list '(1 2 3 4))   ; -> (1 4 9 16)
\end{lstlisting}


\small\color{black!55}\hyperref[sol:69]{$\rightarrow$ Soluzione}\normalsize\color{black}
\end{essbox}

\label{ex:70}
\begin{stdbox}{70}{Filtro con remove-if-not}{1}
\exlabel{Implementare} \fn{(even-list lst)} usando \fn{remove-if-not}.
\nopagebreak
\begin{lstlisting}
(even-list '(1 2 3 4 5 6))   ; -> (2 4 6)
\end{lstlisting}


\small\color{black!55}\hyperref[sol:70]{$\rightarrow$ Soluzione}\normalsize\color{black}
\end{stdbox}

\label{ex:71}
\begin{colbox}{71}{Stringhe che iniziano con vocale}{2}
\exlabel{Implementare} \fn{(vocali-iniziali lst)} usando \fn{remove-if-not}.
\nopagebreak
\begin{lstlisting}
(vocali-iniziali '("anna" "bob" "elena" "carlo"))
; -> ("anna" "elena")
\end{lstlisting}


\small\color{black!55}\hyperref[sol:71]{$\rightarrow$ Soluzione}\normalsize\color{black}
\end{colbox}

\label{ex:72}
\begin{essbox}{72}{Somma con reduce}{2}
\exlabel{Implementare} \fn{(sum-reduce lst)} con \fn{reduce}.
\nopagebreak
\begin{lstlisting}
(sum-reduce '(1 2 3 4))   ; -> 10
(sum-reduce '())          ; -> 0
\end{lstlisting}


\small\color{black!55}\hyperref[sol:72]{$\rightarrow$ Soluzione}\normalsize\color{black}
\end{essbox}

\label{ex:73}
\begin{colbox}{73}{Prodotto con reduce}{2}
\exlabel{Implementare} \fn{(product-reduce lst)} con \fn{reduce}.
\nopagebreak
\begin{lstlisting}
(product-reduce '(2 3 4))   ; -> 24
\end{lstlisting}


\small\color{black!55}\hyperref[sol:73]{$\rightarrow$ Soluzione}\normalsize\color{black}
\end{colbox}

\label{ex:74}
\begin{colbox}{74}{Min e max con reduce}{2}
\exlabel{Implementare} \fn{(reduce-min lst)} e \fn{(reduce-max lst)}.
\nopagebreak
\begin{lstlisting}
(reduce-min '(8 3 5 1 9))   ; -> 1
(reduce-max '(8 3 5 1 9))   ; -> 9
\end{lstlisting}
Gestire la lista vuota.


\small\color{black!55}\hyperref[sol:74]{$\rightarrow$ Soluzione}\normalsize\color{black}
\end{colbox}

\label{ex:75}
\begin{stdbox}{75}{Lista di cifre in numero intero}{3}
\exlabel{Implementare} \fn{(cifre->intero lst)} con \fn{reduce}.
\nopagebreak
\begin{lstlisting}
(cifre->intero '(1 5 3 4))   ; -> 1534
\end{lstlisting}
Suggerimento: a ogni passo moltiplicare l'accumulatore per 10
e sommare la cifra corrente.



\small\color{black!55}\hyperref[sol:75]{$\rightarrow$ Soluzione}\enspace\textcolor{black!30}{|}\enspace\hyperref[hint:75]{$\rightarrow$ Suggerimento}\normalsize\color{black}
\end{stdbox}

\label{ex:76}
\begin{colbox}{76}{Appiattisce un livello}{3}
\exlabel{Implementare} \fn{(appiattisci-un-livello lst)} con \fn{reduce}
e \fn{append}.
\begin{lstlisting}
(appiattisci-un-livello '((a b) (c d) (e)))
; -> (A B C D E)
\end{lstlisting}


\small\color{black!55}\hyperref[sol:76]{$\rightarrow$ Soluzione}\normalsize\color{black}
\end{colbox}

\label{ex:77}
\begin{colbox}{77}{Regola di Horner}{3}
\exlabel{Implementare} \fn{(horner coefficienti x)} con \fn{reduce}.
\[
p(x) = (\ldots((c_n \cdot x + c_{n-1}) \cdot x + c_{n-2})\ldots + c_0)
\]
\begin{lstlisting}
; 2x^2 + 3x + 1  con x=2  -> 15
(horner '(2 3 1) 2)   ; -> 15
\end{lstlisting}


\small\color{black!55}\hyperref[sol:77]{$\rightarrow$ Soluzione}\enspace\textcolor{black!30}{|}\enspace\hyperref[hint:77]{$\rightarrow$ Suggerimento}\normalsize\color{black}
\end{colbox}

\label{ex:78}
\begin{stdbox}{78}{Prodotto scalare con apply}{2}
\exlabel{Implementare} \fn{(prodotto-scalare v1 v2)}.
\nopagebreak
\begin{lstlisting}
(prodotto-scalare '(1 2 3) '(4 5 6))   ; -> 32
\end{lstlisting}
Vincolo: usare \fn{mapcar} per i prodotti componente per componente,
poi \fn{apply} con \fn{+} per la somma.


\small\color{black!55}\hyperref[sol:78]{$\rightarrow$ Soluzione}\normalsize\color{black}
\end{stdbox}

\label{ex:79}
\begin{stdbox}{79}{mapcan}{2}
\exlabel{Implementare} \fn{(raddoppia lst)} con \fn{mapcan}.
\nopagebreak
\begin{lstlisting}
(raddoppia '(1 2 3))   ; -> (1 1 2 2 3 3)
\end{lstlisting}
Spiegare la differenza tra \fn{mapcan} e \fn{mapcar} in un commento.


\small\color{black!55}\hyperref[sol:79]{$\rightarrow$ Soluzione}\normalsize\color{black}
\end{stdbox}

\label{ex:80}
\begin{colbox}{80}{Ordina per chiave}{2}
\exlabel{Implementare} \fn{(ordina-per-punteggio lst)}: ordina coppie
\fn{(nome . punteggio)} in ordine decrescente.
\begin{lstlisting}
(ordina-per-punteggio
  '((mario . 85) (anna . 92) (luca . 78)))
; -> ((ANNA . 92) (MARIO . 85) (LUCA . 78))
\end{lstlisting}
Vincolo: usare \fn{sort} con comparatore custom.


\small\color{black!55}\hyperref[sol:80]{$\rightarrow$ Soluzione}\normalsize\color{black}
\end{colbox}

\label{ex:81}
\begin{colbox}{81}{Tutti e qualcuno}{2}
\exlabel{Implementare} \fn{(tutti-p pred lst)} e \fn{(qualcuno-p pred lst)}
replicando \fn{every}/\fn{some} senza usarli.
\begin{lstlisting}
(tutti-p #'evenp '(2 4 6))    ; -> T
(qualcuno-p #'oddp '(2 3 4))  ; -> T
\end{lstlisting}


\small\color{black!55}\hyperref[sol:81]{$\rightarrow$ Soluzione}\normalsize\color{black}
\end{colbox}

\label{ex:82}
\begin{colbox}{82}{Conta con predicato}{2}
\exlabel{Implementare} \fn{(count-matching pred lst)} senza usare \fn{count-if}.
\nopagebreak
\begin{lstlisting}
(count-matching #'evenp '(1 2 3 4 5))   ; -> 2
\end{lstlisting}


\small\color{black!55}\hyperref[sol:82]{$\rightarrow$ Soluzione}\normalsize\color{black}
\end{colbox}

\label{ex:83}
\begin{colbox}{83}{Raggruppa per chiave}{3}
\exlabel{Implementare} \fn{(raggruppa-per funzione lst)}: restituisce
una alist con chiave = risultato della funzione.
\begin{lstlisting}
(raggruppa-per #'evenp '(1 2 3 4 5))
; -> ((NIL 1 3 5) (T 2 4))
\end{lstlisting}


\small\color{black!55}\hyperref[sol:83]{$\rightarrow$ Soluzione}\enspace\textcolor{black!30}{|}\enspace\hyperref[hint:83]{$\rightarrow$ Suggerimento}\normalsize\color{black}
\end{colbox}

\label{ex:84}
\begin{colbox}{84}{Applica a tutti}{2}
\exlabel{Implementare} \fn{(applica-a-tutti funzioni x)}.
\nopagebreak
\begin{lstlisting}
(applica-a-tutti (list #'1+ #'1- #'zerop) 0)
; -> (1 -1 T)
\end{lstlisting}


\small\color{black!55}\hyperref[sol:84]{$\rightarrow$ Soluzione}\normalsize\color{black}
\end{colbox}

\label{ex:85}
\begin{essbox}{85}{Composizione di funzioni}{3}
\exlabel{Implementare} \fn{(compose f g)}: restituisce $f(g(x))$ come closure.
\nopagebreak
\begin{lstlisting}
(defparameter *f* (compose #'1+ #'1+))
(funcall *f* 10)   ; -> 12
\end{lstlisting}


\small\color{black!55}\hyperref[sol:85]{$\rightarrow$ Soluzione}\enspace\textcolor{black!30}{|}\enspace\hyperref[hint:85]{$\rightarrow$ Suggerimento}\normalsize\color{black}
\end{essbox}

\label{ex:86}
\begin{stdbox}{86}{Applicazione parziale}{3}
\exlabel{Implementare} \fn{(partial funzione \&rest argomenti)}.
\nopagebreak
\begin{lstlisting}
(defparameter *add10* (partial #'+ 10))
(funcall *add10* 5)    ; -> 15
(funcall *add10* 20)   ; -> 30
\end{lstlisting}


\small\color{black!55}\hyperref[sol:86]{$\rightarrow$ Soluzione}\enspace\textcolor{black!30}{|}\enspace\hyperref[hint:86]{$\rightarrow$ Suggerimento}\normalsize\color{black}
\end{stdbox}

\label{ex:87}
\begin{essbox}{87}{Closure con stato: contatore}{3}
\exlabel{Implementare} \fn{(make-counter \&optional (start 0) (step 1))}.
\nopagebreak
\begin{lstlisting}
(defparameter *c* (make-counter))
(funcall *c*)   ; -> 0
(funcall *c*)   ; -> 1

(defparameter *c2* (make-counter 10 5))
(funcall *c2*)  ; -> 10
(funcall *c2*)  ; -> 15
\end{lstlisting}


\small\color{black!55}\hyperref[sol:87]{$\rightarrow$ Soluzione}\enspace\textcolor{black!30}{|}\enspace\hyperref[hint:87]{$\rightarrow$ Suggerimento}\normalsize\color{black}
\end{essbox}

\label{ex:88}
\begin{colbox}{88}{Accumulatore}{3}
\exlabel{Implementare} \fn{(make-accumulator iniziale)}: ogni chiamata
somma un valore al totale.
\begin{lstlisting}
(defparameter *acc* (make-accumulator 5))
(funcall *acc* 10)   ; -> 15
(funcall *acc* 3)    ; -> 18
\end{lstlisting}


\small\color{black!55}\hyperref[sol:88]{$\rightarrow$ Soluzione}\normalsize\color{black}
\end{colbox}

\label{ex:89}
\begin{colbox}{89}{Generatore Fibonacci con closure}{4}
\exlabel{Implementare} \fn{(make-fib-generator)}: restituisce una funzione
che ad ogni chiamata produce il successivo numero di Fibonacci.
\begin{lstlisting}
(defparameter *fib* (make-fib-generator))
(funcall *fib*)   ; -> 0
(funcall *fib*)   ; -> 1
(funcall *fib*)   ; -> 1
(funcall *fib*)   ; -> 2
\end{lstlisting}


\small\color{black!55}\hyperref[sol:89]{$\rightarrow$ Soluzione}\enspace\textcolor{black!30}{|}\enspace\hyperref[hint:89]{$\rightarrow$ Suggerimento}\normalsize\color{black}
\end{colbox}

\label{ex:90}
\begin{colbox}{90}{Pipeline}{3}
\exlabel{Implementare} \fn{(pipeline valore funzioni)}.
\nopagebreak
\begin{lstlisting}
(pipeline 5 (list #'1+ #'1+ (lambda (x) (* x 2))))
; -> 14
\end{lstlisting}


\small\color{black!55}\hyperref[sol:90]{$\rightarrow$ Soluzione}\normalsize\color{black}
\end{colbox}

\label{ex:91}
\begin{stdbox}{91}{Memoization}{4}
\exlabel{Implementare} \fn{(memoize funzione)} che restituisce una versione
con cache.

Vincoli: hash table nella closure; funzione originale invariata.
\begin{lstlisting}
(defparameter *fib-m*
  (memoize #'fibonacci))
(funcall *fib-m* 35)   ; veloce
\end{lstlisting}


\small\color{black!55}\hyperref[sol:91]{$\rightarrow$ Soluzione}\enspace\textcolor{black!30}{|}\enspace\hyperref[hint:91]{$\rightarrow$ Suggerimento}\normalsize\color{black}
\end{stdbox}

\label{ex:92}
\begin{colbox}{92}{Mini libreria funzionale}{4}
Raccogliere in un unico file, con docstring:
\fn{map-list}, \fn{filter-list}, \fn{count-matching},
\fn{find-matching}, \fn{compose}, \fn{partial}.

Aggiungere almeno 3 funzioni a scelta (es.\ \fn{take}, \fn{drop}, \fn{zip}).

Requisiti: nessuno stato globale; almeno un test per funzione.


\small\color{black!55}\hyperref[sol:92]{$\rightarrow$ Soluzione}\normalsize\color{black}
\end{colbox}

% =============================================================
\setcounter{section}{4}
\section{Strutture dati}
% =============================================================

\noindent
Association list, hash table, array 2D, \fn{defstruct},
alberi binari e BST, pila e coda come ADT, grafi.

\bigskip

\label{ex:93}
\begin{essbox}{93}{Association list}{2}
\exlabel{Implementare} \fn{(lookup key alist)} e \fn{(set-key key value alist)}.
\nopagebreak
\begin{lstlisting}
(lookup 'age '((name . "Mario") (age . 23)))
; -> 23
(set-key 'age 24 '((name . "Mario") (age . 23)))
; -> ((NAME . "Mario") (AGE . 24))
\end{lstlisting}
Usare \fn{assoc} in \fn{lookup}; non modificare la lista originale.


\small\color{black!55}\hyperref[sol:93]{$\rightarrow$ Soluzione}\normalsize\color{black}
\end{essbox}

\label{ex:94}
\begin{stdbox}{94}{Frequenza elementi}{2}
\exlabel{Implementare} \fn{(frequencies lst)}.
\nopagebreak
\begin{lstlisting}
(frequencies '(a b a c b a))
; -> ((A . 3) (B . 2) (C . 1))
\end{lstlisting}


\small\color{black!55}\hyperref[sol:94]{$\rightarrow$ Soluzione}\enspace\textcolor{black!30}{|}\enspace\hyperref[hint:94]{$\rightarrow$ Suggerimento}\normalsize\color{black}
\end{stdbox}

\label{ex:95}
\begin{essbox}{95}{Hash table: operazioni base}{2}
\exlabel{Implementare} \fn{(ht-get key ht)}, \fn{(ht-set key value ht)},
\fn{(ht-delete key ht)}, \fn{(ht-keys ht)}.
\begin{lstlisting}
(defparameter *ht* (make-hash-table))
(ht-set 'foo 42 *ht*)
(ht-get 'foo *ht*)     ; -> 42
(ht-delete 'foo *ht*)
(ht-get 'foo *ht*)     ; -> NIL
\end{lstlisting}


\small\color{black!55}\hyperref[sol:95]{$\rightarrow$ Soluzione}\normalsize\color{black}
\end{essbox}

\label{ex:96}
\begin{colbox}{96}{Frequenza con hash table}{2}
Riscrivere \fn{(frequencies-hash lst)} usando una hash table.

\exnota{Obiettivo: confrontare alist vs hash table.}


\small\color{black!55}\hyperref[sol:96]{$\rightarrow$ Soluzione}\normalsize\color{black}
\end{colbox}

\label{ex:97}
\begin{stdbox}{97}{Frequenza dei caratteri}{3}
\exlabel{Implementare} \fn{(frequenza-caratteri stringa)}.
\nopagebreak
\begin{lstlisting}
(frequenza-caratteri "banana")
; hash table: #\b->1, #\a->3, #\n->2
\end{lstlisting}


\small\color{black!55}\hyperref[sol:97]{$\rightarrow$ Soluzione}\enspace\textcolor{black!30}{|}\enspace\hyperref[hint:97]{$\rightarrow$ Suggerimento}\normalsize\color{black}
\end{stdbox}

\label{ex:98}
\begin{colbox}{98}{Unione di hash table}{3}
\exlabel{Implementare} \fn{(unisci-hash ht1 ht2)}: chiavi duplicate
$\to$ somma dei valori.



\small\color{black!55}\hyperref[sol:98]{$\rightarrow$ Soluzione}\enspace\textcolor{black!30}{|}\enspace\hyperref[hint:98]{$\rightarrow$ Suggerimento}\normalsize\color{black}
\end{colbox}

\label{ex:99}
\begin{colbox}{99}{Inverti hash table}{3}
\exlabel{Implementare} \fn{(inverti-hash ht)}: scambia chiavi e valori.
Assumere valori unici.



\small\color{black!55}\hyperref[sol:99]{$\rightarrow$ Soluzione}\enspace\textcolor{black!30}{|}\enspace\hyperref[hint:99]{$\rightarrow$ Suggerimento}\normalsize\color{black}
\end{colbox}

\label{ex:100}
\begin{stdbox}{100}{Conversione alist/hash}{2}
\exlabel{Implementare} \fn{(alist->hash alist)} e \fn{(hash->alist ht)}.
\nopagebreak
\begin{lstlisting}
(alist->hash '((a . 1) (b . 2)))
; -> hash table a->1, b->2
\end{lstlisting}


\small\color{black!55}\hyperref[sol:100]{$\rightarrow$ Soluzione}\normalsize\color{black}
\end{stdbox}

\label{ex:101}
\begin{essbox}{101}{Matrice con array 2D}{2}
Creare una matrice $3 \times 3$ con \fn{make-array} e implementare:
\begin{lstlisting}
(mat-get m r c)
(mat-set m r c val)
(stampa-matrice m)
\end{lstlisting}
\begin{lstlisting}
(defparameter *m* (make-array '(3 3) :initial-element 0))
(mat-set *m* 1 1 5)
(mat-get *m* 1 1)   ; -> 5
\end{lstlisting}


\small\color{black!55}\hyperref[sol:101]{$\rightarrow$ Soluzione}\normalsize\color{black}
\end{essbox}

\label{ex:102}
\begin{colbox}{102}{Trasposta di matrice}{3}
\exlabel{Implementare} \fn{(trasposta m n)}: trasposta di una matrice $n \times n$.



\small\color{black!55}\hyperref[sol:102]{$\rightarrow$ Soluzione}\enspace\textcolor{black!30}{|}\enspace\hyperref[hint:102]{$\rightarrow$ Suggerimento}\normalsize\color{black}
\end{colbox}

\label{ex:103}
\begin{stdbox}{103}{Inversione array in loco}{2}
\exlabel{Implementare} \fn{(inverti-array! v)} che inverte il vettore
modificandolo con \fn{setf}/\fn{aref}, senza liste ausiliarie.
\begin{lstlisting}
(defparameter *v*
  (make-array 5 :initial-contents '(1 2 3 4 5)))
(inverti-array! *v*)
; *v* ora e': #(5 4 3 2 1)
\end{lstlisting}


\small\color{black!55}\hyperref[sol:103]{$\rightarrow$ Soluzione}\normalsize\color{black}
\end{stdbox}

\label{ex:104}
\begin{colbox}{104}{Moltiplicazione di matrici (opzionale)}{4}
\exlabel{Implementare} \fn{(moltiplica-matrici a b n)}: prodotto di
matrici quadrate $n \times n$.



\small\color{black!55}\hyperref[sol:104]{$\rightarrow$ Soluzione}\enspace\textcolor{black!30}{|}\enspace\hyperref[hint:104]{$\rightarrow$ Suggerimento}\normalsize\color{black}
\end{colbox}

\label{ex:105}
\begin{stdbox}{105}{Struttura con defstruct}{2}
Definire:
\begin{lstlisting}
(defstruct persona
  nome
  (eta 18)
  citta)
\end{lstlisting}
\exlabel{Implementare} \fn{(crea-persona nome eta citta)} che valida l'eta'
(positiva) e segnala errore altrimenti.
\begin{lstlisting}
(crea-persona "Mario" 23 "Roma")
; -> #S(PERSONA :NOME "Mario" :ETA 23 :CITTA "Roma")
\end{lstlisting}


\small\color{black!55}\hyperref[sol:105]{$\rightarrow$ Soluzione}\normalsize\color{black}
\end{stdbox}

\label{ex:106}
\begin{colbox}{106}{Struttura nodo-albero}{2}
Definire:
\begin{lstlisting}
(defstruct nodo
  valore
  (sinistra nil)
  (destra   nil))
\end{lstlisting}
Costruire manualmente l'albero:
\begin{lstlisting}
;       5
;      / \
;     3   8
;    / \
;   1   4
\end{lstlisting}


\small\color{black!55}\hyperref[sol:106]{$\rightarrow$ Soluzione}\normalsize\color{black}
\end{colbox}

\label{ex:107}
\begin{colbox}{107}{Albero con liste}{2}
Rappresentare lo stesso albero con la notazione
\fn{(valore sinistra destra)} e implementare:
\begin{lstlisting}
(make-leaf v)
(make-node v left right)
\end{lstlisting}


\small\color{black!55}\hyperref[sol:107]{$\rightarrow$ Soluzione}\normalsize\color{black}
\end{colbox}

\label{ex:108}
\begin{stdbox}{108}{BST: inserimento e ricerca}{4}
Implementare un albero binario di ricerca:
\begin{lstlisting}
(bst-insert value tree)
(bst-find value tree)
\end{lstlisting}
\begin{lstlisting}
(bst-find 4 (bst-insert 4 (bst-insert 5 nil)))
; -> T
\end{lstlisting}


\small\color{black!55}\hyperref[sol:108]{$\rightarrow$ Soluzione}\enspace\textcolor{black!30}{|}\enspace\hyperref[hint:108]{$\rightarrow$ Suggerimento}\normalsize\color{black}
\end{stdbox}

\label{ex:109}
\begin{stdbox}{109}{Traversal dell'albero}{3}
\exlabel{Implementare} \fn{(preorder tree)}, \fn{(inorder tree)},
\fn{(postorder tree)}.

Verificare che \fn{inorder} su un BST produca sequenza ordinata.
\begin{lstlisting}
(inorder albero)   ; -> (1 3 4 5 8)
\end{lstlisting}


\small\color{black!55}\hyperref[sol:109]{$\rightarrow$ Soluzione}\enspace\textcolor{black!30}{|}\enspace\hyperref[hint:109]{$\rightarrow$ Suggerimento}\normalsize\color{black}
\end{stdbox}

\label{ex:110}
\begin{colbox}{110}{Altezza e conteggio foglie}{3}
\exlabel{Implementare} \fn{(tree-height tree)} e \fn{(count-leaves tree)}.
\nopagebreak
\begin{lstlisting}
(tree-height albero)    ; -> 3
(count-leaves albero)   ; -> 3
\end{lstlisting}


\small\color{black!55}\hyperref[sol:110]{$\rightarrow$ Soluzione}\normalsize\color{black}
\end{colbox}

\label{ex:111}
\begin{colbox}{111}{Specchia l'albero}{3}
\exlabel{Implementare} \fn{(specchia tree)}: scambia ricorsivamente
sottoalberi sinistro e destro.


\small\color{black!55}\hyperref[sol:111]{$\rightarrow$ Soluzione}\normalsize\color{black}
\end{colbox}

\label{ex:112}
\begin{colbox}{112}{Albero bilanciato?}{4}
\exlabel{Implementare} \fn{(bilanciato-p tree)}: le altezze dei sottoalberi
differiscono al piu' di 1 per ogni nodo.



\small\color{black!55}\hyperref[sol:112]{$\rightarrow$ Soluzione}\enspace\textcolor{black!30}{|}\enspace\hyperref[hint:112]{$\rightarrow$ Suggerimento}\normalsize\color{black}
\end{colbox}

\label{ex:113}
\begin{stdbox}{113}{Pila (stack) come ADT}{2}
Implementare con \fn{setf}:
\begin{lstlisting}
(make-stack)
(stack-push item stack)
(stack-pop stack)
(stack-peek stack)
(stack-empty-p stack)
\end{lstlisting}
\begin{lstlisting}
(defparameter *s* (make-stack))
(stack-push 1 *s*)
(stack-push 2 *s*)
(stack-pop *s*)    ; -> 2
(stack-peek *s*)   ; -> 1
\end{lstlisting}


\small\color{black!55}\hyperref[sol:113]{$\rightarrow$ Soluzione}\normalsize\color{black}
\end{stdbox}

\label{ex:114}
\begin{stdbox}{114}{Coda (queue) come ADT}{3}
\exlabel{Implementare:}
\nopagebreak
\begin{lstlisting}
(make-queue)
(enqueue item queue)
(dequeue queue)
(queue-empty-p queue)
\end{lstlisting}


\small\color{black!55}\hyperref[sol:114]{$\rightarrow$ Soluzione}\enspace\textcolor{black!30}{|}\enspace\hyperref[hint:114]{$\rightarrow$ Suggerimento}\normalsize\color{black}
\end{stdbox}

\label{ex:115}
\begin{stdbox}{115}{Valida parentesi}{3}
\exlabel{Implementare} \fn{(valida-parentesi stringa)} per
\fn{()}, \fn{[]}, \fn{\{\}}.
\begin{lstlisting}
(valida-parentesi "([]{})")   ; -> T
(valida-parentesi "([)]")     ; -> NIL
(valida-parentesi "((")       ; -> NIL
\end{lstlisting}
Vincolo: usare una pila.



\small\color{black!55}\hyperref[sol:115]{$\rightarrow$ Soluzione}\enspace\textcolor{black!30}{|}\enspace\hyperref[hint:115]{$\rightarrow$ Suggerimento}\normalsize\color{black}
\end{stdbox}

\label{ex:116}
\begin{colbox}{116}{Rubrica con hash table}{3}
\exlabel{Implementare:}
\nopagebreak
\begin{lstlisting}
(add-contact name phone)
(find-contact name)
(remove-contact name)
(list-contacts)
\end{lstlisting}
Gestire il caso contatto non trovato.


\small\color{black!55}\hyperref[sol:116]{$\rightarrow$ Soluzione}\normalsize\color{black}
\end{colbox}

\label{ex:117}
\begin{stdbox}{117}{Grafo come lista di adiacenza}{3}
\exlabel{Implementare:}
\nopagebreak
\begin{lstlisting}
(neighbors node graph)
(add-edge node1 node2 graph)
\end{lstlisting}
\begin{lstlisting}
(neighbors 'a '((a b c) (b a d) (c a d) (d b c)))
; -> (B C)
\end{lstlisting}


\small\color{black!55}\hyperref[sol:117]{$\rightarrow$ Soluzione}\normalsize\color{black}
\end{stdbox}

\label{ex:118}
\begin{colbox}{118}{DFS su grafo}{4}
\exlabel{Implementare} \fn{(dfs graph start)}: visita in profondita'.
\nopagebreak
\begin{lstlisting}
(dfs '((a b c) (b a d) (c a) (d b)) 'a)
; -> (A B D C)
\end{lstlisting}
Gestire i cicli con un insieme di nodi gia' visitati.



\small\color{black!55}\hyperref[sol:118]{$\rightarrow$ Soluzione}\enspace\textcolor{black!30}{|}\enspace\hyperref[hint:118]{$\rightarrow$ Suggerimento}\normalsize\color{black}
\end{colbox}

\label{ex:119}
\begin{colbox}{119}{BFS su grafo}{4}
\exlabel{Implementare} \fn{(bfs graph start)}: visita in ampiezza.
\nopagebreak
\begin{lstlisting}
(bfs '((a b c) (b a d) (c a) (d b)) 'a)
; -> (A B C D)
\end{lstlisting}
Vincolo: usare una coda.



\small\color{black!55}\hyperref[sol:119]{$\rightarrow$ Soluzione}\enspace\textcolor{black!30}{|}\enspace\hyperref[hint:119]{$\rightarrow$ Suggerimento}\normalsize\color{black}
\end{colbox}

\label{ex:120}
\begin{colbox}{120}{Conta parole}{3}
\exlabel{Implementare} \fn{(conta-parole testo)} e
\fn{(parola-piu-frequente testo)}.
\begin{lstlisting}
(conta-parole "il gatto e' sul tetto e il tetto")
; hash table: "il"->2, ...
\end{lstlisting}


\small\color{black!55}\hyperref[sol:120]{$\rightarrow$ Soluzione}\enspace\textcolor{black!30}{|}\enspace\hyperref[hint:120]{$\rightarrow$ Suggerimento}\normalsize\color{black}
\end{colbox}

% =============================================================
\setcounter{section}{5}
\section{Macro}
% =============================================================

\noindent
Le macro operano a livello sintattico: trasformano codice
prima che venga valutato. Questo capitolo mostra come costruire
nuove strutture di controllo.

\bigskip

\label{ex:121}
\begin{stdbox}{121}{Macro quando}{3}
\exlabel{Implementare:}
\nopagebreak
\begin{lstlisting}
(defmacro quando (test &body corpo))
\end{lstlisting}
Si comporta come \fn{when}: esegue \fn{corpo} solo se \fn{test} e' vero.
\begin{lstlisting}
(quando (> x 0)
  (print "positivo")
  (print "ok"))
\end{lstlisting}
Verificare l'espansione con \fn{macroexpand-1}.

\exnota{Obiettivo: capire la differenza macro/funzione.}



\small\color{black!55}\hyperref[sol:121]{$\rightarrow$ Soluzione}\enspace\textcolor{black!30}{|}\enspace\hyperref[hint:121]{$\rightarrow$ Suggerimento}\normalsize\color{black}
\end{stdbox}

\label{ex:122}
\begin{stdbox}{122}{Macro while}{3}
\exlabel{Implementare:}
\nopagebreak
\begin{lstlisting}
(defmacro mio-while (test &body corpo))
\end{lstlisting}
\begin{lstlisting}
(let ((i 0))
  (mio-while (< i 5)
    (print i)
    (setf i (1+ i))))
; stampa 0 1 2 3 4
\end{lstlisting}


\small\color{black!55}\hyperref[sol:122]{$\rightarrow$ Soluzione}\enspace\textcolor{black!30}{|}\enspace\hyperref[hint:122]{$\rightarrow$ Suggerimento}\normalsize\color{black}
\end{stdbox}

\label{ex:123}
\begin{colbox}{123}{Macro for}{4}
\exlabel{Implementare:}
\nopagebreak
\begin{lstlisting}
(defmacro mio-for ((var da a) &body corpo))
\end{lstlisting}
\begin{lstlisting}
(mio-for (i 1 5)
  (format t "~A~%" i))
; stampa 1 2 3 4 5
\end{lstlisting}
\exnota{Obiettivo: usare \fn{gensym} per evitare cattura della variabile.}



\small\color{black!55}\hyperref[sol:123]{$\rightarrow$ Soluzione}\enspace\textcolor{black!30}{|}\enspace\hyperref[hint:123]{$\rightarrow$ Suggerimento}\normalsize\color{black}
\end{colbox}

% =============================================================
\setcounter{section}{6}
\section{Backtracking e algoritmi}
% =============================================================

\noindent
Algoritmi che esplorano uno spazio di ricerca: generazione di
combinazioni, ordinamento e problemi classici di backtracking.

\bigskip

\label{ex:124}
\begin{stdbox}{124}{Tutti i sottoinsiemi}{3}
\exlabel{Implementare} \fn{(subsets lst)}: genera tutti i $2^n$ sottoinsiemi.
\nopagebreak
\begin{lstlisting}
(subsets '(a b))
; -> (() (A) (B) (A B))
\end{lstlisting}


\small\color{black!55}\hyperref[sol:124]{$\rightarrow$ Soluzione}\enspace\textcolor{black!30}{|}\enspace\hyperref[hint:124]{$\rightarrow$ Suggerimento}\normalsize\color{black}
\end{stdbox}

\label{ex:125}
\begin{stdbox}{125}{Permutazioni}{4}
\exlabel{Implementare} \fn{(permutations lst)}.
\nopagebreak
\begin{lstlisting}
(length (permutations '(a b c)))   ; -> 6
\end{lstlisting}


\small\color{black!55}\hyperref[sol:125]{$\rightarrow$ Soluzione}\enspace\textcolor{black!30}{|}\enspace\hyperref[hint:125]{$\rightarrow$ Suggerimento}\normalsize\color{black}
\end{stdbox}

\label{ex:126}
\begin{stdbox}{126}{Combinazioni}{3}
\exlabel{Implementare} \fn{(combinations lst k)}.
\nopagebreak
\begin{lstlisting}
(combinations '(a b c d) 2)
; -> ((A B) (A C) (A D) (B C) (B D) (C D))
\end{lstlisting}


\small\color{black!55}\hyperref[sol:126]{$\rightarrow$ Soluzione}\enspace\textcolor{black!30}{|}\enspace\hyperref[hint:126]{$\rightarrow$ Suggerimento}\normalsize\color{black}
\end{stdbox}

\label{ex:127}
\begin{stdbox}{127}{Somma a target}{3}
\exlabel{Implementare} \fn{(somma-a-target numeri target)}.
\nopagebreak
\begin{lstlisting}
(somma-a-target '(1 2 3 4 5) 6)
; -> ((1 5) (2 4) (1 2 3))
\end{lstlisting}


\small\color{black!55}\hyperref[sol:127]{$\rightarrow$ Soluzione}\enspace\textcolor{black!30}{|}\enspace\hyperref[hint:127]{$\rightarrow$ Suggerimento}\normalsize\color{black}
\end{stdbox}

\label{ex:128}
\begin{colbox}{128}{Parentesi bilanciate generate}{3}
\exlabel{Implementare} \fn{(genera-parentesi n)}.
\nopagebreak
\begin{lstlisting}
(genera-parentesi 3)
; -> ("((()))" "(()())" "(())()" "()(())" "()()()")
\end{lstlisting}


\small\color{black!55}\hyperref[sol:128]{$\rightarrow$ Soluzione}\enspace\textcolor{black!30}{|}\enspace\hyperref[hint:128]{$\rightarrow$ Suggerimento}\normalsize\color{black}
\end{colbox}

\label{ex:129}
\begin{stdbox}{129}{Merge sort}{4}
\exlabel{Implementare} \fn{(merge-sorted lst1 lst2)} e \fn{(merge-sort* lst)}.
\nopagebreak
\begin{lstlisting}
(merge-sort* '(7 2 9 1 5 3))   ; -> (1 2 3 5 7 9)
\end{lstlisting}
Vincoli: non usare \fn{sort}; complessita' $O(n \log n)$.



\small\color{black!55}\hyperref[sol:129]{$\rightarrow$ Soluzione}\enspace\textcolor{black!30}{|}\enspace\hyperref[hint:129]{$\rightarrow$ Suggerimento}\normalsize\color{black}
\end{stdbox}

\label{ex:130}
\begin{stdbox}{130}{Quicksort}{4}
\exlabel{Implementare} \fn{(quicksort* lst)} senza usare \fn{sort}.
\nopagebreak
\begin{lstlisting}
(quicksort* '(7 2 9 1 5 3))   ; -> (1 2 3 5 7 9)
\end{lstlisting}


\small\color{black!55}\hyperref[sol:130]{$\rightarrow$ Soluzione}\enspace\textcolor{black!30}{|}\enspace\hyperref[hint:130]{$\rightarrow$ Suggerimento}\normalsize\color{black}
\end{stdbox}

\label{ex:131}
\begin{stdbox}{131}{Torri di Hanoi}{3}
\exlabel{Implementare} \fn{(hanoi n da a tramite)} che stampa le mosse.
\nopagebreak
\begin{lstlisting}
(hanoi 3 'A 'C 'B)
; Sposta disco 1: A -> C
; ...
\end{lstlisting}


\small\color{black!55}\hyperref[sol:131]{$\rightarrow$ Soluzione}\enspace\textcolor{black!30}{|}\enspace\hyperref[hint:131]{$\rightarrow$ Suggerimento}\normalsize\color{black}
\end{stdbox}

\label{ex:132}
\begin{colbox}{132}{Problema dello zaino}{4}
\exlabel{Implementare} \fn{(knapsack oggetti capacita)}.
Oggetti: lista di \fn{(peso valore)}.
\begin{lstlisting}
(knapsack '((2 3) (3 4) (4 5) (5 8)) 5)
; -> 7
\end{lstlisting}


\small\color{black!55}\hyperref[sol:132]{$\rightarrow$ Soluzione}\enspace\textcolor{black!30}{|}\enspace\hyperref[hint:132]{$\rightarrow$ Suggerimento}\normalsize\color{black}
\end{colbox}

\label{ex:133}
\begin{colbox}{133}{Labirinto}{4}
Rappresentare un labirinto come matrice 2D (0 libero, 1 muro).
\exlabel{Implementare} \fn{(find-path maze start goal)}.
\nopagebreak
\begin{lstlisting}
(find-path maze '(0 0) '(4 4))
; -> ((0 0) (1 0) ... (4 4))
\end{lstlisting}


\small\color{black!55}\hyperref[sol:133]{$\rightarrow$ Soluzione}\enspace\textcolor{black!30}{|}\enspace\hyperref[hint:133]{$\rightarrow$ Suggerimento}\normalsize\color{black}
\end{colbox}

\label{ex:134}
\begin{stdbox}{134}{N-regine}{5}
\exlabel{Implementare} \fn{(n-queens n)}: tutte le configurazioni valide.
\nopagebreak
\begin{lstlisting}
(length (n-queens 4))   ; -> 2
(length (n-queens 8))   ; -> 92
\end{lstlisting}
Vincoli: backtracking con pruning; rappresentazione compatta.



\small\color{black!55}\hyperref[sol:134]{$\rightarrow$ Soluzione}\enspace\textcolor{black!30}{|}\enspace\hyperref[hint:134]{$\rightarrow$ Suggerimento}\normalsize\color{black}
\end{stdbox}

\label{ex:135}
\begin{colbox}{135}{Sudoku}{5}
\exlabel{Implementare} \fn{(solve-sudoku board)}.
\fn{board}: array $9 \times 9$ con 0 per celle vuote.

Algoritmo: per ogni cella vuota, provare i valori 1--9,
verificare vincoli riga/colonna/blocco, procedere ricorsivamente.



\small\color{black!55}\hyperref[sol:135]{$\rightarrow$ Soluzione}\enspace\textcolor{black!30}{|}\enspace\hyperref[hint:135]{$\rightarrow$ Suggerimento}\normalsize\color{black}
\end{colbox}

\label{ex:136}
\begin{colbox}{136}{Cammino del cavallo}{5}
\exlabel{Implementare} \fn{(knights-tour n)}: cammino del cavallo su
scacchiera $n \times n$ visitando ogni casella esattamente una volta.
\begin{lstlisting}
(knights-tour 5)   ; -> lista di mosse, o NIL
\end{lstlisting}
Suggerimento: euristica di Warnsdorff.



\small\color{black!55}\hyperref[sol:136]{$\rightarrow$ Soluzione}\enspace\textcolor{black!30}{|}\enspace\hyperref[hint:136]{$\rightarrow$ Suggerimento}\normalsize\color{black}
\end{colbox}

% =============================================================
\setcounter{section}{7}
\section{Stringhe e I/O}
% =============================================================

\noindent
Operazioni sulle stringhe e tecniche di I/O.

\bigskip

\label{ex:137}
\begin{stdbox}{137}{Inverti una stringa}{2}
\exlabel{Implementare} \fn{(reverse-string s)} senza applicare \fn{reverse}
direttamente; usare \fn{coerce}.
\begin{lstlisting}
(reverse-string "hello")   ; -> "olleh"
\end{lstlisting}


\small\color{black!55}\hyperref[sol:137]{$\rightarrow$ Soluzione}\normalsize\color{black}
\end{stdbox}

\label{ex:138}
\begin{stdbox}{138}{Stringa palindroma}{2}
\exlabel{Implementare} \fn{(palindrome-p s)}.
\nopagebreak
\begin{lstlisting}
(palindrome-p "anna")    ; -> T
(palindrome-p "casa")    ; -> NIL
\end{lstlisting}


\small\color{black!55}\hyperref[sol:138]{$\rightarrow$ Soluzione}\normalsize\color{black}
\end{stdbox}

\label{ex:139}
\begin{colbox}{139}{Palindromo robusto}{3}
\exlabel{Implementare} \fn{(palindrome-robust-p s)}: ignora
maiuscole/minuscole, spazi, punteggiatura.
\begin{lstlisting}
(palindrome-robust-p "I topi non avevano nipoti")
; -> T
\end{lstlisting}


\small\color{black!55}\hyperref[sol:139]{$\rightarrow$ Soluzione}\enspace\textcolor{black!30}{|}\enspace\hyperref[hint:139]{$\rightarrow$ Suggerimento}\normalsize\color{black}
\end{colbox}

\label{ex:140}
\begin{stdbox}{140}{Anagrammi}{3}
\exlabel{Implementare} \fn{(anagram-p s1 s2)}, ignorando maiuscole.
\nopagebreak
\begin{lstlisting}
(anagram-p "roma" "amor")   ; -> T
\end{lstlisting}


\small\color{black!55}\hyperref[sol:140]{$\rightarrow$ Soluzione}\enspace\textcolor{black!30}{|}\enspace\hyperref[hint:140]{$\rightarrow$ Suggerimento}\normalsize\color{black}
\end{stdbox}

\label{ex:141}
\begin{colbox}{141}{Conta caratteri}{2}
\exlabel{Implementare} \fn{(count-char c s)}.
\nopagebreak
\begin{lstlisting}
(count-char #\a "banana")   ; -> 3
\end{lstlisting}


\small\color{black!55}\hyperref[sol:141]{$\rightarrow$ Soluzione}\normalsize\color{black}
\end{colbox}

% =============================================================
\setcounter{section}{8}
\section{Parsing e interprete}
% =============================================================

\noindent
Un valutatore di espressioni simboliche senza usare \fn{eval}:
il codice come dato, pattern fondamentale di Lisp.

\bigskip

\label{ex:142}
\begin{stdbox}{142}{Calcolatrice RPN}{4}
\exlabel{Implementare} \fn{(rpn-eval espressione)}: valuta notazione
polacca inversa.
\begin{lstlisting}
(rpn-eval '(3 4 +))          ; -> 7
(rpn-eval '(3 4 * 2 +))      ; -> 14
(rpn-eval '(5 1 2 + 4 * + 3 -)) ; -> 14
\end{lstlisting}
Vincolo: usare una pila.



\small\color{black!55}\hyperref[sol:142]{$\rightarrow$ Soluzione}\enspace\textcolor{black!30}{|}\enspace\hyperref[hint:142]{$\rightarrow$ Suggerimento}\normalsize\color{black}
\end{stdbox}

\label{ex:143}
\begin{stdbox}{143}{Valutatore di espressioni}{4}
\exlabel{Implementare} \fn{(evaluate-expression expr)}: espressioni
in notazione prefissa, senza usare \fn{eval}.
\begin{lstlisting}
(evaluate-expression '(+ 2 3))        ; -> 5
(evaluate-expression '(* 4 (+ 2 3)))  ; -> 20
\end{lstlisting}


\small\color{black!55}\hyperref[sol:143]{$\rightarrow$ Soluzione}\enspace\textcolor{black!30}{|}\enspace\hyperref[hint:143]{$\rightarrow$ Suggerimento}\normalsize\color{black}
\end{stdbox}

\label{ex:144}
\begin{colbox}{144}{Mini-interprete con variabili}{5}
Estendere l'es.~143 con variabili (\fn{set x 10}) e riferimenti.
\begin{lstlisting}
(mio-eval '(set x 10) env)
(mio-eval '(+ x 5) env)   ; -> 15
\end{lstlisting}
Ambiente separato come alist; non usare \fn{eval}.



\small\color{black!55}\hyperref[sol:144]{$\rightarrow$ Soluzione}\enspace\textcolor{black!30}{|}\enspace\hyperref[hint:144]{$\rightarrow$ Suggerimento}\normalsize\color{black}
\end{colbox}

% =============================================================
\setcounter{section}{9}
\section{Progetti}
% =============================================================

\noindent
Progetti completi che integrano strutture dati, I/O, gestione
degli errori e organizzazione del codice.

\bigskip

\label{ex:145}
\begin{stdbox}{145}{Indovina il numero}{3}
Il programma sceglie un numero casuale tra 1 e 100,
l'utente lo indovina con feedback troppo alto/troppo basso.
\begin{lstlisting}
(gioca)
; Indovina il numero (1-100):
; > 50
; Troppo alto!
; ...
; Esatto! Ci hai messo 5 tentativi.
\end{lstlisting}
Vincolo: \fn{loop} per il ciclo; \fn{random} per la generazione.


\small\color{black!55}\hyperref[sol:145]{$\rightarrow$ Soluzione}\normalsize\color{black}
\end{stdbox}

\label{ex:146}
\begin{colbox}{146}{Tris (tic-tac-toe)}{4}
Gioco del tris a due giocatori da tastiera.
\begin{itemize}[leftmargin=2em]
  \item Griglia $3 \times 3$ con array.
  \item Verifica vittoria: righe, colonne, diagonali.
  \item Turno alternato; rilevamento pareggio.
\end{itemize}
Separare logica del gioco, verifica vittoria e I/O.



\small\color{black!55}\hyperref[sol:146]{$\rightarrow$ Soluzione}\enspace\textcolor{black!30}{|}\enspace\hyperref[hint:146]{$\rightarrow$ Suggerimento}\normalsize\color{black}
\end{colbox}

\label{ex:147}
\begin{colbox}{147}{Gioco della vita (Conway)}{4}
\exlabel{Implementare} \fn{(evolvi griglia)}: un passo di evoluzione.
Griglia: array 2D di 0 (morto) e 1 (vivo).
Regole: sopravvive con 2--3 vicini vivi; nasce con esattamente 3.
\begin{lstlisting}
(loop repeat 5
      do (setf g (evolvi g))
         (stampa-griglia g))
\end{lstlisting}


\small\color{black!55}\hyperref[sol:147]{$\rightarrow$ Soluzione}\enspace\textcolor{black!30}{|}\enspace\hyperref[hint:147]{$\rightarrow$ Suggerimento}\normalsize\color{black}
\end{colbox}

\label{ex:148}
\begin{stdbox}{148}{Mini sistema gestionale}{5}
Programma completo (studenti, libri, prodotti, film o contatti).

Requisiti minimi:
\begin{itemize}[leftmargin=2em]
  \item \fn{defstruct} o CLOS;
  \item hash table o lista come archivio;
  \item inserimento, ricerca, modifica, eliminazione,
        filtraggio, ordinamento;
  \item salvataggio e caricamento da file;
  \item gestione degli errori;
  \item docstring su ogni funzione pubblica.
\end{itemize}

Struttura minima:
\begin{lstlisting}
;; MODEL  -- strutture dati
;; LOGIC  -- operazioni
;; IO     -- input/output e file
\end{lstlisting}



\small\color{black!55}\hyperref[sol:148]{$\rightarrow$ Soluzione}\enspace\textcolor{black!30}{|}\enspace\hyperref[hint:148]{$\rightarrow$ Suggerimento}\normalsize\color{black}
\end{stdbox}








% =============================================================
%  03_suggerimenti.tex — Suggerimenti per esercizi ★★★ e superiori
% =============================================================
\chapter*{Suggerimenti}
\addcontentsline{toc}{chapter}{Suggerimenti}
\markboth{Suggerimenti}{Suggerimenti}
\setcounter{hintsec}{0}  % reset contatore sezioni hint

\noindent\small\color{black!65}
Suggerimenti per gli esercizi \treStelle\ o superiore.
Prima di leggere, fai almeno un tentativo autonomo: anche una soluzione parziale è più utile della risposta immediata.
\color{black}

\vspace{8pt}

% =============================================================
\hintsection{Fondamenti}
% =============================================================

\label{hint:17}
\begin{hintbox}{17}{Somma delle cifre}
\small\color{black!55}\hyperref[ex:17]{$\leftarrow$ Esercizio}\normalsize\color{black}

Estrai le cifre una alla volta con \fn{mod} e \fn{floor}:
\fn{(mod n 10)} da' l'ultima cifra; \fn{(floor n 10)} rimuove
l'ultima cifra. Continua finche' \fn{n} e' zero.

\small\color{black!55}\hyperref[sol:17]{$\rightarrow$ Soluzione}\normalsize\color{black}
\end{hintbox}

% =============================================================
\hintsection{Liste}
% =============================================================

\label{hint:40}
\begin{hintbox}{40}{Append}
\small\color{black!55}\hyperref[ex:40]{$\leftarrow$ Esercizio}\normalsize\color{black}

Casi: se \fn{lst1} e' vuota, restituisci \fn{lst2}.
Altrimenti costruisci con \fn{cons}: il nuovo primo elemento
e' \fn{(car lst1)}, il resto e' la ricorsione su \fn{(cdr lst1)}
e \fn{lst2}.
Non modificare \fn{lst2}: non servono copie, \fn{cons} crea
una nuova struttura automaticamente.

\small\color{black!55}\hyperref[sol:40]{$\rightarrow$ Soluzione}\normalsize\color{black}
\end{hintbox}

\label{hint:41}
\begin{hintbox}{41}{Remove-first}
\small\color{black!55}\hyperref[ex:41]{$\leftarrow$ Esercizio}\normalsize\color{black}

Caso base 1: lista vuota $\to$ \fn{NIL}.
Caso base 2: \fn{(car lst)} uguale a \fn{item} $\to$ restituisci
\fn{(cdr lst)} (salta solo questo elemento).
Caso ricorsivo: \fn{cons} di \fn{(car lst)} con la ricorsione
sul resto. Non rimuovere altre occorrenze: basta non continuare
a cercare dopo aver trovato la prima.

\small\color{black!55}\hyperref[sol:41]{$\rightarrow$ Soluzione}\normalsize\color{black}
\end{hintbox}

\label{hint:49}
\begin{hintbox}{49}{Run-length encoding}
\small\color{black!55}\hyperref[ex:49]{$\leftarrow$ Esercizio}\normalsize\color{black}

\hintschema{Encode:} tieni traccia dell'elemento corrente e del contatore.
Quando l'elemento cambia (o la lista e' finita), aggiungi
la coppia \fn{(contatore elemento)} al risultato.
Una struttura con accumulatori e' piu' naturale di una
puramente ricorsiva.

\hintschema{Decode:} per ogni coppia \fn{(n elem)}, genera \fn{n}
copie di \fn{elem}. Usa un loop interno o una funzione ausiliaria
che costruisce \fn{n} ripetizioni.

\small\color{black!55}\hyperref[sol:49]{$\rightarrow$ Soluzione}\normalsize\color{black}
\end{hintbox}

\label{hint:50}
\begin{hintbox}{50}{Remove-duplicates*}
\small\color{black!55}\hyperref[ex:50]{$\leftarrow$ Esercizio}\normalsize\color{black}

Visita la lista da sinistra. Per ogni elemento, controllane
la presenza nel risultato accumulato (o nella parte gia'
processata). Se non e' presente, includilo; altrimenti saltalo.
Usa \fn{member} (o la tua \fn{member*}) per il controllo.

\small\color{black!55}\hyperref[sol:50]{$\rightarrow$ Soluzione}\normalsize\color{black}
\end{hintbox}

% =============================================================
\hintsection{Ricorsione avanzata}
% =============================================================

\label{hint:52}
\begin{hintbox}{52}{Reverse tail-recursive}
\small\color{black!55}\hyperref[ex:52]{$\leftarrow$ Esercizio}\normalsize\color{black}

Definire con \fn{labels} una funzione ausiliaria con due parametri:
\fn{lst} (la lista da rovesciare) e \fn{acc} (l'accumulatore, inizialmente \fn{NIL}).
Ad ogni passo: se \fn{lst} è vuota restituire \fn{acc};
altrimenti richiamare la funzione con \fn{(cdr lst)} e
\fn{(cons (car lst) acc)}.
La chiamata ricorsiva è l'ultima operazione: nessun calcolo al ritorno.

\small\color{black!55}\hyperref[sol:52]{$\rightarrow$ Soluzione}\normalsize\color{black}
\end{hintbox}

\label{hint:54}
\begin{hintbox}{54}{Fibonacci O(n)}
\small\color{black!55}\hyperref[ex:54]{$\leftarrow$ Esercizio}\normalsize\color{black}

Usa due accumulatori \fn{a} e \fn{b} che rappresentano
$F(i)$ e $F(i+1)$. A ogni passo: il nuovo \fn{a} diventa
il vecchio \fn{b}; il nuovo \fn{b} diventa \fn{a+b}.
Alla fine, \fn{a} e' $F(n)$.

Caso base: se \fn{n = 0}, restituisci \fn{a} (inizializzato a 0).
Inizia con \fn{(iter n 0 1)}.

\small\color{black!55}\hyperref[sol:54]{$\rightarrow$ Soluzione}\normalsize\color{black}
\end{hintbox}

\label{hint:55}
\begin{hintbox}{55}{Flatten*}
\small\color{black!55}\hyperref[ex:55]{$\leftarrow$ Esercizio}\normalsize\color{black}

Tratta ogni elemento della lista come il nodo di un albero binario
implicito: se e' un \textit{atomo}, e' una foglia (mettila nel risultato);
se e' una \textit{lista}, e' un sottoalbero (appiattiscila ricorsivamente).

Schema:
\begin{lstlisting}
(defun flatten* (tree)
  (cond
    ((null tree) '())
    ((atom tree) (list tree))
    (t (append (flatten* (car tree))
               (flatten* (cdr tree))))))
\end{lstlisting}
\hintschema{Attenzione:} ricorrere su \fn{car} e \fn{cdr} separatamente,
non su \fn{(car tree)} e \fn{(cdr tree)} come lista.

\small\color{black!55}\hyperref[sol:55]{$\rightarrow$ Soluzione}\normalsize\color{black}
\end{hintbox}

\label{hint:56}
\begin{hintbox}{56}{Tree-depth}
\small\color{black!55}\hyperref[ex:56]{$\leftarrow$ Esercizio}\normalsize\color{black}

La profondita' di una lista vuota e' 0.
La profondita' di una lista non vuota e' il massimo tra:
\begin{itemize}[leftmargin=2em]
  \item la profondita' del \fn{car} $+1$ (se il \fn{car} e' una lista);
  \item la profondita' del \fn{cdr}.
\end{itemize}
Usa \fn{max} per confrontare i due rami.

\small\color{black!55}\hyperref[sol:56]{$\rightarrow$ Soluzione}\normalsize\color{black}
\end{hintbox}

\label{hint:60}
\begin{hintbox}{60}{Sostituzione profonda}
\small\color{black!55}\hyperref[ex:60]{$\leftarrow$ Esercizio}\normalsize\color{black}

Struttura simile a \fn{flatten*}: ricorri su \fn{car} e \fn{cdr}.
Se il \fn{car} e' uguale a \fn{old}, sostituiscilo con \fn{new};
se e' una lista, ricorri dentro; altrimenti lascialo invariato.

\small\color{black!55}\hyperref[sol:60]{$\rightarrow$ Soluzione}\normalsize\color{black}
\end{hintbox}

\label{hint:61}
\begin{hintbox}{61}{Power-fast}
\small\color{black!55}\hyperref[ex:61]{$\leftarrow$ Esercizio}\normalsize\color{black}

Tre casi: \fn{exp = 0} $\to$ 1; \fn{exp} pari $\to$ calcola
\fn{(power-fast base (/ exp 2))} e metti al quadrato;
\fn{exp} dispari $\to$ moltiplica per \fn{base} e richiama
con \fn{(1- exp)}.

Con \fn{let} puoi evitare la doppia chiamata ricorsiva
nel caso pari:
\begin{lstlisting}
(let ((half (power-fast base (/ exp 2))))
  (* half half))
\end{lstlisting}

\small\color{black!55}\hyperref[sol:61]{$\rightarrow$ Soluzione}\normalsize\color{black}
\end{hintbox}

\label{hint:64}
\begin{hintbox}{64}{Crivello di Eratostene}
\small\color{black!55}\hyperref[ex:64]{$\leftarrow$ Esercizio}\normalsize\color{black}

\begin{enumerate}[leftmargin=2em]
  \item Crea un array booleano \fn{sieve} di dimensione $n+1$,
        inizializzato a \fn{T} (tutti candidati).
  \item Marca 0 e 1 come non-primi.
  \item Per ogni \fn{i} da 2 a $\sqrt{n}$: se \fn{(aref sieve i)}
        e' \fn{T}, marca tutti i multipli di \fn{i} (da $i^2$)
        come \fn{NIL}.
  \item Raccogli tutti gli indici per cui il valore e' ancora \fn{T}.
\end{enumerate}

\small\color{black!55}\hyperref[sol:64]{$\rightarrow$ Soluzione}\normalsize\color{black}
\end{hintbox}

\label{hint:65}
\begin{hintbox}{65}{Fattori primi}
\small\color{black!55}\hyperref[ex:65]{$\leftarrow$ Esercizio}\normalsize\color{black}

Prova a dividere \fn{n} per ogni intero \fn{d} partendo da 2.
Se \fn{(mod n d) = 0}, \fn{d} e' un fattore: aggiungilo al risultato
e continua con \fn{n/d} (non avanzare \fn{d}!).
Se \fn{d} non divide, avanza \fn{d} di 1.
Termina quando \fn{d * d > n}: se \fn{n > 1} e' rimasto
un fattore primo.

\small\color{black!55}\hyperref[sol:65]{$\rightarrow$ Soluzione}\normalsize\color{black}
\end{hintbox}

\label{hint:67}
\begin{hintbox}{67}{Raggruppa in blocchi}
\small\color{black!55}\hyperref[ex:67]{$\leftarrow$ Esercizio}\normalsize\color{black}

Scrivi una funzione ausiliaria \fn{prendi-n lst n} che restituisce
\fn{(primi-n resto)}: i primi \fn{n} elementi e il resto della lista.
Poi \fn{raggruppa} chiama \fn{prendi-n}, prende il primo gruppo,
e ricorre sul resto.

\small\color{black!55}\hyperref[sol:67]{$\rightarrow$ Soluzione}\normalsize\color{black}
\end{hintbox}

% =============================================================
\hintsection{Higher-order e closure}
% =============================================================

\label{hint:75}
\begin{hintbox}{75}{Cifre in intero}
\small\color{black!55}\hyperref[ex:75]{$\leftarrow$ Esercizio}\normalsize\color{black}

Usa \fn{reduce} con una funzione che, dato l'accumulatore \fn{acc}
e la cifra corrente \fn{d}, calcola \fn{(+ (* acc 10) d)}.
Con \fn{:initial-value 0} funziona anche per lista vuota.

\small\color{black!55}\hyperref[sol:75]{$\rightarrow$ Soluzione}\normalsize\color{black}
\end{hintbox}

\label{hint:77}
\begin{hintbox}{77}{Regola di Horner}
\small\color{black!55}\hyperref[ex:77]{$\leftarrow$ Esercizio}\normalsize\color{black}

La regola di Horner valuta $p(x) = c_n x^n + \ldots + c_0$
come $(\ldots((c_n \cdot x + c_{n-1}) \cdot x + c_{n-2})\ldots + c_0)$.

Con \fn{reduce}: la funzione riceve l'accumulatore \fn{acc}
(che parte da $c_n$) e il coefficiente corrente \fn{c},
e calcola \fn{(+ (* acc x) c)}.
I coefficienti devono essere nell'ordine dal grado piu' alto
al piu' basso.

\small\color{black!55}\hyperref[sol:77]{$\rightarrow$ Soluzione}\normalsize\color{black}
\end{hintbox}

\label{hint:83}
\begin{hintbox}{83}{Raggruppa per chiave}
\small\color{black!55}\hyperref[ex:83]{$\leftarrow$ Esercizio}\normalsize\color{black}

Per ogni elemento, calcola la chiave. Cerca nella alist risultato
se quella chiave esiste gia'. Se si', aggiungici l'elemento;
se no, crea una nuova coppia \fn{(chiave elemento)}.
Usa \fn{assoc} per la ricerca e aggiorna la alist con \fn{setf}/\fn{aref}
o ricostruendola.

\small\color{black!55}\hyperref[sol:83]{$\rightarrow$ Soluzione}\normalsize\color{black}
\end{hintbox}

\label{hint:85}
\begin{hintbox}{85}{Compose}
\small\color{black!55}\hyperref[ex:85]{$\leftarrow$ Esercizio}\normalsize\color{black}

Una closure cattura \fn{f} e \fn{g} nell'ambiente lessicale.
Restituisce una \fn{lambda} che prende \fn{x}, applica prima
\fn{g}, poi \fn{f}:
\begin{lstlisting}
(lambda (x) (funcall f (funcall g x)))
\end{lstlisting}
Ricorda che le funzioni si passano con \fn{\#'}.

\small\color{black!55}\hyperref[sol:85]{$\rightarrow$ Soluzione}\normalsize\color{black}
\end{hintbox}

\label{hint:86}
\begin{hintbox}{86}{Partial application}
\small\color{black!55}\hyperref[ex:86]{$\leftarrow$ Esercizio}\normalsize\color{black}

Cattura \fn{funzione} e \fn{argomenti} nella closure.
La lambda restituita prende \fn{\&rest rest-args} e chiama
\fn{funzione} con la lista \fn{(append argomenti rest-args)}.
Usa \fn{apply} per passare la lista combinata:
\begin{lstlisting}
(lambda (&rest rest-args)
  (apply funzione (append argomenti rest-args)))
\end{lstlisting}

\small\color{black!55}\hyperref[sol:86]{$\rightarrow$ Soluzione}\normalsize\color{black}
\end{hintbox}

\label{hint:87}
\begin{hintbox}{87}{Closure con stato}
\small\color{black!55}\hyperref[ex:87]{$\leftarrow$ Esercizio}\normalsize\color{black}

Lo stato vive in un \fn{let} attorno alla \fn{lambda}.
Ogni chiamata alla lambda legge e aggiorna la variabile.
\begin{lstlisting}
(let ((curr start))
  (lambda ()
    (prog1 curr
      (incf curr step))))
\end{lstlisting}
\fn{prog1} restituisce il valore prima dell'incremento.

\small\color{black!55}\hyperref[sol:87]{$\rightarrow$ Soluzione}\normalsize\color{black}
\end{hintbox}

\label{hint:89}
\begin{hintbox}{89}{Generatore Fibonacci con closure}
\small\color{black!55}\hyperref[ex:89]{$\leftarrow$ Esercizio}\normalsize\color{black}

Due variabili nello stesso \fn{let} condividono la stessa closure:
\begin{lstlisting}
(let ((a 0) (b 1))
  (lambda ()
    (prog1 a
      (psetf a b b (+ a b)))))
\end{lstlisting}
\fn{psetf} aggiorna le variabili simultaneamente (senza
che il nuovo valore di \fn{a} influenzi il calcolo di \fn{b}).

\small\color{black!55}\hyperref[sol:89]{$\rightarrow$ Soluzione}\normalsize\color{black}
\end{hintbox}

\label{hint:91}
\begin{hintbox}{91}{Memoize}
\small\color{black!55}\hyperref[ex:91]{$\leftarrow$ Esercizio}\normalsize\color{black}

La struttura base:
\begin{lstlisting}
(let ((cache (make-hash-table :test #'equal)))
  (lambda (&rest args)
    (multiple-value-bind (val found)
        (gethash args cache)
      (if found
          val
          (setf (gethash args cache)
                (apply funzione args))))))
\end{lstlisting}
La chiave della cache e' la lista degli argomenti: usa
\fn{:test \#'equal} per confrontare liste.
\fn{multiple-value-bind} permette di distinguere ``valore NIL
trovato'' da ``chiave non presente''.

\small\color{black!55}\hyperref[sol:91]{$\rightarrow$ Soluzione}\normalsize\color{black}
\end{hintbox}

% =============================================================
\hintsection{Strutture dati}
% =============================================================

\label{hint:94}
\begin{hintbox}{94}{Frequenze con alist}
\small\color{black!55}\hyperref[ex:94]{$\leftarrow$ Esercizio}\normalsize\color{black}

Per ogni elemento: cerca nella alist con \fn{assoc}.
Se la coppia esiste, incrementa il \fn{cdr}; altrimenti,
aggiungi una nuova coppia \fn{(elemento . 1)}.
Con \fn{setf} puoi modificare il \fn{cdr} di una coppia
trovata con \fn{assoc}: \fn{(setf (cdr coppia) (1+ (cdr coppia)))}.

\small\color{black!55}\hyperref[sol:94]{$\rightarrow$ Soluzione}\normalsize\color{black}
\end{hintbox}

\label{hint:97}
\begin{hintbox}{97}{Frequenza dei caratteri}
\small\color{black!55}\hyperref[ex:97]{$\leftarrow$ Esercizio}\normalsize\color{black}

Crea una hash table con \fn{:test \#'eql} (i caratteri si
confrontano con \fn{eql}).
Itera sui caratteri con \fn{(coerce stringa 'list)}
o \fn{(loop for c across stringa ...)}.
Per ogni carattere: incrementa il valore nella tabella
(\fn{(incf (gethash c ht 0))}).

\small\color{black!55}\hyperref[sol:97]{$\rightarrow$ Soluzione}\normalsize\color{black}
\end{hintbox}

\label{hint:98}
\begin{hintbox}{98}{Unione di hash table}
\small\color{black!55}\hyperref[ex:98]{$\leftarrow$ Esercizio}\normalsize\color{black}

Crea una nuova hash table. Itera su \fn{ht1} con \fn{maphash}
e copia tutte le coppie. Poi itera su \fn{ht2}: se la chiave
esiste gia', somma i valori; altrimenti inserisci.

\small\color{black!55}\hyperref[sol:98]{$\rightarrow$ Soluzione}\normalsize\color{black}
\end{hintbox}

\label{hint:99}
\begin{hintbox}{99}{Inverti hash table}
\small\color{black!55}\hyperref[ex:99]{$\leftarrow$ Esercizio}\normalsize\color{black}

Crea una nuova hash table. Itera con \fn{maphash} sulla
tabella originale: per ogni coppia \fn{(k v)}, inserisci
\fn{(v k)} nella nuova tabella.

\small\color{black!55}\hyperref[sol:99]{$\rightarrow$ Soluzione}\normalsize\color{black}
\end{hintbox}

\label{hint:102}
\begin{hintbox}{102}{Trasposta di matrice}
\small\color{black!55}\hyperref[ex:102]{$\leftarrow$ Esercizio}\normalsize\color{black}

Crea una nuova matrice $n \times n$. Con due loop annidati
\fn{(loop for i ...)} e \fn{(loop for j ...)}:
\fn{(setf (aref risultato i j) (aref originale j i))}.

\small\color{black!55}\hyperref[sol:102]{$\rightarrow$ Soluzione}\normalsize\color{black}
\end{hintbox}

\label{hint:104}
\begin{hintbox}{104}{Moltiplicazione di matrici (opzionale)}
\small\color{black!55}\hyperref[ex:104]{$\leftarrow$ Esercizio}\normalsize\color{black}

Il prodotto $C = A \times B$ con $C_{ij} = \sum_k A_{ik} \cdot B_{kj}$.
Tre loop annidati: indici $i$, $j$, $k$.
\begin{lstlisting}
(loop for i below n do
  (loop for j below n do
    (setf (aref c i j)
      (loop for k below n
            sum (* (aref a i k)
                   (aref b k j))))))
\end{lstlisting}

\small\color{black!55}\hyperref[sol:104]{$\rightarrow$ Soluzione}\normalsize\color{black}
\end{hintbox}

\label{hint:108}
\begin{hintbox}{108}{BST: inserimento e ricerca}
\small\color{black!55}\hyperref[ex:108]{$\leftarrow$ Esercizio}\normalsize\color{black}

\textbf{Insert:} se l'albero e' \fn{nil}, crea un nodo foglia.
Se il valore e' minore del nodo corrente, ricorri a sinistra;
altrimenti a destra. Costruisci il nuovo albero con \fn{make-nodo}.

\textbf{Find:} se l'albero e' \fn{nil}, non trovato.
Se il valore e' uguale al nodo, trovato.
Se minore, vai a sinistra; se maggiore, vai a destra.

\small\color{black!55}\hyperref[sol:108]{$\rightarrow$ Soluzione}\normalsize\color{black}
\end{hintbox}

\label{hint:109}
\begin{hintbox}{109}{Traversal dell'albero}
\small\color{black!55}\hyperref[ex:109]{$\leftarrow$ Esercizio}\normalsize\color{black}

I tre traversal differiscono solo per l'ordine in cui si
visita la radice rispetto ai sottoalberi:
\begin{itemize}[leftmargin=2em]
  \item \textbf{preorder:} radice, sinistra, destra.
  \item \textbf{inorder:} sinistra, radice, destra.
  \item \textbf{postorder:} sinistra, destra, radice.
\end{itemize}
Usa \fn{append} per concatenare i risultati dei due sottoalberi.

\small\color{black!55}\hyperref[sol:109]{$\rightarrow$ Soluzione}\normalsize\color{black}
\end{hintbox}

\label{hint:112}
\begin{hintbox}{112}{Albero bilanciato?}
\small\color{black!55}\hyperref[ex:112]{$\leftarrow$ Esercizio}\normalsize\color{black}

Scrivi una funzione ausiliaria che restituisce l'altezza
se l'albero e' bilanciato, oppure $-1$ come segnale di
sbilanciamento.
\begin{lstlisting}
; ritorna altezza se bilanciato, -1 altrimenti
(defun check (tree)
  (if (null tree)
      0
      (let ((hl (check (node-left tree)))
            (hr (check (node-right tree))))
        (if (or (= hl -1) (= hr -1)
                (> (abs (- hl hr)) 1))
            -1
            (1+ (max hl hr))))))
\end{lstlisting}

\small\color{black!55}\hyperref[sol:112]{$\rightarrow$ Soluzione}\normalsize\color{black}
\end{hintbox}

\label{hint:114}
\begin{hintbox}{114}{Coda FIFO}
\small\color{black!55}\hyperref[ex:114]{$\leftarrow$ Esercizio}\normalsize\color{black}

Rappresenta la coda come una struttura con due liste:
\fn{testa} (elementi da estrarre) e \fn{coda} (elementi aggiunti).
\begin{itemize}[leftmargin=2em]
  \item \textbf{Enqueue:} aggiungi in \fn{coda} con \fn{cons}.
  \item \textbf{Dequeue:} prendi da \fn{testa}. Se \fn{testa}
        e' vuota, ribalta \fn{coda} in \fn{testa} con \fn{reverse}.
\end{itemize}
Questa rappresentazione da' dequeue ammortizzato $O(1)$.

\small\color{black!55}\hyperref[sol:114]{$\rightarrow$ Soluzione}\normalsize\color{black}
\end{hintbox}

\label{hint:115}
\begin{hintbox}{115}{Valida parentesi}
\small\color{black!55}\hyperref[ex:115]{$\leftarrow$ Esercizio}\normalsize\color{black}

Itera sui caratteri della stringa. Per ogni carattere:
\begin{itemize}[leftmargin=2em]
  \item parentesi aperta \fn{(}, \fn{[}, \fn{\{}: push sulla pila.
  \item parentesi chiusa \fn{)}, \fn{]}, \fn{\}}: se la pila
        e' vuota o il top non corrisponde, restituisci \fn{NIL}.
        Altrimenti pop.
\end{itemize}
Alla fine, la stringa e' valida se e solo se la pila e' vuota.
Usa una alist per mappare ogni chiusa alla corrispondente aperta.

\small\color{black!55}\hyperref[sol:115]{$\rightarrow$ Soluzione}\normalsize\color{black}
\end{hintbox}

\label{hint:118}
\begin{hintbox}{118}{DFS su grafo}
\small\color{black!55}\hyperref[ex:118]{$\leftarrow$ Esercizio}\normalsize\color{black}

Usa una lista \fn{visitati} per tenere traccia dei nodi gia' visti.
Schema ricorsivo:
\begin{enumerate}[leftmargin=2em]
  \item Se il nodo e' gia' visitato, restituisci \fn{NIL}.
  \item Aggiungi il nodo a \fn{visitati}.
  \item Per ogni vicino non ancora visitato, ricorri.
  \item Accumula l'ordine di visita.
\end{enumerate}
Passa \fn{visitati} come parametro o usala come variabile
catturata da una closure/\fn{labels}.

\small\color{black!55}\hyperref[sol:118]{$\rightarrow$ Soluzione}\normalsize\color{black}
\end{hintbox}

\label{hint:119}
\begin{hintbox}{119}{BFS su grafo}
\small\color{black!55}\hyperref[ex:119]{$\leftarrow$ Esercizio}\normalsize\color{black}

Usa una coda (o semplicemente una lista come coda FIFO).
Schema:
\begin{enumerate}[leftmargin=2em]
  \item Inserisci il nodo di partenza nella coda.
  \item Finche' la coda non e' vuota:
        estrai il primo nodo; se non ancora visitato,
        aggiungilo ai visitati e inserisci i suoi vicini in coda.
\end{enumerate}
Con una lista, usa \fn{append} per aggiungere in coda e
\fn{car}/\fn{cdr} per estrarre dalla testa.

\small\color{black!55}\hyperref[sol:119]{$\rightarrow$ Soluzione}\normalsize\color{black}
\end{hintbox}

\label{hint:120}
\begin{hintbox}{120}{Conta parole}
\small\color{black!55}\hyperref[ex:120]{$\leftarrow$ Esercizio}\normalsize\color{black}

Per spezzare la stringa in parole, implementa un tokenizer
semplice: scansiona carattere per carattere, accumulando
caratteri non-spazio in una parola corrente.
Quando trovi uno spazio (o la fine), se la parola corrente
non e' vuota, aggiungila alla lista.
Poi usa una hash table con \fn{:test \#'equal} (le stringhe
si confrontano con \fn{equal}) per contare le occorrenze.

\small\color{black!55}\hyperref[sol:120]{$\rightarrow$ Soluzione}\normalsize\color{black}
\end{hintbox}

% =============================================================
\hintsection{Macro}
% =============================================================

\label{hint:121}
\begin{hintbox}{121}{Macro quando}
\small\color{black!55}\hyperref[ex:121]{$\leftarrow$ Esercizio}\normalsize\color{black}

Lo schema base di una macro con body:
\begin{lstlisting}
(defmacro quando (test &body corpo)
  `(if ,test
       (progn ,@corpo)
       nil))
\end{lstlisting}
Il backquote costruisce la lista quasi-letteralmente;
\fn{,test} inserisce il valore di \fn{test};
\fn{,@corpo} inserisce gli elementi di \fn{corpo}
al livello corrente (splice).

Verifica con: \fn{(macroexpand-1 '(quando (> x 0) (print x)))}.

\small\color{black!55}\hyperref[sol:121]{$\rightarrow$ Soluzione}\normalsize\color{black}
\end{hintbox}

\label{hint:122}
\begin{hintbox}{122}{Macro while}
\small\color{black!55}\hyperref[ex:122]{$\leftarrow$ Esercizio}\normalsize\color{black}

Un'implementazione semplice usa \fn{loop} nell'espansione:
\begin{lstlisting}
(defmacro mio-while (test &body corpo)
  `(loop while ,test do ,@corpo))
\end{lstlisting}
Oppure, per capire meglio il meccanismo di basso livello,
usa \fn{tagbody}/\fn{go}:
\begin{lstlisting}
(defmacro mio-while (test &body corpo)
  (let ((inizio (gensym)) (fine (gensym)))
    `(tagbody
       ,inizio
       (unless ,test (go ,fine))
       ,@corpo
       (go ,inizio)
       ,fine)))
\end{lstlisting}
\fn{gensym} genera simboli unici per evitare conflitti di nomi.

\small\color{black!55}\hyperref[sol:122]{$\rightarrow$ Soluzione}\normalsize\color{black}
\end{hintbox}

\label{hint:123}
\begin{hintbox}{123}{Macro for}
\small\color{black!55}\hyperref[ex:123]{$\leftarrow$ Esercizio}\normalsize\color{black}

Usa \fn{gensym} per la variabile di ciclo e per evitare che
valutazioni multiple di \fn{da}/\fn{a} causino effetti collaterali.
\begin{lstlisting}
(defmacro mio-for ((var da a) &body corpo)
  (let ((start (gensym)) (end (gensym)))
    `(let ((,start ,da) (,end ,a))
       (loop for ,var from ,start to ,end
             do ,@corpo))))
\end{lstlisting}
Nota: qui \fn{var} non e' avvolto in \fn{gensym} perche'
e' il nome che l'utente vuole usare nel corpo.

\small\color{black!55}\hyperref[sol:123]{$\rightarrow$ Soluzione}\normalsize\color{black}
\end{hintbox}

% =============================================================
\hintsection{Backtracking e algoritmi}
% =============================================================

\label{hint:124}
\begin{hintbox}{124}{Subsets}
\small\color{black!55}\hyperref[ex:124]{$\leftarrow$ Esercizio}\normalsize\color{black}

Ogni elemento ha due scelte: incluso o non incluso.
Schema ricorsivo:
\begin{itemize}[leftmargin=2em]
  \item Caso base: lista vuota $\to$ \fn{'(())}.
  \item Caso ricorsivo: calcola i sottoinsiemi del \fn{cdr}.
        I sottoinsiemi totali sono quelli senza \fn{(car lst)},
        piu' quelli con \fn{(car lst)} aggiunto in testa.
\end{itemize}

\small\color{black!55}\hyperref[sol:124]{$\rightarrow$ Soluzione}\normalsize\color{black}
\end{hintbox}

\label{hint:125}
\begin{hintbox}{125}{Permutazioni}
\small\color{black!55}\hyperref[ex:125]{$\leftarrow$ Esercizio}\normalsize\color{black}

Per generare le permutazioni di \fn{lst}:
\begin{enumerate}[leftmargin=2em]
  \item Per ogni elemento \fn{x} nella lista:
  \item Rimuovi \fn{x} dalla lista (ottenendo \fn{rest}).
  \item Genera le permutazioni di \fn{rest}.
  \item Aggiungi \fn{x} in testa a ogni permutazione.
\end{enumerate}
Usa una funzione ausiliaria che rimuove un elemento per posizione
(non per valore, per gestire duplicati correttamente).

\small\color{black!55}\hyperref[sol:125]{$\rightarrow$ Soluzione}\normalsize\color{black}
\end{hintbox}

\label{hint:126}
\begin{hintbox}{126}{Combinazioni}
\small\color{black!55}\hyperref[ex:126]{$\leftarrow$ Esercizio}\normalsize\color{black}

Per scegliere \fn{k} elementi da \fn{lst}:
\begin{itemize}[leftmargin=2em]
  \item Se \fn{k = 0}: restituisci \fn{'(())}.
  \item Se la lista e' vuota: \fn{NIL}.
  \item Altrimenti: unisci le combinazioni che \textit{includono}
        \fn{(car lst)} (combinazioni di \fn{k-1} dal resto)
        con quelle che \textit{non lo includono} (combinazioni
        di \fn{k} dal resto).
\end{itemize}

\small\color{black!55}\hyperref[sol:126]{$\rightarrow$ Soluzione}\normalsize\color{black}
\end{hintbox}

\label{hint:127}
\begin{hintbox}{127}{Somma a target}
\small\color{black!55}\hyperref[ex:127]{$\leftarrow$ Esercizio}\normalsize\color{black}

Adatta il pattern dei sottoinsiemi: a ogni passo, scegli
di includere o meno il numero corrente.
Se il \fn{target} diventa 0, hai trovato una soluzione.
Se il \fn{target} diventa negativo o la lista e' vuota
senza aver raggiunto 0, il ramo fallisce.

\small\color{black!55}\hyperref[sol:127]{$\rightarrow$ Soluzione}\normalsize\color{black}
\end{hintbox}

\label{hint:128}
\begin{hintbox}{128}{Parentesi bilanciate generate}
\small\color{black!55}\hyperref[ex:128]{$\leftarrow$ Esercizio}\normalsize\color{black}

Costruisci la stringa con due contatori: \fn{aperte} e \fn{chiuse}.
A ogni passo puoi aggiungere \fn{(} se \fn{aperte < n};
puoi aggiungere \fn{)} se \fn{chiuse < aperte}.
Quando \fn{aperte = chiuse = n}, hai una stringa valida.

\small\color{black!55}\hyperref[sol:128]{$\rightarrow$ Soluzione}\normalsize\color{black}
\end{hintbox}

\label{hint:129}
\begin{hintbox}{129}{Merge sort}
\small\color{black!55}\hyperref[ex:129]{$\leftarrow$ Esercizio}\normalsize\color{black}

\textbf{Merge-sorted:} due puntatori alle teste delle liste.
Prendi il minore dei due \fn{car}, avanza il puntatore corrispondente.
Caso base: quando una lista e' vuota, aggiungi il resto dell'altra.

\textbf{Merge-sort:} dividi la lista a meta' (usa \fn{length} e
\fn{nthcdr} per trovare il punto di divisione),
ordina le due meta' ricorsivamente, poi unisci con \fn{merge-sorted}.

\small\color{black!55}\hyperref[sol:129]{$\rightarrow$ Soluzione}\normalsize\color{black}
\end{hintbox}

\label{hint:130}
\begin{hintbox}{130}{Quicksort}
\small\color{black!55}\hyperref[ex:130]{$\leftarrow$ Esercizio}\normalsize\color{black}

Scegli il primo elemento come pivot.
Con \fn{remove-if} (o la tua \fn{filter-list}), dividi
il resto in elementi minori e maggiori-o-uguali al pivot.
Ordina ricorsivamente le due parti e concatena:
\fn{minori ++ [pivot] ++ maggiori}.

\small\color{black!55}\hyperref[sol:130]{$\rightarrow$ Soluzione}\normalsize\color{black}
\end{hintbox}

\label{hint:131}
\begin{hintbox}{131}{Torri di Hanoi}
\small\color{black!55}\hyperref[ex:131]{$\leftarrow$ Esercizio}\normalsize\color{black}

La soluzione e' elegantemente ricorsiva:
\begin{enumerate}[leftmargin=2em]
  \item Sposta $n-1$ dischi da \fn{da} a \fn{tramite} (usando \fn{a}).
  \item Sposta il disco $n$ da \fn{da} a \fn{a}.
  \item Sposta $n-1$ dischi da \fn{tramite} a \fn{a} (usando \fn{da}).
\end{enumerate}
Caso base: \fn{n = 0}, non fare nulla.
Il numero di mosse e' $2^n - 1$.

\small\color{black!55}\hyperref[sol:131]{$\rightarrow$ Soluzione}\normalsize\color{black}
\end{hintbox}

\label{hint:132}
\begin{hintbox}{132}{Problema dello zaino}
\small\color{black!55}\hyperref[ex:132]{$\leftarrow$ Esercizio}\normalsize\color{black}

Schema di ricerca esaustiva:
Per ogni oggetto, hai due scelte: prenderlo o lasciarlo.
\begin{itemize}[leftmargin=2em]
  \item Prenderlo: sottrai il peso dalla capacita', aggiungi
        il valore, ricorri sul resto degli oggetti.
  \item Non prenderlo: ricorri sul resto con capacita' invariata.
\end{itemize}
Restituisci il massimo tra i due rami.
Attenzione: prendi l'oggetto solo se il suo peso non supera
la capacita' residua.

\small\color{black!55}\hyperref[sol:132]{$\rightarrow$ Soluzione}\normalsize\color{black}
\end{hintbox}

\label{hint:133}
\begin{hintbox}{133}{Labirinto}
\small\color{black!55}\hyperref[ex:133]{$\leftarrow$ Esercizio}\normalsize\color{black}

Schema DFS con backtracking:
\begin{enumerate}[leftmargin=2em]
  \item Se la cella corrente e' il goal, restituisci il percorso.
  \item Marca la cella come visitata.
  \item Per ogni direzione (su, giu', sinistra, destra):
        se la cella adiacente e' valida, libera e non visitata,
        ricorri. Se restituisce un percorso, propagalo.
  \item Se nessuna direzione funziona, demarks e restituisci \fn{NIL}.
\end{enumerate}
Usa un array separato per i visitati, oppure marca nel labirinto stesso
(ricordando di ripristinare al backtrack).

\small\color{black!55}\hyperref[sol:133]{$\rightarrow$ Soluzione}\normalsize\color{black}
\end{hintbox}

\label{hint:134}
\begin{hintbox}{134}{N-regine}
\small\color{black!55}\hyperref[ex:134]{$\leftarrow$ Esercizio}\normalsize\color{black}

\hintschema{Rappresentazione:} usa un vettore di \fn{n} elementi dove
\fn{(aref board col)} e' la riga in cui si trova la regina
nella colonna \fn{col}. Stai posizionando una regina per colonna.

\hintschema{Pruning:} prima di posizionare la regina nella colonna \fn{col}
alla riga \fn{row}, verifica che nessuna delle regine gia'
posizionate attacchi questa cella:
\begin{itemize}[leftmargin=2em]
  \item stessa riga: \fn{(= (aref board prev-col) row)};
  \item stessa diagonale: \fn{(= (abs (- row (aref board prev-col))) (- col prev-col))}.
\end{itemize}

\textbf{Struttura:} per ogni colonna, prova tutte le righe 0..(n-1).
Se valida, assegna e ricorri sulla colonna successiva.
Se \fn{col = n}, hai trovato una soluzione.

\small\color{black!55}\hyperref[sol:134]{$\rightarrow$ Soluzione}\normalsize\color{black}
\end{hintbox}

\label{hint:135}
\begin{hintbox}{135}{Sudoku}
\small\color{black!55}\hyperref[ex:135]{$\leftarrow$ Esercizio}\normalsize\color{black}

\begin{enumerate}[leftmargin=2em]
  \item Trova la prima cella vuota (valore 0).
  \item Se non ce ne sono, il puzzle e' risolto.
  \item Per ogni valore da 1 a 9:
        verifica che non compaia nella stessa riga, colonna,
        e blocco $3 \times 3$.
        Se valido, assegna il valore e ricorri.
        Se la ricorsione fallisce, rimetti 0 (backtrack).
\end{enumerate}

Per ottimizzare, scegli la cella con meno valori possibili
(``minimum remaining values heuristic'').

\small\color{black!55}\hyperref[sol:135]{$\rightarrow$ Soluzione}\normalsize\color{black}
\end{hintbox}

\label{hint:136}
\begin{hintbox}{136}{Cammino del cavallo}
\small\color{black!55}\hyperref[ex:136]{$\leftarrow$ Esercizio}\normalsize\color{black}

Le 8 mosse possibili del cavallo in termini di $(dx, dy)$:
$(\pm 1, \pm 2)$ e $(\pm 2, \pm 1)$.

\textbf{Backtracking puro:} prova ogni mossa valida (dentro la
scacchiera e non gia' visitata); se porta a una soluzione,
propagala; altrimenti annulla e prova la successiva.

\textbf{Euristica di Warnsdorff:} ordina le mosse in base al numero
di uscite disponibili dalla cella successiva (meno uscite = prima).
Riduce drasticamente il backtracking.

\small\color{black!55}\hyperref[sol:136]{$\rightarrow$ Soluzione}\normalsize\color{black}
\end{hintbox}

% =============================================================
\hintsection{Stringhe e I/O}
% =============================================================

\label{hint:139}
\begin{hintbox}{139}{Palindromo robusto}
\small\color{black!55}\hyperref[ex:139]{$\leftarrow$ Esercizio}\normalsize\color{black}

\begin{enumerate}[leftmargin=2em]
  \item Converti in minuscolo con \fn{string-downcase}.
  \item Filtra i soli caratteri alfabetici con \fn{remove-if-not}
        e \fn{alpha-char-p}.
  \item Confronta la stringa filtrata con la sua inversa.
\end{enumerate}

\small\color{black!55}\hyperref[sol:139]{$\rightarrow$ Soluzione}\normalsize\color{black}
\end{hintbox}

\label{hint:140}
\begin{hintbox}{140}{Anagrammi}
\small\color{black!55}\hyperref[ex:140]{$\leftarrow$ Esercizio}\normalsize\color{black}

Due stringhe sono anagrammi se, ordinate lettera per lettera,
sono uguali. Converti entrambe in minuscolo, poi in lista
di caratteri con \fn{coerce}, ordina con \fn{sort}, e confronta.

\small\color{black!55}\hyperref[sol:140]{$\rightarrow$ Soluzione}\normalsize\color{black}
\end{hintbox}

% =============================================================
\hintsection{Parsing e interprete}
% =============================================================

\label{hint:142}
\begin{hintbox}{142}{Calcolatrice RPN}
\small\color{black!55}\hyperref[ex:142]{$\leftarrow$ Esercizio}\normalsize\color{black}

Itera sugli elementi dell'espressione:
\begin{itemize}[leftmargin=2em]
  \item Se e' un numero: push sulla pila.
  \item Se e' un operatore: pop due valori, applica
        l'operatore, push il risultato.
\end{itemize}
Alla fine, la pila contiene un solo elemento: il risultato.
Usa \fn{numberp} per distinguere numeri da operatori;
usa una alist o \fn{case} per mappare simbolo $\to$ funzione.

\small\color{black!55}\hyperref[sol:142]{$\rightarrow$ Soluzione}\normalsize\color{black}
\end{hintbox}

\label{hint:143}
\begin{hintbox}{143}{Valutatore di espressioni}
\small\color{black!55}\hyperref[ex:143]{$\leftarrow$ Esercizio}\normalsize\color{black}

La struttura dell'espressione rispecchia la struttura del codice:
\begin{itemize}[leftmargin=2em]
  \item Se e' un numero: restituiscilo direttamente.
  \item Se e' una lista: il \fn{car} e' l'operatore,
        il \fn{cadr} e il \fn{caddr} sono i due operandi.
        Valutali ricorsivamente, poi applica l'operatore.
\end{itemize}
Usa \fn{funcall} con una funzione ricavata dal simbolo
(es. \fn{(symbol-function operator)}) oppure un \fn{case}.

\small\color{black!55}\hyperref[sol:143]{$\rightarrow$ Soluzione}\normalsize\color{black}
\end{hintbox}

\label{hint:144}
\begin{hintbox}{144}{Mini-interprete con variabili}
\small\color{black!55}\hyperref[ex:144]{$\leftarrow$ Esercizio}\normalsize\color{black}

L'ambiente e' una alist \fn{((nome . valore) ...)}.

Estendi il valutatore dell'es.~143:
\begin{itemize}[leftmargin=2em]
  \item Se l'espressione e' un simbolo: cercalo nell'ambiente
        con \fn{assoc}.
  \item Se e' \fn{(set var val)}: valuta \fn{val}, aggiorna
        l'ambiente (aggiungi o modifica la coppia).
  \item Operatori aritmetici: come prima.
\end{itemize}
Gestisci l'ambiente come parametro (passato per riferimento
o restituito insieme al valore).

\small\color{black!55}\hyperref[sol:144]{$\rightarrow$ Soluzione}\normalsize\color{black}
\end{hintbox}

% =============================================================
\hintsection{Progetti}
% =============================================================

\label{hint:146}
\begin{hintbox}{146}{Tris}
\small\color{black!55}\hyperref[ex:146]{$\leftarrow$ Esercizio}\normalsize\color{black}

Struttura consigliata:
\begin{itemize}[leftmargin=2em]
  \item \fn{crea-griglia}: array $3 \times 3$ di \fn{NIL}.
  \item \fn{fai-mossa griglia riga col giocatore}: verifica
        che la cella sia libera e assegna.
  \item \fn{vinto-p griglia giocatore}: controlla le 8
        combinazioni vincenti (3 righe, 3 colonne, 2 diagonali).
  \item \fn{pareggio-p griglia}: nessuna cella vuota e nessuno ha vinto.
  \item \fn{gioca}: loop che alterna i turni fino a vittoria o pareggio.
\end{itemize}

\small\color{black!55}\hyperref[sol:146]{$\rightarrow$ Soluzione}\normalsize\color{black}
\end{hintbox}

\label{hint:147}
\begin{hintbox}{147}{Gioco della vita}
\small\color{black!55}\hyperref[ex:147]{$\leftarrow$ Esercizio}\normalsize\color{black}

Per ogni cella $(r, c)$, conta i vicini vivi (le 8 celle adiacenti).
Le regole:
\begin{itemize}[leftmargin=2em]
  \item Viva con 2 o 3 vicini $\to$ sopravvive.
  \item Morta con esattamente 3 vicini $\to$ nasce.
  \item Altrimenti $\to$ muore/resta morta.
\end{itemize}
Calcola la nuova generazione in un array separato (non
modificare quello corrente mentre lo leggi).

\small\color{black!55}\hyperref[sol:147]{$\rightarrow$ Soluzione}\normalsize\color{black}
\end{hintbox}

\label{hint:148}
\begin{hintbox}{148}{Mini sistema gestionale}
\small\color{black!55}\hyperref[ex:148]{$\leftarrow$ Esercizio}\normalsize\color{black}

Struttura minima consigliata:

\textbf{Model:} definisci la struttura con \fn{defstruct}.

\textbf{Archivio:} una hash table globale \fn{*archivio*}
con chiave = identificatore univoco.

\textbf{Salvataggio:} con \fn{with-open-file} scrivi ogni
record come S-expression con \fn{print} (leggibile con \fn{read}).
\begin{lstlisting}
(with-open-file (s file :direction :output
                        :if-exists :supersede)
  (maphash (lambda (k v)
             (print (list k v) s))
           *archivio*))
\end{lstlisting}

\textbf{Caricamento:} leggi le S-expression con \fn{read}
in un loop finche' non si raggiunge \fn{EOF}.

\small\color{black!55}\hyperref[sol:148]{$\rightarrow$ Soluzione}\normalsize\color{black}
\end{hintbox}


% =============================================================
%  04_soluzioni.tex — Soluzioni (148 esercizi)
% =============================================================
\chapter*{Soluzioni}
\addcontentsline{toc}{chapter}{Soluzioni}
\markboth{Soluzioni}{Soluzioni}
% Sezioni non numerate per evitare conflitti con Esercizi nell'indice

\noindent\small\color{black!65}
\unaStella/\dueStelle: solo codice.\enspace
\treStelle: commenti inline.\enspace
\quattroStelle/\cinqueStelle: box ragionamento + codice commentato.
\color{black}

\vspace{8pt}

% =============================================================
\section*{Fondamenti}
% =============================================================

\label{sol:1}
\begin{solbox}{1}
\small\color{black!55}\hyperref[ex:1]{$\leftarrow$ Esercizio}\normalsize\color{black}
\begin{lstlisting}
(defun add (a b) (+ a b))
\end{lstlisting}
\end{solbox}

\label{sol:2}
\begin{solbox}{2}
\small\color{black!55}\hyperref[ex:2]{$\leftarrow$ Esercizio}\normalsize\color{black}
\begin{lstlisting}
(defun circle-area (r)
  (* pi r r))

(defun circle-circumference (r)
  (* 2 pi r))
\end{lstlisting}
\end{solbox}

\label{sol:3}
\begin{solbox}{3}
\small\color{black!55}\hyperref[ex:3]{$\leftarrow$ Esercizio}\normalsize\color{black}
\begin{lstlisting}
(defun celsius->fahrenheit (c)
  (+ (* c 9/5) 32))

(defun fahrenheit->celsius (f)
  (* (- f 32) 5/9))
\end{lstlisting}
\end{solbox}

\label{sol:4}
\begin{solbox}{4}
\small\color{black!55}\hyperref[ex:4]{$\leftarrow$ Esercizio}\normalsize\color{black}
\begin{lstlisting}
(defun ipotenusa (a b)
  (let ((a2 (* a a))
        (b2 (* b b)))
    (sqrt (+ a2 b2))))
\end{lstlisting}
\end{solbox}

\label{sol:5}
\begin{solbox}{5}
\small\color{black!55}\hyperref[ex:5]{$\leftarrow$ Esercizio}\normalsize\color{black}
\begin{lstlisting}
(defun distanza-euclidea (x1 y1 x2 y2)
  (let ((dx (- x2 x1))
        (dy (- y2 y1)))
    (sqrt (+ (* dx dx) (* dy dy)))))
\end{lstlisting}
\end{solbox}

\label{sol:6}
\begin{solbox}{6}
\small\color{black!55}\hyperref[ex:6]{$\leftarrow$ Esercizio}\normalsize\color{black}
\begin{lstlisting}
(defun pari-p (n)   (= (mod n 2) 0))
(defun dispari-p (n) (= (mod n 2) 1))
\end{lstlisting}
\end{solbox}

\label{sol:7}
\begin{solbox}{7}
\small\color{black!55}\hyperref[ex:7]{$\leftarrow$ Esercizio}\normalsize\color{black}
\begin{lstlisting}
(defun divisibile-p (a b)
  (= (mod a b) 0))
\end{lstlisting}
\end{solbox}

\label{sol:8}
\begin{solbox}{8}
\small\color{black!55}\hyperref[ex:8]{$\leftarrow$ Esercizio}\normalsize\color{black}
\begin{lstlisting}
(defun segno (n)
  (cond ((< n 0) 'negativo)
        ((= n 0) 'zero)
        (t       'positivo)))
\end{lstlisting}
\end{solbox}

\label{sol:9}
\begin{solbox}{9}
\small\color{black!55}\hyperref[ex:9]{$\leftarrow$ Esercizio}\normalsize\color{black}
\begin{lstlisting}
(defun max-tre (a b c)
  (cond ((and (>= a b) (>= a c)) a)
        ((>= b c)                 b)
        (t                        c)))
\end{lstlisting}
\end{solbox}

\label{sol:10}
\begin{solbox}{10}
\small\color{black!55}\hyperref[ex:10]{$\leftarrow$ Esercizio}\normalsize\color{black}
\begin{lstlisting}
(defun triangolo-valido-p (a b c)
  (and (< a (+ b c))
       (< b (+ a c))
       (< c (+ a b))))
\end{lstlisting}
\end{solbox}

\label{sol:11}
\begin{solbox}{11}
\small\color{black!55}\hyperref[ex:11]{$\leftarrow$ Esercizio}\normalsize\color{black}
\begin{lstlisting}
(defun bisestile-p (anno)
  (or (and (= (mod anno 4)   0)
           (/= (mod anno 100) 0))
      (= (mod anno 400) 0)))
\end{lstlisting}
\end{solbox}

\label{sol:12}
\begin{solbox}{12}
\small\color{black!55}\hyperref[ex:12]{$\leftarrow$ Esercizio}\normalsize\color{black}
\begin{lstlisting}
(defun giorni-mese (mese anno)
  (cond ((member mese '(1 3 5 7 8 10 12)) 31)
        ((member mese '(4 6 9 11))        30)
        ((bisestile-p anno)               29)
        (t                                28)))
\end{lstlisting}
\end{solbox}

\label{sol:13}
\begin{solbox}{13}
\small\color{black!55}\hyperref[ex:13]{$\leftarrow$ Esercizio}\normalsize\color{black}
\begin{lstlisting}
(defun voto->giudizio (p)
  (cond ((>= p 90) "ottimo")
        ((>= p 75) "buono")
        ((>= p 60) "sufficiente")
        (t         "insufficiente")))
\end{lstlisting}
\end{solbox}

\label{sol:14}
\begin{solbox}{14}
\small\color{black!55}\hyperref[ex:14]{$\leftarrow$ Esercizio}\normalsize\color{black}
\begin{lstlisting}
(defun terna-pitagorica-p (a b c)
  (= (* c c) (+ (* a a) (* b b))))
\end{lstlisting}
\end{solbox}

\label{sol:15}
\begin{solbox}{15}
\small\color{black!55}\hyperref[ex:15]{$\leftarrow$ Esercizio}\normalsize\color{black}
\begin{lstlisting}
(defun radici-quadratiche (a b c)
  (let* ((delta (- (* b b) (* 4 a c))))
    (if (< delta 0)
        nil
        (let* ((sd (sqrt delta))
               (r1 (/ (+ (- b) sd) (* 2 a)))
               (r2 (/ (- (- b) sd) (* 2 a))))
          (list r1 r2)))))
\end{lstlisting}
\end{solbox}

\label{sol:16}
\begin{solbox}{16}
\small\color{black!55}\hyperref[ex:16]{$\leftarrow$ Esercizio}\normalsize\color{black}
\begin{lstlisting}
(defun calcola (op a b)
  (case op
    (+  (+ a b))
    (-  (- a b))
    (*  (* a b))
    (/  (if (= b 0) nil (/ a b)))))
\end{lstlisting}
\end{solbox}

\label{sol:17}
\begin{solbox}{17}
\small\color{black!55}\hyperref[ex:17]{$\leftarrow$ Esercizio}\enspace\textcolor{black!30}{|}\enspace\hyperref[hint:17]{$\leftarrow$ Suggerimento}\normalsize\color{black}
\begin{lstlisting}
(defun somma-cifre (n)
  (if (= n 0)
      0
      (+ (mod n 10)
         (somma-cifre (floor n 10)))))
\end{lstlisting}
\end{solbox}

\label{sol:18}
\begin{solbox}{18}
\small\color{black!55}\hyperref[ex:18]{$\leftarrow$ Esercizio}\normalsize\color{black}
\begin{lstlisting}
(defun conta-fino-a (n)
  (dotimes (i n)
    (print (1+ i))))
\end{lstlisting}
\end{solbox}

\label{sol:19}
\begin{solbox}{19}
\small\color{black!55}\hyperref[ex:19]{$\leftarrow$ Esercizio}\normalsize\color{black}
\begin{lstlisting}
(defun tabellina (n)
  (dotimes (i 10)
    (format t "~D x ~D = ~D~%"
            n (1+ i) (* n (1+ i)))))
\end{lstlisting}
\end{solbox}

\label{sol:20}
\begin{solbox}{20}
\small\color{black!55}\hyperref[ex:20]{$\leftarrow$ Esercizio}\normalsize\color{black}
\begin{lstlisting}
(defun sconto (prezzo)
  (cond ((>= prezzo 200) (* prezzo 0.8))
        ((>= prezzo 100) (* prezzo 0.9))
        (t               prezzo)))
\end{lstlisting}
\end{solbox}

\label{sol:21}
\begin{solbox}{21}
\small\color{black!55}\hyperref[ex:21]{$\leftarrow$ Esercizio}\normalsize\color{black}
\begin{lstlisting}
(defun paga-settimanale (ore tariffa)
  (if (<= ore 40)
      (* ore tariffa)
      (+ (* 40 tariffa)
         (* (- ore 40) tariffa 3/2))))
\end{lstlisting}
\end{solbox}

% =============================================================
\section*{Liste}
% =============================================================

\label{sol:22}
\begin{solbox}{22}
\small\color{black!55}\hyperref[ex:22]{$\leftarrow$ Esercizio}\normalsize\color{black}
\begin{lstlisting}
;; Parte A: con cons
(cons 'a (cons (cons 'b (cons 'c nil)) (cons 'd nil)))
; -> (A (B C) D)

;; Parte B: con list
(list 'a (list 'b 'c) 'd)
; -> (A (B C) D)
\end{lstlisting}
\end{solbox}

\label{sol:23}
\begin{solbox}{23}
\small\color{black!55}\hyperref[ex:23]{$\leftarrow$ Esercizio}\normalsize\color{black}
\begin{lstlisting}
(car (cdr (cdr '(a b c d e))))
; -> C
\end{lstlisting}
\end{solbox}

\label{sol:24}
\begin{solbox}{24}
\small\color{black!55}\hyperref[ex:24]{$\leftarrow$ Esercizio}\normalsize\color{black}
\begin{lstlisting}
(defun first-element (lst)
  (if (null lst) nil (nth 0 lst)))

(defun rest-elements (lst)
  (if (null lst) nil (nthcdr 1 lst)))
\end{lstlisting}
\end{solbox}

\label{sol:25}
\begin{solbox}{25}
\small\color{black!55}\hyperref[ex:25]{$\leftarrow$ Esercizio}\normalsize\color{black}
\begin{lstlisting}
(defun primo-e-ultimo (lst)
  (if (null lst)
      nil
      (list (car lst) (car (last lst)))))
\end{lstlisting}
\end{solbox}

\label{sol:26}
\begin{solbox}{26}
\small\color{black!55}\hyperref[ex:26]{$\leftarrow$ Esercizio}\normalsize\color{black}
\begin{lstlisting}
(defun scambia-primi-due (lst)
  (if (or (null lst) (null (cdr lst)))
      lst
      (cons (cadr lst)
            (cons (car lst)
                  (cddr lst)))))
\end{lstlisting}
\end{solbox}

\label{sol:27}
\begin{solbox}{27}
\small\color{black!55}\hyperref[ex:27]{$\leftarrow$ Esercizio}\normalsize\color{black}
\begin{lstlisting}
(defun length* (lst)
  (if (null lst)
      0
      (1+ (length* (cdr lst)))))
\end{lstlisting}
\end{solbox}

\label{sol:28}
\begin{solbox}{28}
\small\color{black!55}\hyperref[ex:28]{$\leftarrow$ Esercizio}\normalsize\color{black}
\begin{lstlisting}
(defun sum-list (lst)
  (if (null lst)
      0
      (+ (car lst) (sum-list (cdr lst)))))
\end{lstlisting}
\end{solbox}

\label{sol:29}
\begin{solbox}{29}
\small\color{black!55}\hyperref[ex:29]{$\leftarrow$ Esercizio}\normalsize\color{black}
\begin{lstlisting}
(defun product-list (lst)
  (if (null lst)
      1
      (* (car lst) (product-list (cdr lst)))))
\end{lstlisting}
\end{solbox}

\label{sol:30}
\begin{solbox}{30}
\small\color{black!55}\hyperref[ex:30]{$\leftarrow$ Esercizio}\normalsize\color{black}
\begin{lstlisting}
(defun max-list (lst)
  (cond ((null lst) (error "lista vuota"))
        ((null (cdr lst)) (car lst))
        (t (max (car lst) (max-list (cdr lst))))))
\end{lstlisting}
\end{solbox}

\label{sol:31}
\begin{solbox}{31}
\small\color{black!55}\hyperref[ex:31]{$\leftarrow$ Esercizio}\normalsize\color{black}
\begin{lstlisting}
(defun member* (item lst)
  (cond ((null lst)          nil)
        ((equal (car lst) item) t)
        (t (member* item (cdr lst)))))
\end{lstlisting}
\end{solbox}

\label{sol:32}
\begin{solbox}{32}
\small\color{black!55}\hyperref[ex:32]{$\leftarrow$ Esercizio}\normalsize\color{black}
\begin{lstlisting}
(defun count-element (item lst)
  (cond ((null lst) 0)
        ((equal (car lst) item)
         (1+ (count-element item (cdr lst))))
        (t (count-element item (cdr lst)))))
\end{lstlisting}
\end{solbox}

\label{sol:33}
\begin{solbox}{33}
\small\color{black!55}\hyperref[ex:33]{$\leftarrow$ Esercizio}\normalsize\color{black}
\begin{lstlisting}
(defun position* (item lst)
  (labels ((iter (lst i)
             (cond ((null lst) nil)
                   ((equal (car lst) item) i)
                   (t (iter (cdr lst) (1+ i))))))
    (iter lst 0)))
\end{lstlisting}
\end{solbox}

\label{sol:34}
\begin{solbox}{34}
\small\color{black!55}\hyperref[ex:34]{$\leftarrow$ Esercizio}\normalsize\color{black}
\begin{lstlisting}
(defun last-element (lst)
  (cond ((null lst) nil)
        ((null (cdr lst)) (car lst))
        (t (last-element (cdr lst)))))
\end{lstlisting}
\end{solbox}

\label{sol:35}
\begin{solbox}{35}
\small\color{black!55}\hyperref[ex:35]{$\leftarrow$ Esercizio}\normalsize\color{black}
\begin{lstlisting}
(defun incrementa-tutti (lst)
  (if (null lst)
      nil
      (cons (1+ (car lst))
            (incrementa-tutti (cdr lst)))))
\end{lstlisting}
\end{solbox}

\label{sol:36}
\begin{solbox}{36}
\small\color{black!55}\hyperref[ex:36]{$\leftarrow$ Esercizio}\normalsize\color{black}
\begin{lstlisting}
(defun solo-pari (lst)
  (cond ((null lst) nil)
        ((evenp (car lst))
         (cons (car lst) (solo-pari (cdr lst))))
        (t (solo-pari (cdr lst)))))
\end{lstlisting}
\end{solbox}

\label{sol:37}
\begin{solbox}{37}
\small\color{black!55}\hyperref[ex:37]{$\leftarrow$ Esercizio}\normalsize\color{black}
\begin{lstlisting}
(defun elimina-tutte (item lst)
  (cond ((null lst) nil)
        ((equal (car lst) item)
         (elimina-tutte item (cdr lst)))
        (t (cons (car lst)
                 (elimina-tutte item (cdr lst))))))
\end{lstlisting}
\end{solbox}

\label{sol:38}
\begin{solbox}{38}
\small\color{black!55}\hyperref[ex:38]{$\leftarrow$ Esercizio}\normalsize\color{black}
\begin{lstlisting}
(defun seleziona-dispari (lst)
  (cond ((null lst)       nil)
        ((null (cdr lst)) (list (car lst)))
        (t (cons (car lst)
                 (seleziona-dispari (cddr lst))))))
\end{lstlisting}
\end{solbox}

\label{sol:39}
\begin{solbox}{39}
\small\color{black!55}\hyperref[ex:39]{$\leftarrow$ Esercizio}\normalsize\color{black}
\begin{lstlisting}
;; Versione ricorsiva semplice (non tail-recursive)
(defun reverse*-naive (lst)
  (if (null lst)
      nil
      (append (reverse*-naive (cdr lst))
              (list (car lst)))))

;; Versione con accumulatore (tail-recursive)
(defun reverse* (lst)
  (labels ((iter (lst acc)
             (if (null lst)
                 acc
                 (iter (cdr lst) (cons (car lst) acc)))))
    (iter lst nil)))
\end{lstlisting}
\end{solbox}

\label{sol:40}
\begin{solbox}{40}
\small\color{black!55}\hyperref[ex:40]{$\leftarrow$ Esercizio}\enspace\textcolor{black!30}{|}\enspace\hyperref[hint:40]{$\leftarrow$ Suggerimento}\normalsize\color{black}
\begin{lstlisting}
(defun append* (lst1 lst2)
  (if (null lst1)
      lst2
      (cons (car lst1)
            (append* (cdr lst1) lst2))))
\end{lstlisting}
\end{solbox}

\label{sol:41}
\begin{solbox}{41}
\small\color{black!55}\hyperref[ex:41]{$\leftarrow$ Esercizio}\enspace\textcolor{black!30}{|}\enspace\hyperref[hint:41]{$\leftarrow$ Suggerimento}\normalsize\color{black}
\begin{lstlisting}
(defun remove-first (item lst)
  (cond ((null lst) nil)
        ((equal (car lst) item) (cdr lst))
        (t (cons (car lst)
                 (remove-first item (cdr lst))))))
\end{lstlisting}
\end{solbox}

\label{sol:42}
\begin{solbox}{42}
\small\color{black!55}\hyperref[ex:42]{$\leftarrow$ Esercizio}\normalsize\color{black}
\begin{lstlisting}
(defun replace-first (old new lst)
  (cond ((null lst) nil)
        ((equal (car lst) old)
         (cons new (cdr lst)))
        (t (cons (car lst)
                 (replace-first old new (cdr lst))))))
\end{lstlisting}
\end{solbox}

\label{sol:43}
\begin{solbox}{43}
\small\color{black!55}\hyperref[ex:43]{$\leftarrow$ Esercizio}\normalsize\color{black}
\begin{lstlisting}
(defun palindroma-p (lst)
  (equal lst (reverse* lst)))
\end{lstlisting}
\end{solbox}

\label{sol:44}
\begin{solbox}{44}
\small\color{black!55}\hyperref[ex:44]{$\leftarrow$ Esercizio}\normalsize\color{black}
\begin{lstlisting}
(defun range (a b)
  (if (> a b)
      nil
      (cons a (range (1+ a) b))))
\end{lstlisting}
\end{solbox}

\label{sol:45}
\begin{solbox}{45}
\small\color{black!55}\hyperref[ex:45]{$\leftarrow$ Esercizio}\normalsize\color{black}
\begin{lstlisting}
(defun duplica-elementi (lst)
  (if (null lst)
      nil
      (cons (car lst)
            (cons (car lst)
                  (duplica-elementi (cdr lst))))))
\end{lstlisting}
\end{solbox}

\label{sol:46}
\begin{solbox}{46}
\small\color{black!55}\hyperref[ex:46]{$\leftarrow$ Esercizio}\normalsize\color{black}
\begin{lstlisting}
(defun schiaccia-coppia (alist)
  (mapcar #'car alist))
\end{lstlisting}
\end{solbox}

\label{sol:47}
\begin{solbox}{47}
\small\color{black!55}\hyperref[ex:47]{$\leftarrow$ Esercizio}\normalsize\color{black}
\begin{lstlisting}
(defun prefisso-p (sub lst)
  (cond ((null sub) t)
        ((null lst)  nil)
        ((equal (car sub) (car lst))
         (prefisso-p (cdr sub) (cdr lst)))
        (t nil)))
\end{lstlisting}
\end{solbox}

\label{sol:48}
\begin{solbox}{48}
\small\color{black!55}\hyperref[ex:48]{$\leftarrow$ Esercizio}\normalsize\color{black}
\begin{lstlisting}
(defun intercala (lst1 lst2)
  (if (null lst1)
      nil
      (cons (car lst1)
            (cons (car lst2)
                  (intercala (cdr lst1) (cdr lst2))))))
\end{lstlisting}
\end{solbox}

\label{sol:49}
\begin{solbox}{49}
\small\color{black!55}\hyperref[ex:49]{$\leftarrow$ Esercizio}\enspace\textcolor{black!30}{|}\enspace\hyperref[hint:49]{$\leftarrow$ Suggerimento}\normalsize\color{black}
\begin{lstlisting}
(defun rle-encode (lst)
  (if (null lst)
      nil
      (labels ((iter (lst curr count acc)
                 (cond
                   ((null lst)
                    (reverse (cons (list count curr) acc)))
                   ((equal (car lst) curr)
                    (iter (cdr lst) curr (1+ count) acc))
                   (t
                    (iter (cdr lst) (car lst) 1
                          (cons (list count curr) acc))))))
        (iter (cdr lst) (car lst) 1 nil))))

(defun rle-decode (lst)
  (if (null lst)
      nil
      (let ((n (caar lst))
            (e (cadar lst)))
        (append (loop repeat n collect e)
                (rle-decode (cdr lst))))))
\end{lstlisting}
\end{solbox}

\label{sol:50}
\begin{solbox}{50}
\small\color{black!55}\hyperref[ex:50]{$\leftarrow$ Esercizio}\enspace\textcolor{black!30}{|}\enspace\hyperref[hint:50]{$\leftarrow$ Suggerimento}\normalsize\color{black}
\begin{lstlisting}
(defun remove-duplicates* (lst)
  (labels ((iter (lst seen)
             (cond
               ((null lst) nil)
               ((member (car lst) seen)
                (iter (cdr lst) seen))
               (t (cons (car lst)
                        (iter (cdr lst)
                              (cons (car lst) seen)))))))
    (iter lst nil)))
\end{lstlisting}
\end{solbox}

% =============================================================
\section*{Ricorsione avanzata}
% =============================================================

\label{sol:51}
\begin{solbox}{51}
\small\color{black!55}\hyperref[ex:51]{$\leftarrow$ Esercizio}\normalsize\color{black}
\begin{lstlisting}
(defun fattoriale (n)
  ;; labels definisce iter localmente (tail-recursive)
  (labels ((iter (n acc)
             (if (= n 0)
                 acc                        ; caso base: restituisci acc
                 (iter (1- n) (* n acc))))) ; accumula n nel prodotto
    (iter n 1)))                            ; avvia con acc = 1
\end{lstlisting}
\end{solbox}

\label{sol:52}
\begin{solbox}{52}
\small\color{black!55}\hyperref[ex:52]{$\leftarrow$ Esercizio}\normalsize\color{black}
\begin{lstlisting}
(defun reverse-tr (lst)
  ;; l'accumulatore accumula in ordine inverso
  (labels ((iter (lst acc)
             (if (null lst)
                 acc
                 (iter (cdr lst) (cons (car lst) acc)))))
    (iter lst nil)))
\end{lstlisting}
\end{solbox}

\label{sol:53}
\begin{solbox}{53}
\small\color{black!55}\hyperref[ex:53]{$\leftarrow$ Esercizio}\normalsize\color{black}
\begin{lstlisting}
(defun somma-range (a b)
  (labels ((iter (i acc)
             (if (> i b)
                 acc
                 (iter (1+ i) (+ acc i)))))
    (iter a 0)))
\end{lstlisting}
\end{solbox}

\label{sol:54}
\begin{solbox}{54}
\small\color{black!55}\hyperref[ex:54]{$\leftarrow$ Esercizio}\enspace\textcolor{black!30}{|}\enspace\hyperref[hint:54]{$\leftarrow$ Suggerimento}\normalsize\color{black}
\begin{lstlisting}
(defun fibonacci-fast (n)
  ;; a = F(i), b = F(i+1); a ogni passo a <- b, b <- a+b
  (labels ((iter (i a b)
             (if (= i 0)
                 a
                 (iter (1- i) b (+ a b)))))
    (iter n 0 1)))
\end{lstlisting}
\end{solbox}

\label{sol:55}
\begin{solbox}{55}
\small\color{black!55}\hyperref[ex:55]{$\leftarrow$ Esercizio}\enspace\textcolor{black!30}{|}\enspace\hyperref[hint:55]{$\leftarrow$ Suggerimento}\normalsize\color{black}
\begin{lstlisting}
(defun flatten* (tree)
  (cond
    ((null tree) nil)
    ;; atomo (foglia): avvolgilo in lista
    ((atom tree) (list tree))
    ;; nodo interno: appiattisci car e cdr separatamente
    (t (append (flatten* (car tree))
               (flatten* (cdr tree))))))
\end{lstlisting}
\end{solbox}

\label{sol:56}
\begin{solbox}{56}
\small\color{black!55}\hyperref[ex:56]{$\leftarrow$ Esercizio}\enspace\textcolor{black!30}{|}\enspace\hyperref[hint:56]{$\leftarrow$ Suggerimento}\normalsize\color{black}
\begin{lstlisting}
(defun tree-depth (tree)
  (cond
    ((null tree) 0)
    ((atom tree) 0)
    ;; profondita' = max(1 + profondita' car, profondita' cdr)
    (t (max (1+ (tree-depth (car tree)))
            (tree-depth (cdr tree))))))
\end{lstlisting}
\end{solbox}

\label{sol:57}
\begin{solbox}{57}
\small\color{black!55}\hyperref[ex:57]{$\leftarrow$ Esercizio}\normalsize\color{black}
\begin{lstlisting}
(defun count-atoms (tree)
  (cond ((null tree) 0)
        ((atom tree) 1)
        (t (+ (count-atoms (car tree))
              (count-atoms (cdr tree))))))
\end{lstlisting}
\end{solbox}

\label{sol:58}
\begin{solbox}{58}
\small\color{black!55}\hyperref[ex:58]{$\leftarrow$ Esercizio}\normalsize\color{black}
\begin{lstlisting}
(defun sum-tree (tree)
  (cond ((null tree)   0)
        ((numberp tree) tree)
        ((atom tree)    0)
        (t (+ (sum-tree (car tree))
              (sum-tree (cdr tree))))))
\end{lstlisting}
\end{solbox}

\label{sol:59}
\begin{solbox}{59}
\small\color{black!55}\hyperref[ex:59]{$\leftarrow$ Esercizio}\normalsize\color{black}
\begin{lstlisting}
(defun tree-member-p (item tree)
  (cond ((null tree) nil)
        ((equal tree item) t)
        ((atom tree) nil)
        (t (or (tree-member-p item (car tree))
               (tree-member-p item (cdr tree))))))
\end{lstlisting}
\end{solbox}

\label{sol:60}
\begin{solbox}{60}
\small\color{black!55}\hyperref[ex:60]{$\leftarrow$ Esercizio}\enspace\textcolor{black!30}{|}\enspace\hyperref[hint:60]{$\leftarrow$ Suggerimento}\normalsize\color{black}
\begin{lstlisting}
(defun sostituisci-profondo (old new tree)
  (cond ((null tree) nil)
        ((equal tree old) new)   ; sostituzione
        ((atom tree) tree)       ; atomo diverso: lascia invariato
        (t (cons (sostituisci-profondo old new (car tree))
                 (sostituisci-profondo old new (cdr tree))))))
\end{lstlisting}
\end{solbox}

\label{sol:61}
\begin{solbox}{61}
\small\color{black!55}\hyperref[ex:61]{$\leftarrow$ Esercizio}\enspace\textcolor{black!30}{|}\enspace\hyperref[hint:61]{$\leftarrow$ Suggerimento}\normalsize\color{black}
\begin{lstlisting}
(defun power-fast (base exp)
  (cond ((= exp 0) 1)
        ((evenp exp)
         ;; evita doppia ricorsione con let
         (let ((half (power-fast base (/ exp 2))))
           (* half half)))
        (t (* base (power-fast base (1- exp))))))
\end{lstlisting}
\end{solbox}

\label{sol:62}
\begin{solbox}{62}
\small\color{black!55}\hyperref[ex:62]{$\leftarrow$ Esercizio}\normalsize\color{black}
\begin{lstlisting}
(defun mcd* (a b)
  (if (= b 0) a (mcd* b (mod a b))))

(defun mcm* (a b)
  (/ (* a b) (mcd* a b)))
\end{lstlisting}
\end{solbox}

\label{sol:63}
\begin{solbox}{63}
\small\color{black!55}\hyperref[ex:63]{$\leftarrow$ Esercizio}\normalsize\color{black}
\begin{lstlisting}
(defun primo-p (n)
  (if (< n 2)
      nil
      (loop for d from 2 to (isqrt n)
            never (= (mod n d) 0))))
\end{lstlisting}
\end{solbox}

\label{sol:64}
\begin{solbox}{64}
\small\color{black!55}\hyperref[ex:64]{$\leftarrow$ Esercizio}\enspace\textcolor{black!30}{|}\enspace\hyperref[hint:64]{$\leftarrow$ Suggerimento}\normalsize\color{black}
\begin{lstlisting}
(defun crivello (n)
  ;; array di candidati: T = possibile primo
  (let ((sieve (make-array (1+ n) :initial-element t)))
    (setf (aref sieve 0) nil
          (aref sieve 1) nil)
    ;; elimina multipli di ogni primo trovato
    (loop for i from 2 to (isqrt n)
          when (aref sieve i)
          do (loop for j from (* i i) to n by i
                   do (setf (aref sieve j) nil)))
    ;; raccogli gli indici rimasti T
    (loop for i from 2 to n
          when (aref sieve i) collect i)))
\end{lstlisting}
\end{solbox}

\label{sol:65}
\begin{solbox}{65}
\small\color{black!55}\hyperref[ex:65]{$\leftarrow$ Esercizio}\enspace\textcolor{black!30}{|}\enspace\hyperref[hint:65]{$\leftarrow$ Suggerimento}\normalsize\color{black}
\begin{lstlisting}
(defun prime-factors (n)
  ;; prova divisori da 2 in su; non avanzare d se divide
  (labels ((iter (n d)
             (cond ((> (* d d) n)
                    (if (> n 1) (list n) nil))
                   ((= (mod n d) 0)
                    (cons d (iter (/ n d) d)))
                   (t (iter n (1+ d))))))
    (iter n 2)))
\end{lstlisting}
\end{solbox}

\label{sol:66}
\begin{solbox}{66}
\small\color{black!55}\hyperref[ex:66]{$\leftarrow$ Esercizio}\normalsize\color{black}
\begin{lstlisting}
(defun sublist (lst inizio fine)
  (let ((len (- fine inizio)))
    (if (<= len 0)
        nil
        (loop for i from inizio below fine
              collect (nth i lst)))))
\end{lstlisting}
\end{solbox}

\label{sol:67}
\begin{solbox}{67}
\small\color{black!55}\hyperref[ex:67]{$\leftarrow$ Esercizio}\enspace\textcolor{black!30}{|}\enspace\hyperref[hint:67]{$\leftarrow$ Suggerimento}\normalsize\color{black}
\begin{lstlisting}
(defun raggruppa (lst n)
  (if (null lst)
      nil
      (cons (loop for x in lst
                  for i below n collect x)
            (raggruppa (nthcdr n lst) n))))
\end{lstlisting}
\end{solbox}

\label{sol:68}
\begin{solbox}{68}
\small\color{black!55}\hyperref[ex:68]{$\leftarrow$ Esercizio}\normalsize\color{black}
\begin{lstlisting}
(defun interseca (lst1 lst2)
  (remove-if-not (lambda (x) (member x lst2)) lst1))
\end{lstlisting}
\end{solbox}

% =============================================================
\section*{Higher-order e closure}
% =============================================================

\label{sol:69}
\begin{solbox}{69}
\small\color{black!55}\hyperref[ex:69]{$\leftarrow$ Esercizio}\normalsize\color{black}
\begin{lstlisting}
(defun square-list (lst)
  (mapcar (lambda (x) (* x x)) lst))
\end{lstlisting}
\end{solbox}

\label{sol:70}
\begin{solbox}{70}
\small\color{black!55}\hyperref[ex:70]{$\leftarrow$ Esercizio}\normalsize\color{black}
\begin{lstlisting}
(defun even-list (lst)
  (remove-if-not #'evenp lst))
\end{lstlisting}
\end{solbox}

\label{sol:71}
\begin{solbox}{71}
\small\color{black!55}\hyperref[ex:71]{$\leftarrow$ Esercizio}\normalsize\color{black}
\begin{lstlisting}
(defun vocali-iniziali (lst)
  (remove-if-not
   (lambda (s)
     (member (char-downcase (char s 0))
             '(#\a #\e #\i #\o #\u)))
   lst))
\end{lstlisting}
\end{solbox}

\label{sol:72}
\begin{solbox}{72}
\small\color{black!55}\hyperref[ex:72]{$\leftarrow$ Esercizio}\normalsize\color{black}
\begin{lstlisting}
(defun sum-reduce (lst)
  (reduce #'+ lst :initial-value 0))
\end{lstlisting}
\end{solbox}

\label{sol:73}
\begin{solbox}{73}
\small\color{black!55}\hyperref[ex:73]{$\leftarrow$ Esercizio}\normalsize\color{black}
\begin{lstlisting}
(defun product-reduce (lst)
  (reduce #'* lst :initial-value 1))
\end{lstlisting}
\end{solbox}

\label{sol:74}
\begin{solbox}{74}
\small\color{black!55}\hyperref[ex:74]{$\leftarrow$ Esercizio}\normalsize\color{black}
\begin{lstlisting}
(defun reduce-min (lst)
  (if (null lst)
      (error "lista vuota")
      (reduce #'min lst)))

(defun reduce-max (lst)
  (if (null lst)
      (error "lista vuota")
      (reduce #'max lst)))
\end{lstlisting}
\end{solbox}

\label{sol:75}
\begin{solbox}{75}
\small\color{black!55}\hyperref[ex:75]{$\leftarrow$ Esercizio}\enspace\textcolor{black!30}{|}\enspace\hyperref[hint:75]{$\leftarrow$ Suggerimento}\normalsize\color{black}
\begin{lstlisting}
(defun cifre->intero (lst)
  ;; a ogni passo: acc = acc*10 + cifra
  (reduce (lambda (acc d) (+ (* acc 10) d))
          lst :initial-value 0))
\end{lstlisting}
\end{solbox}

\label{sol:76}
\begin{solbox}{76}
\small\color{black!55}\hyperref[ex:76]{$\leftarrow$ Esercizio}\normalsize\color{black}
\begin{lstlisting}
(defun appiattisci-un-livello (lst)
  (reduce #'append lst :initial-value nil))
\end{lstlisting}
\end{solbox}

\label{sol:77}
\begin{solbox}{77}
\small\color{black!55}\hyperref[ex:77]{$\leftarrow$ Esercizio}\enspace\textcolor{black!30}{|}\enspace\hyperref[hint:77]{$\leftarrow$ Suggerimento}\normalsize\color{black}
\begin{lstlisting}
(defun horner (coefficienti x)
  ;; fold: acc = acc*x + c
  (reduce (lambda (acc c) (+ (* acc x) c))
          coefficienti :initial-value 0))
\end{lstlisting}
\end{solbox}

\label{sol:78}
\begin{solbox}{78}
\small\color{black!55}\hyperref[ex:78]{$\leftarrow$ Esercizio}\normalsize\color{black}
\begin{lstlisting}
(defun prodotto-scalare (v1 v2)
  (apply #'+ (mapcar #'* v1 v2)))
\end{lstlisting}
\end{solbox}

\label{sol:79}
\begin{solbox}{79}
\small\color{black!55}\hyperref[ex:79]{$\leftarrow$ Esercizio}\normalsize\color{black}
\begin{lstlisting}
;; mapcan applica f a ogni elemento e concatena le liste risultanti;
;; mapcar le raccoglie senza appiattire.
(defun raddoppia (lst)
  (mapcan (lambda (x) (list x x)) lst))
\end{lstlisting}
\end{solbox}

\label{sol:80}
\begin{solbox}{80}
\small\color{black!55}\hyperref[ex:80]{$\leftarrow$ Esercizio}\normalsize\color{black}
\begin{lstlisting}
(defun ordina-per-punteggio (lst)
  (sort (copy-list lst)
        (lambda (a b) (> (cdr a) (cdr b)))))
\end{lstlisting}
\end{solbox}

\label{sol:81}
\begin{solbox}{81}
\small\color{black!55}\hyperref[ex:81]{$\leftarrow$ Esercizio}\normalsize\color{black}
\begin{lstlisting}
(defun tutti-p (pred lst)
  (cond ((null lst) t)
        ((funcall pred (car lst))
         (tutti-p pred (cdr lst)))
        (t nil)))

(defun qualcuno-p (pred lst)
  (cond ((null lst) nil)
        ((funcall pred (car lst)) t)
        (t (qualcuno-p pred (cdr lst)))))
\end{lstlisting}
\end{solbox}

\label{sol:82}
\begin{solbox}{82}
\small\color{black!55}\hyperref[ex:82]{$\leftarrow$ Esercizio}\normalsize\color{black}
\begin{lstlisting}
(defun count-matching (pred lst)
  (reduce (lambda (acc x)
            (if (funcall pred x) (1+ acc) acc))
          lst :initial-value 0))
\end{lstlisting}
\end{solbox}

\label{sol:83}
\begin{solbox}{83}
\small\color{black!55}\hyperref[ex:83]{$\leftarrow$ Esercizio}\enspace\textcolor{black!30}{|}\enspace\hyperref[hint:83]{$\leftarrow$ Suggerimento}\normalsize\color{black}
\begin{lstlisting}
(defun raggruppa-per (f lst)
  ;; per ogni elemento, trova o crea il gruppo nella alist
  (let ((result nil))
    (dolist (x lst result)
      (let* ((k (funcall f x))
             (bucket (assoc k result)))
        (if bucket
            (setf (cdr bucket) (append (cdr bucket) (list x)))
            (push (cons k (list x)) result))))))
\end{lstlisting}
\end{solbox}

\label{sol:84}
\begin{solbox}{84}
\small\color{black!55}\hyperref[ex:84]{$\leftarrow$ Esercizio}\normalsize\color{black}
\begin{lstlisting}
(defun applica-a-tutti (funzioni x)
  (mapcar (lambda (f) (funcall f x)) funzioni))
\end{lstlisting}
\end{solbox}

\label{sol:85}
\begin{solbox}{85}
\small\color{black!55}\hyperref[ex:85]{$\leftarrow$ Esercizio}\enspace\textcolor{black!30}{|}\enspace\hyperref[hint:85]{$\leftarrow$ Suggerimento}\normalsize\color{black}
\begin{lstlisting}
(defun compose (f g)
  ;; cattura f e g nella closure
  (lambda (x) (funcall f (funcall g x))))
\end{lstlisting}
\end{solbox}

\label{sol:86}
\begin{solbox}{86}
\small\color{black!55}\hyperref[ex:86]{$\leftarrow$ Esercizio}\enspace\textcolor{black!30}{|}\enspace\hyperref[hint:86]{$\leftarrow$ Suggerimento}\normalsize\color{black}
\begin{lstlisting}
(defun partial (funzione &rest argomenti)
  ;; cattura funzione e argomenti nella closure
  (lambda (&rest rest-args)
    (apply funzione (append argomenti rest-args))))
\end{lstlisting}
\end{solbox}

\label{sol:87}
\begin{solbox}{87}
\small\color{black!55}\hyperref[ex:87]{$\leftarrow$ Esercizio}\enspace\textcolor{black!30}{|}\enspace\hyperref[hint:87]{$\leftarrow$ Suggerimento}\normalsize\color{black}
\begin{lstlisting}
(defun make-counter (&optional (start 0) (step 1))
  ;; curr e' privato alla closure
  (let ((curr start))
    (lambda ()
      (prog1 curr       ; restituisce curr prima dell'incremento
        (incf curr step)))))
\end{lstlisting}
\end{solbox}

\label{sol:88}
\begin{solbox}{88}
\small\color{black!55}\hyperref[ex:88]{$\leftarrow$ Esercizio}\normalsize\color{black}
\begin{lstlisting}
(defun make-accumulator (iniziale)
  (let ((totale iniziale))
    (lambda (x)
      (incf totale x))))
\end{lstlisting}
\end{solbox}

\label{sol:89}
\begin{solbox}{89}
\small\color{black!55}\hyperref[ex:89]{$\leftarrow$ Esercizio}\enspace\textcolor{black!30}{|}\enspace\hyperref[hint:89]{$\leftarrow$ Suggerimento}\normalsize\color{black}
\begin{lstlisting}
(defun make-fib-generator ()
  ;; a e b condividono la stessa closure
  (let ((a 0) (b 1))
    (lambda ()
      (prog1 a
        (psetf a b           ; aggiorna simultaneamente
               b (+ a b))))))
\end{lstlisting}
\end{solbox}

\label{sol:90}
\begin{solbox}{90}
\small\color{black!55}\hyperref[ex:90]{$\leftarrow$ Esercizio}\normalsize\color{black}
\begin{lstlisting}
(defun pipeline (valore funzioni)
  (reduce (lambda (v f) (funcall f v))
          funzioni :initial-value valore))
\end{lstlisting}
\end{solbox}

\label{sol:91}
\begin{solbox}{91}
\small\color{black!55}\hyperref[ex:91]{$\leftarrow$ Esercizio}\enspace\textcolor{black!30}{|}\enspace\hyperref[hint:91]{$\leftarrow$ Suggerimento}\normalsize\color{black}
\begin{lstlisting}
(defun memoize (funzione)
  ;; la cache e' privata alla closure
  (let ((cache (make-hash-table :test #'equal)))
    (lambda (&rest args)
      (multiple-value-bind (val found)
          (gethash args cache)
        (if found
            val
            ;; calcola e memorizza
            (setf (gethash args cache)
                  (apply funzione args)))))))
\end{lstlisting}
\end{solbox}

\label{sol:92}
\begin{solbox}{92}
\small\color{black!55}\hyperref[ex:92]{$\leftarrow$ Esercizio}\normalsize\color{black}
\begin{lstlisting}
;; Mini libreria funzionale (estratto)
(defun map-list (f lst)
  "Applica F a ogni elemento di LST."
  (mapcar f lst))

(defun filter-list (pred lst)
  "Filtra LST mantenendo gli elementi per cui PRED e' vero."
  (remove-if-not pred lst))

(defun find-matching (pred lst)
  "Restituisce il primo elemento che soddisfa PRED."
  (find-if pred lst))

(defun take (n lst)
  "Restituisce i primi N elementi di LST."
  (if (or (= n 0) (null lst))
      nil
      (cons (car lst) (take (1- n) (cdr lst)))))

(defun drop (n lst)
  "Rimuove i primi N elementi di LST."
  (if (or (= n 0) (null lst))
      lst
      (drop (1- n) (cdr lst))))

(defun zip (lst1 lst2)
  "Combina due liste in lista di coppie."
  (if (or (null lst1) (null lst2))
      nil
      (cons (list (car lst1) (car lst2))
            (zip (cdr lst1) (cdr lst2)))))
\end{lstlisting}
\end{solbox}

% =============================================================
\section*{Strutture dati}
% =============================================================

\label{sol:93}
\begin{solbox}{93}
\small\color{black!55}\hyperref[ex:93]{$\leftarrow$ Esercizio}\normalsize\color{black}
\begin{lstlisting}
(defun lookup (key alist)
  (let ((pair (assoc key alist)))
    (if pair (cdr pair) nil)))

(defun set-key (key value alist)
  ;; costruisce nuova alist senza modificare l'originale
  (cons (cons key value)
        (remove-if (lambda (p) (equal (car p) key))
                   alist)))
\end{lstlisting}
\end{solbox}

\label{sol:94}
\begin{solbox}{94}
\small\color{black!55}\hyperref[ex:94]{$\leftarrow$ Esercizio}\enspace\textcolor{black!30}{|}\enspace\hyperref[hint:94]{$\leftarrow$ Suggerimento}\normalsize\color{black}
\begin{lstlisting}
(defun frequencies (lst)
  ;; per ogni elemento, aggiorna o aggiunge la coppia
  (let ((result nil))
    (dolist (x lst result)
      (let ((pair (assoc x result)))
        (if pair
            (incf (cdr pair))
            (push (cons x 1) result))))))
\end{lstlisting}
\end{solbox}

\label{sol:95}
\begin{solbox}{95}
\small\color{black!55}\hyperref[ex:95]{$\leftarrow$ Esercizio}\normalsize\color{black}
\begin{lstlisting}
(defun ht-get (key ht)
  (gethash key ht))

(defun ht-set (key value ht)
  (setf (gethash key ht) value))

(defun ht-delete (key ht)
  (remhash key ht))

(defun ht-keys (ht)
  (loop for k being the hash-keys of ht collect k))
\end{lstlisting}
\end{solbox}

\label{sol:96}
\begin{solbox}{96}
\small\color{black!55}\hyperref[ex:96]{$\leftarrow$ Esercizio}\normalsize\color{black}
\begin{lstlisting}
(defun frequencies-hash (lst)
  (let ((ht (make-hash-table)))
    (dolist (x lst ht)
      (incf (gethash x ht 0)))))
\end{lstlisting}
\end{solbox}

\label{sol:97}
\begin{solbox}{97}
\small\color{black!55}\hyperref[ex:97]{$\leftarrow$ Esercizio}\enspace\textcolor{black!30}{|}\enspace\hyperref[hint:97]{$\leftarrow$ Suggerimento}\normalsize\color{black}
\begin{lstlisting}
(defun frequenza-caratteri (stringa)
  (let ((ht (make-hash-table :test #'eql)))
    (loop for c across stringa
          do (incf (gethash c ht 0)))
    ht))
\end{lstlisting}
\end{solbox}

\label{sol:98}
\begin{solbox}{98}
\small\color{black!55}\hyperref[ex:98]{$\leftarrow$ Esercizio}\enspace\textcolor{black!30}{|}\enspace\hyperref[hint:98]{$\leftarrow$ Suggerimento}\normalsize\color{black}
\begin{lstlisting}
(defun unisci-hash (ht1 ht2)
  (let ((result (make-hash-table :test (hash-table-test ht1))))
    (maphash (lambda (k v) (setf (gethash k result) v)) ht1)
    (maphash (lambda (k v)
               (incf (gethash k result 0) v))
             ht2)
    result))
\end{lstlisting}
\end{solbox}

\label{sol:99}
\begin{solbox}{99}
\small\color{black!55}\hyperref[ex:99]{$\leftarrow$ Esercizio}\enspace\textcolor{black!30}{|}\enspace\hyperref[hint:99]{$\leftarrow$ Suggerimento}\normalsize\color{black}
\begin{lstlisting}
(defun inverti-hash (ht)
  (let ((result (make-hash-table :test #'equal)))
    (maphash (lambda (k v) (setf (gethash v result) k)) ht)
    result))
\end{lstlisting}
\end{solbox}

\label{sol:100}
\begin{solbox}{100}
\small\color{black!55}\hyperref[ex:100]{$\leftarrow$ Esercizio}\normalsize\color{black}
\begin{lstlisting}
(defun alist->hash (alist)
  (let ((ht (make-hash-table)))
    (dolist (pair alist ht)
      (setf (gethash (car pair) ht) (cdr pair)))))

(defun hash->alist (ht)
  (loop for k being the hash-keys of ht
        collect (cons k (gethash k ht))))
\end{lstlisting}
\end{solbox}

\label{sol:101}
\begin{solbox}{101}
\small\color{black!55}\hyperref[ex:101]{$\leftarrow$ Esercizio}\normalsize\color{black}
\begin{lstlisting}
(defun mat-get (m r c) (aref m r c))

(defun mat-set (m r c val) (setf (aref m r c) val))

(defun stampa-matrice (m)
  (let ((n (car (array-dimensions m))))
    (dotimes (r n)
      (dotimes (c n)
        (format t "~4D" (aref m r c)))
      (terpri))))
\end{lstlisting}
\end{solbox}

\label{sol:102}
\begin{solbox}{102}
\small\color{black!55}\hyperref[ex:102]{$\leftarrow$ Esercizio}\enspace\textcolor{black!30}{|}\enspace\hyperref[hint:102]{$\leftarrow$ Suggerimento}\normalsize\color{black}
\begin{lstlisting}
(defun trasposta (m n)
  ;; t[i][j] = m[j][i]
  (let ((result (make-array (list n n))))
    (dotimes (i n result)
      (dotimes (j n)
        (setf (aref result i j)
              (aref m j i))))))
\end{lstlisting}
\end{solbox}

\label{sol:103}
\begin{solbox}{103}
\small\color{black!55}\hyperref[ex:103]{$\leftarrow$ Esercizio}\normalsize\color{black}
\begin{lstlisting}
(defun inverti-array! (v)
  ;; scambia elementi simmetrici senza lista ausiliaria
  (let ((len (length v)))
    (dotimes (i (floor len 2))
      (let ((tmp (aref v i)))
        (setf (aref v i) (aref v (- len 1 i))
              (aref v (- len 1 i)) tmp)))))
\end{lstlisting}
\end{solbox}

\label{sol:104}
\begin{solbox}{104}
\small\color{black!55}\hyperref[ex:104]{$\leftarrow$ Esercizio}\enspace\textcolor{black!30}{|}\enspace\hyperref[hint:104]{$\leftarrow$ Suggerimento}\normalsize\color{black}
\begin{lstlisting}
(defun moltiplica-matrici (a b n)
  (let ((c (make-array (list n n) :initial-element 0)))
    (dotimes (i n c)
      (dotimes (j n)
        (dotimes (k n)
          (incf (aref c i j)
                (* (aref a i k) (aref b k j))))))))
\end{lstlisting}
\end{solbox}

\label{sol:105}
\begin{solbox}{105}
\small\color{black!55}\hyperref[ex:105]{$\leftarrow$ Esercizio}\normalsize\color{black}
\begin{lstlisting}
(defstruct persona
  nome
  (eta 18)
  citta)

(defun crea-persona (nome eta citta)
  (when (not (and (integerp eta) (> eta 0)))
    (error "Eta' non valida: ~A" eta))
  (make-persona :nome nome :eta eta :citta citta))
\end{lstlisting}
\end{solbox}

\label{sol:106}
\begin{solbox}{106}
\small\color{black!55}\hyperref[ex:106]{$\leftarrow$ Esercizio}\normalsize\color{black}
\begin{lstlisting}
(defstruct nodo
  valore
  (sinistra nil)
  (destra   nil))

;; albero:     5
;;            / \
;;           3   8
;;          / \
;;         1   4
(defparameter *albero*
  (make-nodo :valore 5
    :sinistra (make-nodo :valore 3
                :sinistra (make-nodo :valore 1)
                :destra   (make-nodo :valore 4))
    :destra   (make-nodo :valore 8)))
\end{lstlisting}
\end{solbox}

\label{sol:107}
\begin{solbox}{107}
\small\color{black!55}\hyperref[ex:107]{$\leftarrow$ Esercizio}\normalsize\color{black}
\begin{lstlisting}
;; Rappresentazione: (valore sinistra destra)
;; NIL = albero vuoto
(defun make-leaf (v) (list v nil nil))
(defun make-node (v l r) (list v l r))
(defun node-val  (n) (car n))
(defun node-left (n) (cadr n))
(defun node-right (n) (caddr n))
\end{lstlisting}
\end{solbox}

\label{sol:108}
\begin{solbox}{108}
\small\color{black!55}\hyperref[ex:108]{$\leftarrow$ Esercizio}\enspace\textcolor{black!30}{|}\enspace\hyperref[hint:108]{$\leftarrow$ Suggerimento}\normalsize\color{black}
\begin{lstlisting}
(defun bst-insert (value tree)
  ;; restituisce il nuovo albero (non modifica in loco)
  (if (null tree)
      (make-nodo :valore value)
      (cond ((< value (nodo-valore tree))
             (make-nodo
               :valore (nodo-valore tree)
               :sinistra (bst-insert value (nodo-sinistra tree))
               :destra   (nodo-destra tree)))
            ((> value (nodo-valore tree))
             (make-nodo
               :valore (nodo-valore tree)
               :sinistra (nodo-sinistra tree)
               :destra   (bst-insert value (nodo-destra tree))))
            ;; duplicato: non inserire
            (t tree))))

(defun bst-find (value tree)
  (cond ((null tree) nil)
        ((= value (nodo-valore tree)) t)
        ((< value (nodo-valore tree))
         (bst-find value (nodo-sinistra tree)))
        (t (bst-find value (nodo-destra tree)))))
\end{lstlisting}
\end{solbox}

\label{sol:109}
\begin{solbox}{109}
\small\color{black!55}\hyperref[ex:109]{$\leftarrow$ Esercizio}\enspace\textcolor{black!30}{|}\enspace\hyperref[hint:109]{$\leftarrow$ Suggerimento}\normalsize\color{black}
\begin{lstlisting}
(defun preorder (tree)
  (if (null tree) nil
      (append (list (nodo-valore tree))
              (preorder (nodo-sinistra tree))
              (preorder (nodo-destra tree)))))

(defun inorder (tree)
  (if (null tree) nil
      (append (inorder (nodo-sinistra tree))
              (list (nodo-valore tree))
              (inorder (nodo-destra tree)))))

(defun postorder (tree)
  (if (null tree) nil
      (append (postorder (nodo-sinistra tree))
              (postorder (nodo-destra tree))
              (list (nodo-valore tree)))))
\end{lstlisting}
\end{solbox}

\label{sol:110}
\begin{solbox}{110}
\small\color{black!55}\hyperref[ex:110]{$\leftarrow$ Esercizio}\normalsize\color{black}
\begin{lstlisting}
(defun tree-height (tree)
  (if (null tree) 0
      (1+ (max (tree-height (nodo-sinistra tree))
               (tree-height (nodo-destra tree))))))

(defun count-leaves (tree)
  (cond ((null tree) 0)
        ((and (null (nodo-sinistra tree))
              (null (nodo-destra tree))) 1)
        (t (+ (count-leaves (nodo-sinistra tree))
              (count-leaves (nodo-destra tree))))))
\end{lstlisting}
\end{solbox}

\label{sol:111}
\begin{solbox}{111}
\small\color{black!55}\hyperref[ex:111]{$\leftarrow$ Esercizio}\normalsize\color{black}
\begin{lstlisting}
(defun specchia (tree)
  (if (null tree) nil
      (make-nodo :valore   (nodo-valore tree)
                 :sinistra (specchia (nodo-destra tree))
                 :destra   (specchia (nodo-sinistra tree)))))
\end{lstlisting}
\end{solbox}

\label{sol:112}
\begin{solbox}{112}
\small\color{black!55}\hyperref[ex:112]{$\leftarrow$ Esercizio}\enspace\textcolor{black!30}{|}\enspace\hyperref[hint:112]{$\leftarrow$ Suggerimento}\normalsize\color{black}
\begin{lstlisting}
(defun bilanciato-p (tree)
  ;; ritorna altezza se bilanciato, -1 come segnale di errore
  (labels ((check (t)
             (if (null t)
                 0
                 (let ((hl (check (nodo-sinistra t)))
                       (hr (check (nodo-destra t))))
                   (if (or (= hl -1) (= hr -1)
                           (> (abs (- hl hr)) 1))
                       -1
                       (1+ (max hl hr)))))))
    (/= (check tree) -1)))
\end{lstlisting}
\end{solbox}

\label{sol:113}
\begin{solbox}{113}
\small\color{black!55}\hyperref[ex:113]{$\leftarrow$ Esercizio}\normalsize\color{black}
\begin{lstlisting}
;; Pila come lista avvolta in cons-cell per mutabilita'
(defun make-stack () (list nil))

(defun stack-push (item stack)
  (setf (car stack) (cons item (car stack))))

(defun stack-pop (stack)
  (prog1 (caar stack)
    (setf (car stack) (cdar stack))))

(defun stack-peek (stack) (caar stack))

(defun stack-empty-p (stack) (null (car stack)))
\end{lstlisting}
\end{solbox}

\label{sol:114}
\begin{solbox}{114}
\small\color{black!55}\hyperref[ex:114]{$\leftarrow$ Esercizio}\enspace\textcolor{black!30}{|}\enspace\hyperref[hint:114]{$\leftarrow$ Suggerimento}\normalsize\color{black}
\begin{lstlisting}
;; Coda con due liste: testa (estrai) e fondo (aggiungi)
(defstruct queue (head nil) (tail nil))

(defun enqueue (item q)
  (setf (queue-tail q) (cons item (queue-tail q))))

(defun dequeue (q)
  ;; se testa vuota, ribalta il fondo
  (when (null (queue-head q))
    (setf (queue-head q) (reverse (queue-tail q))
          (queue-tail q) nil))
  (prog1 (car (queue-head q))
    (setf (queue-head q) (cdr (queue-head q)))))

(defun queue-empty-p (q)
  (and (null (queue-head q))
       (null (queue-tail q))))
\end{lstlisting}
\end{solbox}

\label{sol:115}
\begin{solbox}{115}
\small\color{black!55}\hyperref[ex:115]{$\leftarrow$ Esercizio}\enspace\textcolor{black!30}{|}\enspace\hyperref[hint:115]{$\leftarrow$ Suggerimento}\normalsize\color{black}
\begin{lstlisting}
(defun valida-parentesi (stringa)
  ;; mappa ogni chiusa alla corrispondente aperta
  (let ((match '((#\) . #\()
                 (#\] . #\[)
                 (#\} . #\{)))
        (stack (make-stack)))
    (loop for c across stringa
          do (cond
               ;; aperta: push
               ((member c '(#\( #\[ #\{))
                (stack-push c stack))
               ;; chiusa: verifica top
               ((assoc c match)
                (if (or (stack-empty-p stack)
                        (not (char= (stack-pop stack)
                                    (cdr (assoc c match)))))
                    (return-from valida-parentesi nil)))))
    (stack-empty-p stack)))
\end{lstlisting}
\end{solbox}

\label{sol:116}
\begin{solbox}{116}
\small\color{black!55}\hyperref[ex:116]{$\leftarrow$ Esercizio}\normalsize\color{black}
\begin{lstlisting}
(defparameter *rubrica* (make-hash-table :test #'equal))

(defun add-contact (name phone)
  (setf (gethash name *rubrica*) phone))

(defun find-contact (name)
  (or (gethash name *rubrica*)
      (error "Contatto non trovato: ~A" name)))

(defun remove-contact (name)
  (remhash name *rubrica*))

(defun list-contacts ()
  (maphash (lambda (n p) (format t "~A: ~A~%" n p))
           *rubrica*))
\end{lstlisting}
\end{solbox}

\label{sol:117}
\begin{solbox}{117}
\small\color{black!55}\hyperref[ex:117]{$\leftarrow$ Esercizio}\normalsize\color{black}
\begin{lstlisting}
(defun neighbors (node graph)
  (cdr (assoc node graph)))

(defun add-edge (n1 n2 graph)
  ;; aggiunge arco bidirezionale
  (let ((g (copy-tree graph)))
    (let ((entry1 (assoc n1 g))
          (entry2 (assoc n2 g)))
      (when entry1 (setf (cdr entry1) (cons n2 (cdr entry1))))
      (when entry2 (setf (cdr entry2) (cons n1 (cdr entry2))))
      g)))
\end{lstlisting}
\end{solbox}

\label{sol:118}
\begin{solbox}{118}
\small\color{black!55}\hyperref[ex:118]{$\leftarrow$ Esercizio}\enspace\textcolor{black!30}{|}\enspace\hyperref[hint:118]{$\leftarrow$ Suggerimento}\normalsize\color{black}
\begin{lstlisting}
(defun dfs (graph start)
  ;; visited e' condivisa dalla closure tramite let
  (let ((visited nil))
    (labels ((visit (node)
               (unless (member node visited)
                 (push node visited)
                 (dolist (n (neighbors node graph))
                   (visit n)))))
      (visit start)
      (reverse visited))))
\end{lstlisting}
\end{solbox}

\label{sol:119}
\begin{solbox}{119}
\small\color{black!55}\hyperref[ex:119]{$\leftarrow$ Esercizio}\enspace\textcolor{black!30}{|}\enspace\hyperref[hint:119]{$\leftarrow$ Suggerimento}\normalsize\color{black}
\begin{lstlisting}
(defun bfs (graph start)
  (let ((visited nil)
        (queue   (list start)))
    (loop while queue
          for node = (pop queue)
          unless (member node visited)
          do (push node visited)
             (setf queue
                   (append queue
                           (remove-if (lambda (n) (member n visited))
                                      (neighbors node graph)))))
    (reverse visited)))
\end{lstlisting}
\end{solbox}

\label{sol:120}
\begin{solbox}{120}
\small\color{black!55}\hyperref[ex:120]{$\leftarrow$ Esercizio}\enspace\textcolor{black!30}{|}\enspace\hyperref[hint:120]{$\leftarrow$ Suggerimento}\normalsize\color{black}
\begin{lstlisting}
(defun tokenize (testo)
  ;; divide per spazi
  (let ((parole nil) (corrente ""))
    (loop for c across testo
          do (if (char= c #\Space)
                 (unless (string= corrente "")
                   (push corrente parole)
                   (setf corrente ""))
                 (setf corrente (concatenate 'string corrente (string c)))))
    (unless (string= corrente "")
      (push corrente parole))
    (reverse parole)))

(defun conta-parole (testo)
  (let ((ht (make-hash-table :test #'equal)))
    (dolist (w (tokenize testo) ht)
      (incf (gethash w ht 0)))))

(defun parola-piu-frequente (testo)
  (let ((ht (conta-parole testo))
        (best-w nil) (best-n 0))
    (maphash (lambda (w n)
               (when (> n best-n)
                 (setf best-w w best-n n)))
             ht)
    best-w))
\end{lstlisting}
\end{solbox}

% =============================================================
\section*{Macro}
% =============================================================

\label{sol:121}
\begin{solbox}{121}
\small\color{black!55}\hyperref[ex:121]{$\leftarrow$ Esercizio}\enspace\textcolor{black!30}{|}\enspace\hyperref[hint:121]{$\leftarrow$ Suggerimento}\normalsize\color{black}
\begin{lstlisting}
(defmacro quando (test &body corpo)
  ;; backquote: ,test inserisce il valore di test
  ;; ,@corpo: splice della lista corpo nel progn
  `(if ,test
       (progn ,@corpo)
       nil))

;; Verifica espansione:
;; (macroexpand-1 '(quando (> x 0) (print "ok") x))
;; -> (IF (> X 0) (PROGN (PRINT "ok") X) NIL)
\end{lstlisting}
\end{solbox}

\label{sol:122}
\begin{solbox}{122}
\small\color{black!55}\hyperref[ex:122]{$\leftarrow$ Esercizio}\enspace\textcolor{black!30}{|}\enspace\hyperref[hint:122]{$\leftarrow$ Suggerimento}\normalsize\color{black}
\begin{lstlisting}
(defmacro mio-while (test &body corpo)
  ;; usa loop nell'espansione: semplice ed idiomatico
  `(loop while ,test do ,@corpo))

;; Alternativa con tagbody/go (piu' primitiva):
;; (defmacro mio-while (test &body corpo)
;;   (let ((inizio (gensym)) (fine (gensym)))
;;     `(tagbody
;;        ,inizio
;;        (unless ,test (go ,fine))
;;        ,@corpo
;;        (go ,inizio)
;;        ,fine)))
\end{lstlisting}
\end{solbox}

\label{sol:123}
\begin{solbox}{123}
\small\color{black!55}\hyperref[ex:123]{$\leftarrow$ Esercizio}\enspace\textcolor{black!30}{|}\enspace\hyperref[hint:123]{$\leftarrow$ Suggerimento}\normalsize\color{black}
\begin{lstlisting}
(defmacro mio-for ((var da a) &body corpo)
  ;; gensym per start/end: evita valutazioni multiple
  ;; di da e a (effetti collaterali)
  (let ((start (gensym "START"))
        (end   (gensym "END")))
    `(let ((,start ,da)
           (,end   ,a))
       (loop for ,var from ,start to ,end
             do ,@corpo))))
\end{lstlisting}
\end{solbox}

% =============================================================
\section*{Backtracking e algoritmi}
% =============================================================

\label{sol:124}
\begin{solbox}{124}
\small\color{black!55}\hyperref[ex:124]{$\leftarrow$ Esercizio}\enspace\textcolor{black!30}{|}\enspace\hyperref[hint:124]{$\leftarrow$ Suggerimento}\normalsize\color{black}
\begin{lstlisting}
(defun subsets (lst)
  ;; ogni elemento: incluso o non incluso
  (if (null lst)
      '(())
      (let ((rest (subsets (cdr lst))))
        (append rest
                (mapcar (lambda (s)
                          (cons (car lst) s))
                        rest)))))
\end{lstlisting}
\end{solbox}

\label{sol:125}
\begin{solbox}{125}
\small\color{black!55}\hyperref[ex:125]{$\leftarrow$ Esercizio}\enspace\textcolor{black!30}{|}\enspace\hyperref[hint:125]{$\leftarrow$ Suggerimento}\normalsize\color{black}
\begin{lstlisting}
(defun permutations (lst)
  (if (null lst)
      '(())
      ;; per ogni elemento x: costruisci le permutazioni del resto
      (mapcan (lambda (x)
                (let ((rest (remove x lst :count 1)))
                  (mapcar (lambda (p) (cons x p))
                          (permutations rest))))
              lst)))
\end{lstlisting}
\end{solbox}

\label{sol:126}
\begin{solbox}{126}
\small\color{black!55}\hyperref[ex:126]{$\leftarrow$ Esercizio}\enspace\textcolor{black!30}{|}\enspace\hyperref[hint:126]{$\leftarrow$ Suggerimento}\normalsize\color{black}
\begin{lstlisting}
(defun combinations (lst k)
  (cond ((= k 0)    '(()))
        ((null lst)  nil)
        ;; con (car lst): combina k-1 dal resto; senza: combina k dal resto
        (t (append
            (mapcar (lambda (c) (cons (car lst) c))
                    (combinations (cdr lst) (1- k)))
            (combinations (cdr lst) k)))))
\end{lstlisting}
\end{solbox}

\label{sol:127}
\begin{solbox}{127}
\small\color{black!55}\hyperref[ex:127]{$\leftarrow$ Esercizio}\enspace\textcolor{black!30}{|}\enspace\hyperref[hint:127]{$\leftarrow$ Suggerimento}\normalsize\color{black}
\begin{lstlisting}
(defun somma-a-target (numeri target)
  (cond ((= target 0) '(()))
        ((or (null numeri) (< target 0)) nil)
        (t (let ((x (car numeri))
                 (rest (cdr numeri)))
             ;; includi x
             (append
              (mapcar (lambda (s) (cons x s))
                      (somma-a-target rest (- target x)))
              ;; non includere x
              (somma-a-target rest target))))))
\end{lstlisting}
\end{solbox}

\label{sol:128}
\begin{solbox}{128}
\small\color{black!55}\hyperref[ex:128]{$\leftarrow$ Esercizio}\enspace\textcolor{black!30}{|}\enspace\hyperref[hint:128]{$\leftarrow$ Suggerimento}\normalsize\color{black}
\begin{lstlisting}
(defun genera-parentesi (n)
  (let ((risultati nil))
    (labels ((gen (s aperte chiuse)
               (cond
                 ((= (length s) (* 2 n))
                  (push s risultati))
                 (t
                  (when (< aperte n)
                    (gen (concatenate 'string s "(")
                         (1+ aperte) chiuse))
                  (when (< chiuse aperte)
                    (gen (concatenate 'string s ")")
                         aperte (1+ chiuse)))))))
      (gen "" 0 0)
      (reverse risultati))))
\end{lstlisting}
\end{solbox}

\label{sol:129}
\begin{solbox}{129}
\small\color{black!55}\hyperref[ex:129]{$\leftarrow$ Esercizio}\enspace\textcolor{black!30}{|}\enspace\hyperref[hint:129]{$\leftarrow$ Suggerimento}\normalsize\color{black}
\begin{lstlisting}
(defun merge-sorted (lst1 lst2)
  (cond ((null lst1) lst2)
        ((null lst2) lst1)
        ((< (car lst1) (car lst2))
         (cons (car lst1) (merge-sorted (cdr lst1) lst2)))
        (t
         (cons (car lst2) (merge-sorted lst1 (cdr lst2))))))

(defun merge-sort* (lst)
  (let ((len (length lst)))
    (if (< len 2)
        lst
        (let* ((mid  (floor len 2))
               (left  (subseq lst 0 mid))
               (right (subseq lst mid)))
          (merge-sorted (merge-sort* left)
                        (merge-sort* right))))))
\end{lstlisting}
\end{solbox}

\label{sol:130}
\begin{solbox}{130}
\small\color{black!55}\hyperref[ex:130]{$\leftarrow$ Esercizio}\enspace\textcolor{black!30}{|}\enspace\hyperref[hint:130]{$\leftarrow$ Suggerimento}\normalsize\color{black}
\begin{lstlisting}
(defun quicksort* (lst)
  (if (or (null lst) (null (cdr lst)))
      lst
      (let* ((pivot   (car lst))
             (resto   (cdr lst))
             (minori  (remove-if-not (lambda (x) (< x pivot)) resto))
             (maggiori (remove-if (lambda (x) (< x pivot)) resto)))
        (append (quicksort* minori)
                (list pivot)
                (quicksort* maggiori)))))
\end{lstlisting}
\end{solbox}

\label{sol:131}
\begin{solbox}{131}
\small\color{black!55}\hyperref[ex:131]{$\leftarrow$ Esercizio}\enspace\textcolor{black!30}{|}\enspace\hyperref[hint:131]{$\leftarrow$ Suggerimento}\normalsize\color{black}
\begin{lstlisting}
(defun hanoi (n da a tramite)
  (when (> n 0)
    ;; passo 1: sposta n-1 da DA a TRAMITE (usando A)
    (hanoi (1- n) da tramite a)
    ;; passo 2: sposta il disco n da DA ad A
    (format t "Sposta disco ~D: ~A -> ~A~%" n da a)
    ;; passo 3: sposta n-1 da TRAMITE ad A (usando DA)
    (hanoi (1- n) tramite a da)))
\end{lstlisting}
\end{solbox}

\label{sol:132}
\begin{solbox}{132}
\small\color{black!55}\hyperref[ex:132]{$\leftarrow$ Esercizio}\enspace\textcolor{black!30}{|}\enspace\hyperref[hint:132]{$\leftarrow$ Suggerimento}\normalsize\color{black}
\begin{ragionamentoBox}{132}
La soluzione ricorsiva esplora tutte le $2^n$ combinazioni.
Per ogni oggetto ci sono due scelte: prenderlo o lasciarlo.
Restituiamo il valore massimo tra i due rami.
Complessita': $O(2^n)$ — accettabile per la versione didattica.
Per prestazioni migliori si userebbe la programmazione
dinamica con una tabella $n \times W$.
\end{ragionamentoBox}
\begin{lstlisting}
(defun knapsack (oggetti capacita)
  ;; oggetti: lista di (peso valore)
  (if (or (null oggetti) (= capacita 0))
      0
      (let* ((obj   (car oggetti))
             (resto (cdr oggetti))
             (peso  (car obj))
             (val   (cadr obj))
             ;; senza questo oggetto
             (senza (knapsack resto capacita)))
        (if (> peso capacita)
            senza
            ;; massimo tra prendere e non prendere
            (max senza
                 (+ val (knapsack resto (- capacita peso))))))))
\end{lstlisting}
\end{solbox}

\label{sol:133}
\begin{solbox}{133}
\small\color{black!55}\hyperref[ex:133]{$\leftarrow$ Esercizio}\enspace\textcolor{black!30}{|}\enspace\hyperref[hint:133]{$\leftarrow$ Suggerimento}\normalsize\color{black}
\begin{ragionamentoBox}{133}
DFS con backtracking su griglia. Usiamo un array separato
per i visitati per non modificare il labirinto originale.
La funzione restituisce il percorso se trovato, \fn{NIL} altrimenti.
\end{ragionamentoBox}
\begin{lstlisting}
(defun find-path (maze start goal)
  (let* ((rows (array-dimension maze 0))
         (cols (array-dimension maze 1))
         (visited (make-array (list rows cols) :initial-element nil)))
    (labels
      ((dfs (r c path)
         (cond
           ;; fuori dai bordi o muro o gia' visitata
           ((or (< r 0) (>= r rows) (< c 0) (>= c cols)
                (= (aref maze r c) 1)
                (aref visited r c))
            nil)
           ;; goal raggiunto
           ((equal (list r c) goal)
            (reverse (cons (list r c) path)))
           (t
            (setf (aref visited r c) t)
            ;; prova le 4 direzioni
            (let ((result (or (dfs (1+ r) c (cons (list r c) path))
                              (dfs (1- r) c (cons (list r c) path))
                              (dfs r (1+ c) (cons (list r c) path))
                              (dfs r (1- c) (cons (list r c) path)))))
              (setf (aref visited r c) nil) ; backtrack
              result)))))
      (dfs (car start) (cadr start) nil))))
\end{lstlisting}
\end{solbox}

\label{sol:134}
\begin{solbox}{134}
\small\color{black!55}\hyperref[ex:134]{$\leftarrow$ Esercizio}\enspace\textcolor{black!30}{|}\enspace\hyperref[hint:134]{$\leftarrow$ Suggerimento}\normalsize\color{black}
\begin{ragionamentoBox}{134}
\textbf{Rappresentazione:} vettore \fn{board} di \fn{n} interi, dove
\fn{(aref board col)} e' la riga della regina nella colonna \fn{col}.
Posizioniamo una regina per colonna, eliminando subito i conflitti.

\textbf{Pruning:} per la colonna \fn{col} e la riga \fn{row},
verifichiamo che nessuna colonna precedente \fn{pc} abbia:
\begin{itemize}[leftmargin=2em]
  \item stessa riga: \fn{(= (aref board pc) row)};
  \item stessa diagonale: \fn{(= (abs (- row (aref board pc))) (- col pc))}.
\end{itemize}
\textbf{Complessita':} con pruning, molto meglio di $O(n^n)$;
per $n=8$ si trovano 92 soluzioni in poche ms.
\end{ragionamentoBox}
\begin{lstlisting}
(defun n-queens (n)
  (let ((board (make-array n :initial-element 0))
        (solutions nil))
    (labels
      ((safe-p (col row)
         ;; nessun conflitto con le regine gia' posizionate
         (loop for pc below col
               never (let ((pr (aref board pc)))
                       (or (= pr row)
                           (= (abs (- row pr))
                              (- col pc))))))
       (place (col)
         (if (= col n)
             ;; soluzione trovata: registra una copia
             (push (coerce board 'list) solutions)
             (dotimes (row n)
               (when (safe-p col row)
                 (setf (aref board col) row)
                 (place (1+ col)))))))
      (place 0)
      solutions)))
\end{lstlisting}
\end{solbox}

\label{sol:135}
\begin{solbox}{135}
\small\color{black!55}\hyperref[ex:135]{$\leftarrow$ Esercizio}\enspace\textcolor{black!30}{|}\enspace\hyperref[hint:135]{$\leftarrow$ Suggerimento}\normalsize\color{black}
\solriferimento
\begin{ragionamentoBox}{135}
Backtracking con tre array di flag per righe, colonne e blocchi
$3\times3$ per un check $O(1)$ dei vincoli.
\end{ragionamentoBox}
\begin{lstlisting}
(defun solve-sudoku (board)
  ;; trova prima cella vuota
  (labels
    ((find-empty ()
       (loop for r below 9 do
         (loop for c below 9
               when (= (aref board r c) 0)
               do (return-from find-empty (list r c))))
       nil)
     (valid-p (r c val)
       ;; controlla riga, colonna, blocco 3x3
       (and (loop for i below 9 always (/= (aref board r i) val))
            (loop for i below 9 always (/= (aref board i c) val))
            (let ((br (* 3 (floor r 3)))
                  (bc (* 3 (floor c 3))))
              (loop for dr below 3 always
                (loop for dc below 3 always
                  (/= (aref board (+ br dr) (+ bc dc)) val)))))))
    (let ((empty (find-empty)))
      (if (null empty)
          board   ; risolto
          (let ((r (car empty)) (c (cadr empty)))
            (loop for val from 1 to 9
                  when (valid-p r c val)
                  do (setf (aref board r c) val)
                     (let ((result (solve-sudoku board)))
                       (when result (return-from solve-sudoku result)))
                     (setf (aref board r c) 0))
            nil))))) ; backtrack
\end{lstlisting}
\end{solbox}

\label{sol:136}
\begin{solbox}{136}
\small\color{black!55}\hyperref[ex:136]{$\leftarrow$ Esercizio}\enspace\textcolor{black!30}{|}\enspace\hyperref[hint:136]{$\leftarrow$ Suggerimento}\normalsize\color{black}
\solriferimento
\begin{ragionamentoBox}{136}
Backtracking con euristica di Warnsdorff: tra le mosse
disponibili, scegliamo sempre quella che porta alla cella
con meno uscite (mosse future disponibili).
Riduce drasticamente il backtracking rispetto al DFS puro.
\end{ragionamentoBox}
\begin{lstlisting}
(defparameter *moves*
  '((2 1)(2 -1)(-2 1)(-2 -1)(1 2)(1 -2)(-1 2)(-1 -2)))

(defun knights-tour (n)
  (let ((board (make-array (list n n) :initial-element 0)))
    (labels
      ((valid (r c)
         (and (>= r 0) (< r n) (>= c 0) (< c n)
              (= (aref board r c) 0)))
       (degree (r c)
         (count-if (lambda (m)
                     (valid (+ r (car m)) (+ c (cadr m))))
                   *moves*))
       (solve (r c step)
         (setf (aref board r c) step)
         (if (= step (* n n))
             t
             (let* ((next (remove-if-not
                           (lambda (m)
                             (valid (+ r (car m)) (+ c (cadr m))))
                           *moves*))
                    (sorted (sort next #'<
                                  :key (lambda (m)
                                         (degree (+ r (car m))
                                                 (+ c (cadr m)))))))
               (or (loop for m in sorted
                         thereis (solve (+ r (car m))
                                        (+ c (cadr m))
                                        (1+ step)))
                   (progn (setf (aref board r c) 0) nil))))))
      (when (solve 0 0 1) board))))
\end{lstlisting}
\end{solbox}

% =============================================================
\section*{Stringhe e I/O}
% =============================================================

\label{sol:137}
\begin{solbox}{137}
\small\color{black!55}\hyperref[ex:137]{$\leftarrow$ Esercizio}\normalsize\color{black}
\begin{lstlisting}
(defun reverse-string (s)
  (coerce (reverse (coerce s 'list)) 'string))
\end{lstlisting}
\end{solbox}

\label{sol:138}
\begin{solbox}{138}
\small\color{black!55}\hyperref[ex:138]{$\leftarrow$ Esercizio}\normalsize\color{black}
\begin{lstlisting}
(defun palindrome-p (s)
  (equal s (reverse-string s)))
\end{lstlisting}
\end{solbox}

\label{sol:139}
\begin{solbox}{139}
\small\color{black!55}\hyperref[ex:139]{$\leftarrow$ Esercizio}\enspace\textcolor{black!30}{|}\enspace\hyperref[hint:139]{$\leftarrow$ Suggerimento}\normalsize\color{black}
\begin{lstlisting}
(defun palindrome-robust-p (s)
  ;; normalizza: minuscolo, solo caratteri alfabetici
  (let ((clean (coerce
                (remove-if-not #'alpha-char-p
                               (coerce (string-downcase s) 'list))
                'string)))
    (equal clean (reverse-string clean))))
\end{lstlisting}
\end{solbox}

\label{sol:140}
\begin{solbox}{140}
\small\color{black!55}\hyperref[ex:140]{$\leftarrow$ Esercizio}\enspace\textcolor{black!30}{|}\enspace\hyperref[hint:140]{$\leftarrow$ Suggerimento}\normalsize\color{black}
\begin{lstlisting}
(defun anagram-p (s1 s2)
  ;; ordina le lettere di entrambe e confronta
  (equal (sort (coerce (string-downcase s1) 'list) #'char<)
         (sort (coerce (string-downcase s2) 'list) #'char<)))
\end{lstlisting}
\end{solbox}

\label{sol:141}
\begin{solbox}{141}
\small\color{black!55}\hyperref[ex:141]{$\leftarrow$ Esercizio}\normalsize\color{black}
\begin{lstlisting}
(defun count-char (c s)
  (count c s))
\end{lstlisting}
\end{solbox}

% =============================================================
\section*{Parsing e interprete}
% =============================================================

\label{sol:142}
\begin{solbox}{142}
\small\color{black!55}\hyperref[ex:142]{$\leftarrow$ Esercizio}\enspace\textcolor{black!30}{|}\enspace\hyperref[hint:142]{$\leftarrow$ Suggerimento}\normalsize\color{black}
\begin{lstlisting}
(defun rpn-eval (expr)
  (let ((stack (make-stack)))
    (dolist (token expr (stack-pop stack))
      (case token
        ((+ - * /)
         (let* ((b (stack-pop stack))
                (a (stack-pop stack)))
           (stack-push
            (case token
              (+ (+ a b)) (- (- a b))
              (* (* a b)) (/ (/ a b)))
            stack)))
        (otherwise
         (stack-push token stack))))))
\end{lstlisting}
\end{solbox}

\label{sol:143}
\begin{solbox}{143}
\small\color{black!55}\hyperref[ex:143]{$\leftarrow$ Esercizio}\enspace\textcolor{black!30}{|}\enspace\hyperref[hint:143]{$\leftarrow$ Suggerimento}\normalsize\color{black}
\begin{lstlisting}
(defun evaluate-expression (expr)
  ;; atomo numerico: restituisci direttamente
  (cond ((numberp expr) expr)
        ((symbolp expr) (error "variabile non supportata: ~A" expr))
        ;; lista: (operatore arg1 arg2)
        (t (let ((op  (car expr))
                 (a1  (evaluate-expression (cadr expr)))
                 (a2  (evaluate-expression (caddr expr))))
             (case op
               (+ (+ a1 a2))
               (- (- a1 a2))
               (* (* a1 a2))
               (/ (if (= a2 0)
                      (error "divisione per zero")
                      (/ a1 a2)))
               (otherwise (error "operatore sconosciuto: ~A" op)))))))
\end{lstlisting}
\end{solbox}

\label{sol:144}
\begin{solbox}{144}
\small\color{black!55}\hyperref[ex:144]{$\leftarrow$ Esercizio}\enspace\textcolor{black!30}{|}\enspace\hyperref[hint:144]{$\leftarrow$ Suggerimento}\normalsize\color{black}
\begin{ragionamentoBox}{144}
L'ambiente e' una alist \fn{((simbolo . valore) ...)}.
Il valutatore e' un'estensione dell'es.~143 con due nuovi casi:
simbolo (lookup) e forma \fn{set} (binding).
L'ambiente e' passato per valore; per modificarlo
permanentemente, \fn{set} deve restituire il nuovo ambiente
oppure usare un contenitore mutabile.
\end{ragionamentoBox}
\begin{lstlisting}
(defun mio-eval (expr env)
  "Valuta EXPR nell'ambiente ENV (alist simbolo->valore).
   Restituisce (valori . nuovo-env)."
  (cond
    ;; numero: restituisce se stesso
    ((numberp expr) (cons expr env))
    ;; simbolo: lookup nell'ambiente
    ((symbolp expr)
     (let ((binding (assoc expr env)))
       (if binding
           (cons (cdr binding) env)
           (error "variabile non definita: ~A" expr))))
    ;; forma (set var val)
    ((and (listp expr) (eq (car expr) 'set))
     (let* ((var (cadr expr))
            (res (mio-eval (caddr expr) env))
            (val (car res))
            (e2  (cdr res)))
       (cons val (acons var val e2))))
    ;; espressione aritmetica
    ((listp expr)
     (let* ((op (car expr))
            (r1 (mio-eval (cadr expr)  env))
            (r2 (mio-eval (caddr expr) (cdr r1)))
            (a  (car r1))
            (b  (car r2))
            (e2 (cdr r2)))
       (cons (case op
               (+ (+ a b)) (- (- a b))
               (* (* a b)) (/ (/ a b)))
             e2)))
    (t (error "espressione non riconosciuta: ~A" expr))))
\end{lstlisting}
\end{solbox}

% =============================================================
\section*{Progetti}
% =============================================================

\label{sol:145}
\begin{solbox}{145}
\small\color{black!55}\hyperref[ex:145]{$\leftarrow$ Esercizio}\normalsize\color{black}
\begin{lstlisting}
(defun gioca ()
  (let ((segreto (1+ (random 100)))
        (tentativi 0))
    (format t "Indovina il numero (1-100):~%")
    (loop
      (format t "> ")
      (let ((guess (read)))
        (incf tentativi)
        (cond ((< guess segreto) (format t "Troppo basso!~%"))
              ((> guess segreto) (format t "Troppo alto!~%"))
              (t (format t "Esatto! Tentativi: ~D~%" tentativi)
                 (return)))))))
\end{lstlisting}
\end{solbox}

\label{sol:146}
\begin{solbox}{146}
\small\color{black!55}\hyperref[ex:146]{$\leftarrow$ Esercizio}\enspace\textcolor{black!30}{|}\enspace\hyperref[hint:146]{$\leftarrow$ Suggerimento}\normalsize\color{black}
\begin{ragionamentoBox}{146}
Separazione netta tra Model (griglia), Logic (verifica vittoria,
mossa valida) e IO (input/output). La funzione \fn{vinto-p}
controlla le 8 combinazioni vincenti hard-coded
--- semplice e sufficiente per una griglia $3\times3$.
\end{ragionamentoBox}
\begin{lstlisting}
;;; MODEL
(defun crea-griglia () (make-array '(3 3) :initial-element nil))

;;; LOGIC
(defun mossa-valida-p (g r c)
  (null (aref g r c)))

(defun fai-mossa (g r c giocatore)
  (setf (aref g r c) giocatore))

(defun vinto-p (g p)
  (or ;; righe
      (loop for r below 3 thereis
        (loop for c below 3 always (eq (aref g r c) p)))
      ;; colonne
      (loop for c below 3 thereis
        (loop for r below 3 always (eq (aref g r c) p)))
      ;; diagonali
      (loop for i below 3 always (eq (aref g i i) p))
      (loop for i below 3 always (eq (aref g i (- 2 i)) p))))

(defun pareggio-p (g)
  (loop for r below 3 always
    (loop for c below 3 always (not (null (aref g r c))))))

;;; IO
(defun stampa-griglia (g)
  (dotimes (r 3)
    (dotimes (c 3)
      (format t "~A " (or (aref g r c) ".")))
    (terpri)))

;;; LOOP DI GIOCO
(defun gioca-tris ()
  (let ((g (crea-griglia))
        (turno 0))
    (loop
      (stampa-griglia g)
      (let ((p (if (evenp turno) 'X 'O)))
        (format t "Turno di ~A. Riga colonna: " p)
        (let ((r (read)) (c (read)))
          (if (mossa-valida-p g r c)
              (progn
                (fai-mossa g r c p)
                (cond ((vinto-p g p)
                       (stampa-griglia g)
                       (format t "~A vince!~%" p)
                       (return))
                      ((pareggio-p g)
                       (stampa-griglia g)
                       (format t "Pareggio!~%")
                       (return))
                      (t (incf turno))))
              (format t "Mossa non valida.~%")))))))
\end{lstlisting}
\end{solbox}

\label{sol:147}
\begin{solbox}{147}
\small\color{black!55}\hyperref[ex:147]{$\leftarrow$ Esercizio}\enspace\textcolor{black!30}{|}\enspace\hyperref[hint:147]{$\leftarrow$ Suggerimento}\normalsize\color{black}
\begin{lstlisting}
(defun conta-vicini (g r c rows cols)
  ;; conta le 8 celle adiacenti vive
  (let ((count 0))
    (loop for dr from -1 to 1 do
      (loop for dc from -1 to 1
            unless (and (= dr 0) (= dc 0))
            do (let ((nr (+ r dr)) (nc (+ c dc)))
                 (when (and (>= nr 0) (< nr rows)
                            (>= nc 0) (< nc cols)
                            (= (aref g nr nc) 1))
                   (incf count)))))
    count))

(defun evolvi (griglia)
  (let* ((rows (array-dimension griglia 0))
         (cols (array-dimension griglia 1))
         (nuova (make-array (list rows cols) :initial-element 0)))
    (dotimes (r rows nuova)
      (dotimes (c cols)
        (let ((v  (aref griglia r c))
              (n  (conta-vicini griglia r c rows cols)))
          (setf (aref nuova r c)
                (if (= v 1)
                    (if (member n '(2 3)) 1 0)
                    (if (= n 3) 1 0))))))))

(defun stampa-griglia-vita (g)
  (let ((rows (array-dimension g 0))
        (cols (array-dimension g 1)))
    (dotimes (r rows)
      (dotimes (c cols)
        (format t "~A" (if (= (aref g r c) 1) "#" ".")))
      (terpri))
    (terpri)))
\end{lstlisting}
\end{solbox}

\label{sol:148}
\begin{solbox}{148}
\small\color{black!55}\hyperref[ex:148]{$\leftarrow$ Esercizio}\enspace\textcolor{black!30}{|}\enspace\hyperref[hint:148]{$\leftarrow$ Suggerimento}\normalsize\color{black}
\textit{Soluzione di riferimento --- scheletro minimale per studenti.}

\begin{ragionamentoBox}{148}
Struttura a tre strati: MODEL (defstruct), LOGIC (operazioni),
IO (file e interazione). Il salvataggio usa \fn{print} (S-expression)
per semplicita': e' leggibile con \fn{read} senza parsing custom.
\end{ragionamentoBox}
\begin{lstlisting}
;;; MODEL
(defstruct studente nome cognome voto)

;;; ARCHIVIO
(defparameter *db* (make-hash-table :test #'equal))

;;; LOGIC
(defun inserisci (nome cognome voto)
  "Aggiunge o aggiorna uno studente."
  (let ((id (format nil "~A ~A" nome cognome)))
    (setf (gethash id *db*)
          (make-studente :nome nome :cognome cognome :voto voto))
    id))

(defun cerca (nome cognome)
  "Cerca uno studente; segnala errore se non esiste."
  (let ((id (format nil "~A ~A" nome cognome)))
    (or (gethash id *db*)
        (error "Studente non trovato: ~A" id))))

(defun elimina (nome cognome)
  "Rimuove uno studente."
  (remhash (format nil "~A ~A" nome cognome) *db*))

(defun filtra-per-voto (minimo)
  "Restituisce gli studenti con voto >= MINIMO."
  (loop for s being the hash-values of *db*
        when (>= (studente-voto s) minimo) collect s))

(defun ordina-per-voto ()
  "Restituisce la lista degli studenti ordinata per voto."
  (sort (loop for s being the hash-values of *db* collect s)
        #'> :key #'studente-voto))

;;; IO
(defun salva (file)
  "Salva il database su FILE come S-expressions."
  (with-open-file (s file :direction :output :if-exists :supersede)
    (maphash (lambda (k v)
               (print (list k (studente-nome v)
                              (studente-cognome v)
                              (studente-voto v))
                      s))
             *db*)))

(defun carica (file)
  "Carica il database da FILE."
  (clrhash *db*)
  (with-open-file (s file :direction :input :if-does-not-exist nil)
    (when s
      (loop for record = (read s nil nil)
            while record
            do (inserisci (second record)
                          (third record)
                          (fourth record))))))
\end{lstlisting}
\end{solbox}




\end{document}
